\documentclass[a4paper,11pt,openany]{ltjsbook}
\usepackage{luatexja}
\usepackage{luatexja-fontspec}
\usepackage{amsmath,amssymb,amsfonts,amsthm}
\usepackage{geometry}
\geometry{top=24mm,bottom=24mm,left=24mm,right=24mm,footskip=10mm}
\usepackage{tcolorbox}
\tcbuselibrary{skins,breakable}
\usepackage{hyperref}
\usepackage{fancyhdr}

\definecolor{primaryblue}{RGB}{20, 60, 110}
\definecolor{darkslate}{RGB}{35, 45, 60}
\definecolor{accentcrimson}{RGB}{160, 30, 45}

\hypersetup{
    colorlinks=true,
    linkcolor=primaryblue,
    citecolor=darkslate,
    urlcolor=accentcrimson,
    pdftitle={数学入門 —— 抽象の力と数理的思考の体系（原著全17章 完全翻訳・数学校閲版）},
    pdfauthor={アルフレッド・ノース・ホワイトヘッド}
}

\pagestyle{fancy}
\fancyhf{}
\fancyhead[LE,RO]{\thepage}
\fancyhead[RE]{\leftmark}
\fancyhead[LO]{\rightmark}
\setlength{\headheight}{14pt}
\renewcommand{\headrulewidth}{0.4pt}

\title{
    {\huge\bfseries\gtfamily 数学入門 —— 抽象の力と数理的思考の体系（原著全17章 完全翻訳・数学校閲版）}\\[14pt]
    {\large 原題：\textit{An Introduction to Mathematics}}\\[6pt]
    {\normalsize Project Gutenberg eBook \#41568}
}
\author{\textbf{アルフレッド・ノース・ホワイトヘッド} 著\\[4pt]教育情報分析研究会 訳・数学校閲}
\date{2026年10月 完全決定版}

\begin{document}
\maketitle
\frontmatter
\tableofcontents
\clearpage
\mainmatter

\chapter{第1章：Chapter I: The Abstract Nature of Mathematics}

数学の研究は、しばしば失望から始まる傾向がある。科学の重要な応用、概念的な興味深いアイデア、そして論理的な厳密さのある手法により、迅速に興味深い過程への導入を期待させる。私たちはその科学によって星が重りづけられ、一滴の水中の数十億の分子が数えられるということが教えられる。しかし、ハムレットの父の幽霊のように、この偉大な科学は私たちの精神的な武器による取り組みを逃れてしまう——「ここにいる、そこにもいる、そして消えた」——そして見えてくるものは、幽霊が粗雑な方法で説明できないほどの優れた存在であるという理由では説明されないような虚構性を示す。

この科学が期待通りにならない理由は、その基本的な概念が学生に技術的な手続きから解き放たれて説明されていないためである。したがって、不幸な学習者は、全体的な概念が光っていない多くの詳細を理解しようと必死になる。もちろん、技術的な能力は価値ある精神活動の第一条件である：ミルトンのリズムやシェリーの情熱を理解するまでに、単語を拼写し、個々の文字の形が確実にある間、その詩の美しさを十分に評価することはできない。この意味では、学習には王道は存在しない。しかし、技術的な過程にのみ注目し、全体的な概念を考えることから除外することも誤りである。ここに、教条主義への道がある。

次の各章の目的は、数学を教えることではなく、学生がコースの最初から科学が何であるか、そしてなぜ自然現象に対する正確な思考の基礎となるべきかを理解できるようにすることです。その後で述べられる科学内の詳細な推論への参照は例示のためのみであり、注意深く作られた一般的な議論は、読者が理解できない技術的な過程や記号が挙げられても、その内容が理解できるよう配慮されます。

多くの人々が数学と最初に出会う方法は算術を通じてです。二加二が四であることは通常、誰もが耳にしたことがある単純な数学者の命題の型例として取り上げられます。そのため、算術を考慮することで、可能であれば科学の最も明確な特性を見つけることができます。現在、算術における最初に注目すべき事実は、その適用範囲が広いことです。味や音、リンゴや天使、心のアイデアや体の骨まで何でも適用されます。物事がどのようなものであれ、二加二が四であることには変わりありません。したがって、数学の主要な特性として、物事そのものが理由ではなく、何かが何かであるという性質と概念を扱うことを記述します。そしてこれらの性質やアイデアは、何かが何かであるという点から適用され、それらと直接関連のある感情や感覚とは無関係です。これが数学を抽象的な科学と呼ぶ意味です。

The結果が達成されたことは注目すべきである。抽象的な科学が人間の生活のことでほとんど重要でないと思われるのは自然なことだが、それは実際の興味深い要素を考慮に入れていなかったからである。スウィフトは「ゴルビル物語」においてラプタ王国の数学者の扱いについて二つの見解を持っている。彼はその国の数学家たちを馬鹿で無用な夢想家とし、彼らの注意を覚醒させるために拍子木が必要だと描写している。また、数学的な裁縫師は自分の身長を四分法で測り、他の寸法は規則とコンパスを使って導き出し、非常に不適切な服を作っている。

一方で、ラプタ王国の数学者は空中に浮かぶ磁気島という驚異的な発明によって国を支配し、その支配力を維持している。スウィフトはニュートンの「原理」がまだ書かれている時代には、現代世界を大きく変革する力となった一方で、数学家たちを嘲るような時とは特別に適していなかった。スウィフトが地震を笑うように、彼もまた当時の数学家たちを嘲笑してもよかったかもしれない。

しかし、数学の成果の一覧だけではその重要性を理解するには不十分である。数学がなぜ、その抽象性ゆえに常に最も重要な思考のテーマの一つでなければならないのか、少し考えることでその根本的な理由を明確にする価値がある。我々はなぜ事件の順序を説明する際には必ずしも数学的になる傾向があるのかを理解しようと試みる。

すべての出来事は相互に結びついており、それを考慮する必要があります。雷を観察したときには、それに伴う thunder を聞くことになります；風が吹くと海面に波が立つことを観察します；秋の寒さを感じると葉が落ちます。秩序はあらゆる所で支配しており、特定の状況が見つかった場合、他の状況も同様に存在すると予測できます。科学の進歩はこれらの相互関係を観察し、この動的な世界における出来事が少数の一般的な関係や法則の例であることを示すことにあります。個別の事象の中に普遍性を見出し、瞬時なものの中に恒常性を見出すことが科学的思考の目的です。科学の眼から見れば、リンゴが落ちる現象、太陽繞る惑星の運動、大気の地球への密着は全て重力法則の例として観察されます。最も複雑な瞬時的な状況を様々な恒常的な法則の例に分解する可能性は、現代思考の支配的アイデアです。

では、この科学的理想を完全に実現するために求められる法則の種類について考えてみましょう。私たちは周囲の世界の特定の事実に関する知識を得るために、感覚から情報を得ます。私たちは見たり、聞いたり、味わったり、嗅いだり、熱を感じたり、押し付けられたり、痛みを感じたり、しびれを感じたりします。これらはすべて私個人の感覚です：私の歯痛はあなたの歯痛ではありませんし、私の視覚もあなたの視覚でもありません。しかし、これらの感覚の起源を外部世界の事物間の関係に帰属させます。つまり、歯医者は歯痛を取り除くのではなく、歯そのものを取り除きます。さらに、私たちはすべての人々の観察が基盤となる一つの連続的な世界があることを想像しようと試みます。私の感覚とあなたの感覚のために異なる世界があるわけではありません；私たち両者が存在する一つの世界があります。患者と歯医者にとって同じ歯です。また、私たちは視覚、触覚、そして聴覚で同じ世界を観察します。

従って、これらの外部の事物との関係を特定の感覚に依存しない方法で説明したいということが容易に理解できる。外部の事物の世界における事象の進行が満足する法則は、盲人や耳聋の人々にとっても、私たちの理解を超えた能力を持つ存在にとっても、同じような中立的かつ普遍的な形で記述されるべきである。

しかし、直感的な感覚を捨て去ったとき、残されたものの中で最も有用な部分は、事物の抽象的な形式的な性質に関する一般的な概念から成り立っている。実際には、先に述べた抽象的な数学的概念である。従って、人間は段階的に進み、その過程の完全な意味を理解せずに、宇宙の性質に対する数学者的な説明を求めることになる。なぜなら、この方法以外では、特定の人々や特定の感覚タイプに依存しない事象の進行の大まかな概念を形成することは不可能だからである。

例えば、食事中に問われるかもしれない質問は：
「私のかんじ、あなたの触れる感じ、そして彼のかんじと臭いは何が基盤になっていたのか？」
答えは「リンゴ」である。

しかし最終的には、科学はリンゴを分子の位置と運動から説明しようとし、その説明は私やあなたや彼、視覚や触覚や味覚や嗅覚を無視する。

従って、抽象的な数学的概念が科学的な事象の進行に対する説明に必要なのは、これらの概念が抽象的であるためである。

この点は通常、狭い観点から理解されてきた。
ピタゴラスはそのことを予感し、数が万物の源であると宣言した。
近代においては、ニュートン力学がすべての最終的な説明を含むものだと信じる信念は、科学が完璧に成長するにつれて全てが数学的になるという真実の一端を示していた。

\chapter{第2章：Chapter II: Variables}

数学は、最初に誰か、おそらくギリシャ人によって、任意のものやいくつかのものについての命題を証明したときから科学として始まった。これらの命題は、ギリシャ人が幾何学のために最初に述べたものであり、それに応じて、幾何学は偉大なギリシャ数学の科学となった。幾何学が興隆した後数世紀が経過するまで、代数は実質的な開始すらしなかった。それでも、後のギリシャ数学家中にはいくつかの予兆があった。

代数において、任意の数やある数という概念は文字を使用することで導入されます。具体的な算術の数字ではなく、文字$x$と$y$を使って、例えば $2 + 3 = 3 + 2$ の代わりに、任意の二つの数$x$と$y$について $x + y = y + x$ という一般化をします。

また、$3 > 2$ の代わりに、ある数$x$についてある数やいくつかの数$y$が存在し、$y > x$ であるという一般化を行います。なお、この後者の仮定は厳密な最終形で述べると仮定となるものですが、哲学と数学にとって極めて重要です。これは無限の概念を導入するためのものです。アラビア数字の導入により、文字が具体的な数を表すことを完全に廃止したことで、数学家たちに任意の数やある数という概念のために文字を使用することの技術的な利便性が示唆されました。ローマ人はこの文書が記された年をMDCCCCXと表していましたが、我々は1910と書き、文字は他の用途に使用します。しかしこれは単なる推測です。

代数の発展とともにニュートンとライプニッツによって微分 calculus が発明されましたが、これらの概念に関しては数学的思考の哲学における進歩は一時的に止まりました。そして最近の数年間で初めて、任意の数やある数が数学そのものに極めて根本的な役割を果たすことが認識されました。これによりさらなる数学的研究分野が開拓されることになりました。

以下に簡単な代数的声明をいくつか述べて、これらの基本的な概念がどのように現れるかを理解するために考えてみましょう。

\begin{align}
(1) \& \quad 任意の数$x$について $x + 2 = 2 + x$; \\
(2) \& \quad ある数$x$について $x + 2 = 3$; \\
(3) \& \quad ある数$x$について $x + 2 > 3$.
\end{align}

最初に注意すべきは、「some」という言葉の含意にある可能性です。ここでの「any」はすべての数に対して成り立つことを示していますが、$x + 2 = 2 + x$ は任意の数$x$に対して真です。したがって、ある数$x$についても真となります。「any」という言葉は「some」を含むことから、「some」は「any」を排しないことがわかります。

また、第二の例では実際には一つだけの数$x$が存在し、$x + 2 = 3$ を満たすもので、それは1のみです。したがって、「some」はただ一つの数を指すこともあります。しかし第三の例では、任意の数$x$が1より大きいとき、$x + 2 > 3$ が成り立ちます。これにより、この場合の「some」に対応する無限多種の数がわかります。したがって、「some」は「any」と「one only」の間で何でもあり、両端の境界ケースも含めることができます。

次に自然な形で式（2）と式（3）を質問に置き換えることができます：

式（2′）$x + 2 = 3$ を満たす数$x$は何か？

式（3′）$x + 2 > 3$ を満たす数$x$は何か？

考慮する式（2'）より、$x + 2 = 3$は方程式であり、その解は容易に$x = 3 - 2 = 1$であることが分かります。
方程式 $x + 2 = 3$ の文脈から推測される問いに対する答えにおいて、$x$は未知数と呼ばれます。方程式の解き方は未知数を特定することです。方程式は数学において非常に重要であり、式（2'）は元々の式（2）よりもはるかに根本的な概念を示しているように思われますが、これは完全な誤りです。不定な「変数」という概念が「何か」や「どれか」を使用する際に現れるものこそが数学において真に重要な概念であり、方程式における未知数は速やかに解くべきものとして重要ですが、その役割は補佐的なものです。もちろんこれは非常に重要です。初等代数の大部分が表面上単純である理由の一つは、教科書が方程式の解き方に過度に集中していることにあると言えます。不等式（3'）も元々の式（3）と比較して同様の主張が成り立つことを述べることができます。

特に「何か」の概念が含まれる場合には、多くの興味深い公式は一つ以上の変数を含みます。例えば、$x + y = 1$ を満たす分数または整数のペア $x$ と $y$ の考慮は、二つの関連する変数 $x$ と $y$ の概念を含んでいます。二つの変数が存在する場合、主に二種類の声明が現れます。例えば、任意の数のペア $x$ と $y$ に対して $x + y = y + x$（式（1））、あるいはいくつかの数のペア $x$ と $y$ に対して $x + y = 1$（式（2））です。

第二の種類の命題は、固定された関係によって結び付けられた数のペアの全体を考慮することを招く。この例では、$x + y = 1$という関係によって結ばれた数のペアを考える。

最初のタイプの公式が任意の数のペアに対して真である場合、それらの公式は第二のタイプの公式を無限に多くの等価な形へと変形することができる。例えば、$x + y = 1$という関係は以下の関係に等しい：
\[
y + x = 1,\quad
(x - y) + 2y = 1,\quad
6x + 6y = 6,
\]
などである。したがって、命題を考慮している関係の適切な等価形は、その直近の目的にとって最も便利なものを使い、巧みな数学者はそれを用いる。

一般には、あるペアの項が固定された関係を満たす場合でも、一方の項が与えられたからといってもう一方の項が確定的に定められるわけではありません。例えば、$y^2 = x$という関係において$x = 4$とすると、$y$は$\pm 2$となり得ますので、任意の正の$x$に対して異なる値を持つ$y$が存在します。また、$x + y > 1$という関係において、$x$や$y$のいずれかが与えられたからといってもう一方の項には無数の値が残ります。

また重要な点が一つあります。正の数、整数または分数を考慮する場合、関係式 $x + y = 1$ を考えます。このとき、もし $x$ や $y$ のいずれかが $1$ より大きい値を持つと、もう一方の数が正の数としてこの関係を満たす値は存在しません。したがって、関係式 $x + y = 1$ における $x$ の「領域」は $1$ より小さい数に制限され、同様に $y$ の「領域」も制限されます。

また、整数（正または負）のみを考慮し、関係式 $y^{2} = x$ を満たすような数のペアを考えます。このとき、$y$ に与えられたどんな整数値にもかかわらず、$x$ は一意に対応する整数値を取ることができます。したがって、$y$ の「領域」はこれらの正または負の整数全体で制限されません。しかし、$x$ の「領域」は二つの方法で制限されます。まず第一に $x$ は正である必要がありますし、第二に $y$ が整数であるため $x$ は完全平方数である必要があります。したがって、$x$ の「領域」は $1^{2}$, $2^{2}$, $3^{2}$, $4^{2}$, など、すなわち $1$, $4$, $9$, $16$, などの整数の集合に制限されます。

数値ペア間の関係の一般的な性質を研究する際には、次のように構築された図を使って多くの助けになります。

二直線 $OX$ と $OY$ を垂直に引く。
任意の数 $x$ を $OX$ 上の $x$ 単位の長さ、任意の数 $y$ を $OY$ 上の $y$ 単位の長さで表す（いずれも任意のスケール）。
したがって、$OX$ 上の $OM$ が $x$ 単位の長さを持つとき、$OY$ 上の $ON$ が $y$ 単位の長さを持つとき、平行四辺形 $OMPN$ を完成させることで点 $P$ が得られ、この点は数対 $(x, y)$ に対応する。
各点には一意の数対が対応し、各数対には一意の点が対応する。
この数対は点の座標と呼ばれる。したがって、固定された関係を満たす点の座標を便利な方法で示すことができる。すべての点が一直線上にある場合は直線を引くことができ、領域全体に存在する場合にはその領域を塗りつぶすことができる。
例えば、正の数のみを考えるとき、$x + y = 1$ を満たす座標を持つ点は図 \Fig{1} の直線 $AB$ 上にある。ここで $0A = 1$ かつ $OB = 1$ である。
したがって、この直線 $AB$ の一部は正の数に対する関係の性質を視覚的に表現するものとなる。

二つの変数間の関係の別の例は、一定温度下での一部の気体物質（例えば空気や石炭ガスや蒸気など）の圧力と体積の変動を考慮することで得られます。そのとき、$v$ をその体積が立方フィートで表す数値、$p$ をその圧力が1平方インチあたりポンドで表す数値とします。このとき、圧力と体積の両方の変化を表現するボイルの法則は、常に温度が変わらないという仮定のもとで、積$pv$ が一定であることを述べています。例えば、ガスの量やその他の状況が一定であり、$pv = 1$（等式の右辺の正確な数値は基本的には重要ではありません）となるように仮定することができます。

次に、図2において、直角に交わる線分$OV$と$OP$を取り、$OM$を$OV$上に引き、$v$立方フィートの体積を表すようにします。また、$ON$を$OP$上に引き、$p$ポンド力あたり平方インチの圧力を表すようにします。すると、平行四辺形$OMQN$を完成させた点$Q$は、そのときガスの体積が$v$立方フィートで圧力が$p$ポンド力あたり平方インチである状態を表します。考慮しているガスの状況が$pv = 1$となるような場合であれば、それに対応するすべての点$Q$は、$p$と$v$が正で$pv = 1$を満たすすべての点を含む曲線$ABC$上に存在します。つまり、この曲線は体積と圧力との関係を視覚的に表しています。圧力が非常に大きい場合、対応する点$Q$は$C$に近い位置にあるか、あるいは未描きの部分にあるかもしれません。すると、その体積は非常に小さいことになります。体積が大きい場合、$Q$は$A$に近い位置にあるか、あるいはそれ以遠にあるかもしれません。すると、圧力は小さくなります。注意していただきたいのは、エンジニアや物理学者にとって、特定の体積に対応する圧力を知る必要がある場合があります。つまり、已知数$v$に対する未知数$p$を決定する問題です。しかし、これは特定の場合に限られます。ガスの性質とその行動を一般に考慮するときには、全体の曲線$ABC$の一般的な形状とその一般的な特性を考えなければなりません。言い換えれば、根本的な考えは、関係$pv = 1$を満たす変数のペアであるということです。この例は、変数の概念が数学の理論および応用においても基本的なアイデアであることを示しています。

\chapter{第3章：Chapter III: Methods of Application}

変数が関係を満たすという概念が数学の応用においてどのように現れるかについては、深く考えることが価値があります。そのような概念について少し時間を費やすことで、全体的な主題に対する私たちの思考を明確にすることができます。

最も単純な例から始めましょう：建設費用は1立方フィートあたり£1とします。さらに£20が1ポンドであるとします。新しい家を建てる複雑な状況の中で、所有者、建築家、builders（原文では「builder」ですが、ここでは複数形で表記）、労働者、見物人など、家が完成するまでのさまざまな感情や感覚の間、固定された関係は法によって保たれるものと仮定されます。つまり、$x$が立方フィートの数、£$y$が費用であるとき、$20y = x$という関係が成り立つと仮定します。この$x$と$y$の関係は、どの所有者がどの家を建設する際にも真実であると仮定されます。また、家の体積と費用は特定の感覚や能力によって感知または理解されたものではなく、抽象的な一般的方式で述べられ、所有者が請求書を支払うときの精神状態に関係なく、完全に無関心です。

今更考えてみれば、これらすべては何を意味しているのか。家を建てるという行為は複雑な状況である。法則を適用したりテストするためには、一般的な出来事の中で特定の一連の出来事が家を建設する具体的な例として認識できることが必要だ。つまり、家を見たときにそれが何かを知らなければならないし、その建設に関連する出来事を認識しなければならない。これらの出来事の中から、概念的に他の自然から隔離された場所で、コストと立方体の内容量の両方を決定することが可能である。それらが両方とも決定されると、法則が真実であれば、一般式
\[
20y = x
\]
を満たす。

しかし、その法則は真実でしょうか？家を建設に関わった経験がある人なら、ここでのコストは高いと感じるでしょう。この価格で機能するのは高級なタイプの家だけです。これにより別の点が明確になる。公式 $20y = x$ に関連する数学的な計算をしている間、法則が真実であれ偽であれ、それは私たちにとって関係ありません。実際、$x$ と $y$ の意味付けは立方フィートの数とポンド Sterling の数であるとしても、それ自体には関係ありません。数学的調査中では、実質的にこの変動する数字のペア $x$ と $y$ の間の相関性の性質を考慮しているに過ぎません。結果は、$y$ を漁師の数と解釈し、$x$ を捕れた魚の数と解釈しても同様に適用されます。つまり、平均して1人の漁師が20匹の魚を捕れるという法則が仮定されている場合でも、調査結果は公式 $20y = x$ の相関性の性質を与えるもので、数学的な確実性はその結果にのみ付与されます。具体的な家を建設するコストには何の数学的確実性もありません。法則は完全には真実ではなく、それを提供する結果も正確とは限らない。実際、それは完全に間違っている可能性もあります。

現在これらは確かに非常に明白なことのように思われます。
しかし、より複雑な例では、数学的な計算が prolonged かつ正確に行われたとしても、その結果を自然の事実に適用する確実性は全くありません。議論の結論は、その始まりの仮定よりも決してより確実ではありません。自然の進行に関するすべての数理的計算は、建設コストの仮定された法則のような、何かしらの自然の仮定された法則から始めなければなりません。
したがって、ある出来事が正確に予測されるほど正確な計算を行っても、その法則は真実であるかどうかの疑問は常に残ります。もし法則が具体的な結果を述べているなら、ほぼ確実にそれは正確ではありません；そして、最良の場合でも、厳密に計算された結果はおそらく起こりません。しかし、我々には理想的な精度で観察する能力がないため、結局のところ、不正確な法則でも十分良いものであるかもしれません。

我達现在来考虑一个实际的例子，即牛顿和万有引力定律。この法則は、任意の2つの物体が質量の積に比例し、それと逆比例する距離の平方に引き合う力をもつことを述べています。つまり、2つの物体の質量を$m$と$M$（lbs.で計測）とし、その間の距離を$d$（マイル）とするとき、一方の物体に対する他方からの引力は$\dfrac{mM}{d^{2}}$に比例します；したがって、この力を$k\dfrac{mM}{d^{2}}$と表すことができます。ここで$k$は、この引力の絶対的な大きさと、我々が選んだ力の尺度によって定まる一定の数です。もし1lbs.の質量に対する重力として力を計測したい場合、$k$が代表する数値は非常に小さいことがわかります。なぜなら、$m$、$M$、および$d$をそれぞれ1に置き換えると、$\dfrac{kmM}{d^{2}}$は1lbs.の等しい質量間の距離1マイルでの引力となり、これはほとんど感じられません。

しかし、現在私たちは引力の力の公式を得ました。この力を$F$と呼びますとすれば、その式は$F = k\dfrac{mM}{d^{2}}$となります。これにより変数$F$、$m$、$M$、および$d$との関係が示されます。この法則の発見物語を知っています。

ニュートンは orchard でリンゴが落ちるのを見ていて、その後、普遍的な重力法則が彼の頭に浮かび上がったとされています。最終的な法則の公式が彼の orchard で思いついただけでなく、他でも考えたかもしれません——彼にはどこかにいる必要があります。しかし、私たちの目的からは、この正確な法則を導くために必要な膨大な前準備的思想、多くの頭脳と数世紀分の産物について深く触れることはより有益です。

まず第一に、前の二章で説明した数学的な思考習慣や手続きが生成される必要がありました。それなしではニュートンは、任意の質量間の力を表現する公式を考え出すことはできなかったでしょう。また、使用されている用語の意味、つまり力、質量、距離について考えてみましょう。距離という最も簡単な用語から始めます。私たちは当然のようにすべての物質が一定の幾何学的な全体として考えられることを理解していますし、その各部分間の距離は特定の単位長さ（マイルやヤードなど）で測定できると考えています。これは私たちが最初に思いつく材料構造の最も基本的な側面です。これは幾何学の研究と測量理論の考察の結果として徐々に達成されます。現在でも、特定の場合では他の思考方法が便利なこともあります。山岳地帯では距離はしばしば時間で計測されることがあります。しかし、距離を除く、力と質量という他の用語は遥かに不明瞭です。ニュートンがこれらの言葉で伝えようとした思想の正確な理解はゆっくりと成長し、実際にはニュートン自身が初めて動力学の真の一般的原理を完全に把握した最初の人でした。
\index{重力}\%
\index{ニュートン}\%
\index{動力学}\%

中世を通じて、アリストテレスの影響下で、科学は完全に誤解されていました。ニュートンは、ガリレオを含む一連の偉大な人物の後ろに立つ利点がありました。イタリアでは、前2世紀にガリレオが科学を再構築し、正しい思考方法を発明しました。ニュートンは彼らの仕事を完成させました。その後、力、質量、距離という概念について、頭の中で明確で鮮明な理解を持ち、その重要性とリンゴの落下や惑星の運動にどのように関連しているかを認識したとき、万有引力の法則を見つけて証明しました。

数学的公式の適用における重要なポイントは、観察対象の現象に対する明確な概念と正確な関連性の評価を持つことです。私たち自身と同じように、遠い祖先たちは自然現象の重要性に感銘を受け、事件の順序を規律するために積極的な措置を講じるべきであるという点でも影響を受けました。無関係な考えに基づいて、新月の誕生を助けるための複雑な宗教儀式を行い、日食の危機時に太陽を救うために犠牲を捧げました。彼らが私たちよりも愚かだったとは思われません。しかし当時の段階では、明確で関連性のある概念の積み重ねという時間的な経験がなかったのです。

物理学の諸分野が数学的手法による扱いに適する形になる過程は、電磁気学の科学的発展を通じて例示される。雷暴は大規模な出来事で、人間だけでなく動物も恐怖を引き起こす。古代から、雷暴は狂騒的なそして幻想的な仮説の対象となってきたが、現代の電気に関する科学的発見が原始社会の魔法的な説明よりも驚異的であるかどうかは疑問に値する。

ギリシャ人は、琥珀（ギリシャ語で electron）を擦ると軽いそして乾燥した物体に引き寄せられることがわかると知っていた。1600年頃、コルチェスターのドクター・ジルバートは、どのような科学的手法も用いる最初の著作を出版した。彼は、琥珀と同じような性質を持つ物質の一覧を作成し、電気と磁気の現象が多少明確に結び付けられたことも認められる。

17世紀末から18世紀を通じて知識は進展した。電気機械が作られ、火花が得られ、レインドール Jar が発明され、これらの効果を強化する手段が提供された。組織的な知識が得られつつあったが、まだ関連する数学的概念は見つからなかった。

1752年、フランクリンは空に kite を上げて雷暴が電気であることを証明した。

同時から最古の時代（紀元前2634年頃）より、中国では磁針の特性を用いていましたが、理論的な考えとは結びつけていなかったようです。人類生活における真に深い変化は、その最終的な源を自己のための知識追求に求めることができます。磁針の使用は、12世紀末までヨーロッパには導入されず、中国で初めて使用された後約3000年後のことでありましたが、これは欧州人の優れた実用的偏見によるものではなく、西で電気や磁気現象を抽象的な理論的関心に支配された人々が研究したことに归因されます。

電流の発見は1780年にガルバーニと1792年にボッタによってイタリア人によってなされました。この偉大な発明は新たな現象の調査を可能にしました。科学界には、摩擦による静電気効果、磁気現象、電流による効果という三つの別々だが関連のある事象が手に入りました。18世紀末からこれらの三つの研究路線は迅速に結びつき、現代の電磁力学の科学が構築され、これにより人類生活が変革される可能性があるようになりました。

数学的なアイデアが現れた。1780年から1789年の間に、フランス人であるコルトン \index{コルトン} 誰かが証明したところによれば、磁極は距離の逆平方に比例して互いに引き寄せられたり斥け合ったりする。また、電荷についても同じ法則が成り立つことが示された——これは重力法則と奇妙な類似性を持つ法則である。1820年には、デンマーク人オーレスト \index{オーレスト@オーセスト} が発見したところによれば、電流は磁石に力を与えることが確認された。その直後に、フランス人アムペール \index{アムペール@アムペール} が、この力の数学的な法則を正しく導き出し、また2つの電流が互いに力を及ぼすことを証明した。「アムペールが電流間の力学的相互作用の法則を確立する実験調査は科学における最も輝かしい達成事例の一つである。理論と実験全体は、『電気のニュートン』といわれる人物の脳から一気に生まれたかのように見受けられる。その形式は完璧で、正確性は攻撃不能であり、すべての現象を導き出す公式が含まれており、常に電磁動力学の枢要な公式として残ることだろう。」\footnote{\Title{電気と磁気}, クラック・マクスウェル, 第2巻, 第3章。}

マイケル・ファラデーは1831年から1838年まで、電流間の感 induction法則と磁石との間の感 induction法則を発見しました。
\index{ファラデー}\%
ファラデーに質問がimed: 「この発見は何のために有用ですか?」彼は答えた：「子供は何のために有用か---成長して男になるからである。」
ファラデーの子供は成長し、現在電気のすべての現代的な応用の基礎となっています。
\index{クレッグ・マクスウェル}\%
ファラデーもまた科学の全体的な理論的概念を再組織しました。彼の理解が充分でなかったアイデアは1873年にクレッグ・マクスウェルによって直接数学的形式に拡張されました。
彼の数理的調査結果から、マクスウェルは特定の条件下では電気振動が伝播するべきであると認識しました。彼は即座に、光を形成する振動が電気的なものであるという推測を提案しました。この推測は以来検証され、現在の光の理論は大いなる電気学の一派となっています。
\index{光}\%
また1888年にドイツ人のヘルツはマクスウェルのアイデアに基づき、直接的な電気的方法で電気振動を生成することに成功しました。彼の実験は無線電信の基礎となっています。

最近の年々ではより根本的な発見がなされ、科学は理論的的重要性と実用的的重要性において引き続き成長しています。その急速な進歩は、関連する理論的概念の逐次導入によって示されます。これは実験によって提示され、また新たな実験を提案し、孤立した甚至は平凡な現象が一貫性のある科学に統合される過程を表しています。抽象的な数学的推論から少数の単純な仮定法則を出発点として得られる結果は、事象の進行の複雑な絡み合いに対する説明となります。

最終的に、電磁気学や光の特別な科学を越えて、私たちは視点をさらに一般化し、数学的物理学の成長を一つの大切なる科学思想の一章として注目することができます。まず最初に、その成長の物語はどのようなものだったでしょうか？

それは一門の科学から始まったわけではなく、一人の集団によって生み出されたものでもありませんでした。カルデアの羊飼いたちは星空を観察し、美索不達メラやエジプトの政府職員たちは土地を測定し、祭司と哲学者たちは万物の一般的な性質について考えました。自然の大きな作用は、神秘的な深淵にある力によって引き起こされているように見えていました。「風はその都合で吹く」は当時の現象の詳細に従う稳定的な法則が存在しないことの正確な表現でした。大まかには今でも同様ですが、事件の規則性は明らかでした。しかし、それらの相互関係を細部まで追跡することは不可能であり、そのような科学を構築する方法すら知る手段もありませんでした。

個別の推測や、物事の本質に対する数少ない幸運な或いは不幸な試みが、これ以上のものが生み出された限界でした。

Meanwhile, 地図作成は幾何学を生み出し、
\index{Geometry}\%
そして天体観測は太陽系の正確な規則性を明らかにしました。その後期のギリシャ人たちはアキメデスのような者がいました。彼は水静力学や光学の基本的な現象について正解的な見解を持っていました。実際、数学における天才と物理的洞察力を持つアキメデスは、ほぼ2000年後のニュートンと共に数理物理学の創始者として評価されるべきです。彼はシチリアの大きなギリシャ都市サルデナに住んでいました。ローマ軍が町を包囲したとき（紀元前212年から210年）、彼は鏡を使用して太陽光線を集束させて、敵の船を焼いたという話があります。この話は非常に不確かなものですが、彼が光学に関する知識で当時の同僚たちから高い評価を得ていたことを示す良い証拠です。包囲戦終了後、彼は殺されました。プラUTERSのマセルス伝によると、ある説では、彼は捕らわれた兵士によって見つかり、部屋の砂利に描いた幾何学的な図形を研究していたとされています。直ちに命令に従わなかったため、殺されてしまったということです。ローマ軍総帥の名誉のために言わなければなりませんが、彼を殺す命令は出されていなかったということです。彼に関するもう一つの有名な話の内部証拠は非常に強いです；その発見は彼の数学的・物理的研究における天才性にふさわしいものです。
\PageIndex{Plutarch}\%
\index{Marcellus}\%
\PageIndex{Newton}
\index{Archimedes|EtSeq}\%
\PageIndex{Geometry}

幸いにも、それは十分単純で、ここでは詳細に説明することができます。これは数学的概念を物理学に適用する方法の最も良い簡単な例です。

ヒエロ、シラクサの王が金を金工人に送り、冠の材料とするよう命じた。彼は金工人が部分的な金を取り除き、その場所に価値の低い金属で合金を作成したと疑っていた。ヒエロはアリキメデスに冠を送り、それを検証するように頼んだ。当時の時代では無数の化学的試験が利用可能だったかもしれないが、そのときアリキメデスは問題を再考しなければならなかった。彼は入浴しているときに解き明かす閃きを得た。彼は立ち上がり、宮殿を通って町を駆け抜け、「エウレカ！ エウレカ！」と叫んだ。この日が数学的物理学の誕生の日であるべきだ。ニュートンが果樹園で座っているときには科学は成熟したのである。アリキメデスは実際には大きな発見をした。彼は物体が水に浸かると、周囲の水によって上向きの力を受けることに気づき、その結果的な力は物体が排水する水の重さと等しいという法則であることを理解した。この法則は水静力学の数学的原理から理論的に証明でき、また実験的に検証することもできる。従って、冠の空気中の重量を$W$ポンド、完全に浸かるとき排水する水の重量を$w$ポンドとする場合、$W - w$は冠が水中で吊るされているときに必要な追加の上向き力となる。

今、この上向きの力は水に浸かっている物体を秤で計ることによって容易に測定できます。図に示すように、右側の秤の重さが$F$ポンドである場合、冠の水中的な重量は$F$ポンドとなります。したがって、

\[
F = W - w
\]

となり、

\[
w = W - F,
\]

そして

\[
\frac{W}{w} = \frac{W}{W - F}
\Tag{(A)}
\]

ここで、$W$と$F$は秤で容易に測定し、比較的正確な値を得ることができます。

したがって、式(A)より$\dfrac{W}{w}$は分かることになります。しかし、

\[
\dfrac{W}{w}
\]

は冠の重量に対する等量の水の重量の比です。この比は同じ材料で作られた任意の塊の金属に対して同じであり、現在はその材料の特定密度と呼ばれ、物質の固有の性質にのみ依存し、形状や量には依存しません。したがって、冠が金であるかどうかを検証するためには、疑いの余地のない純金の塊を取り出し、同じ手順で特定密度を求めればよいです。二つの特定密度が一致すれば冠は純金であり、不一致であれば偽物であることがわかります。

この議論は詳しく述べたのは、なぜならこれは数学的概念を物理学に適用する最初の正確な例であり、またこれこそが科学にとって永遠の方法と精神となるべき完璧でシンプルなモデルだからです。

アキメデスの死はローマ兵の手によって象徴的に世界の変化を示しており、抽象的な科学を愛した理論的ギリシャ人が、実用主義的なローマ人によってヨーロッパ世界の指導権を奪われたことを表しています。ベーコンフィールド伯爵は、自前の誤りを実践する者が実用家であると定義しています。ローマ人は偉大な民族でしたが、実用性に伴う不妊性に呪われていました。彼らは先祖の知識を改善せず、全ての進歩はエンジニアリングの細かい技術的詳細に留まりました。十分な夢見がなく、自然の力に対するより根本的な制御につながる新しい視点に至ることができませんでした。ローマ人は数学図形に没頭しているために死ぬことはありませんでした。

\chapter{第4章：Chapter IV: Dynamics}

世界は十八世紀まで待たされ、ギリシャの数学物理学者がその後を引き継いだ。当時の16世紀と17世紀における偉大なイタリア人、特にアーリストオレス（画家）、レオナルド・ダ・ヴィンチ（生年1452、没年1519）とガリレオ（生年1564、没年1642）は、アーキメデスが知っていた秘密を再発見し、抽象的な数学的概念と自然現象の実験的調査との関連性を見つけることに成功した。一方で、数学の進歩と正確な天文学知識の蓄積により、自然哲学者は研究においてより有利な立場に置かれた。また、その時代の自我中心的な自己主張や個人経験への渇望は、思考者たちを自身で現象を見るために動かし、帰納的論理における数学理論と実験との関係の秘密が実質的に解明された。

ガリレオという哲学者がピサの傾斜塔から重りを落としたことは、その時代の特徴的な行為だった。常に思考者と行動者の存在がある。数学物理は、同じ人物の中に思考の衝動と行動の衝動を組み合わせた時代の産物である。

この塔から重さの異なる物体を落とすという事柄は、知識における重要な段階を象徴的に示しています。これは、全分野の基礎となる動力学に関する正しい概念を得た最初のステップに匹敵するものです。議論の焦点は、同じ高さから落とされた異なる重さの物体が同じ時間で落ちるかどうかでした。アリストテレスの定説によれば、その時代まで普遍的に受け入れられていたように、 heavier な物体の方が速く落ちるとされました。ガリレオはこれらが同じ時間で落ちることを主張し、彼の主張を立証するために、傾斜した塔の頂上から重さの異なる物体を落としました。

この規則に対する例外的なケースは、例えば極めて軽量や高速な理由により空気抵抗が重要となる場合に生じます。しかし、空気を無視すれば法則は正確です。

ガリレオの成功した実験は単なる運の良さによるものではなく、運動の不変性と質量に関する彼の正しい概念から生まれたものです。現在ニュートンが述べる最初の運動法則は次の通りです---全ての物体は、強制的な力によってその状態を変えるまで、静止状態または一直線の等速運動の状態を維持します。この法則は単なる乾いた公式以上のものであり、敗れた異端者に対する勝利の讃歌でもあります。法則から「または一直線の等速運動」の表現を取り除くことで、アリストテレスの反対法則として解釈できる形を得ることができます---全ての物体は強制的な力によってその状態を変えるまで静止状態を維持します。

この最後の誤った公式では、力以外に物体は静止状態を維持すると主張され、従って物体が動いている場合、運動を維持するためには力が必要であると述べられている。したがって、力を停止させると運動も停止する。真のニュートンの法則はこの点から全く反対の立場を取ります。力によって作用されない物体の状態は一直線の均一な運動であり、外部の力や影響を原因として探す必要はありません。または、あなたがそう定義したいなら、この均一な直線運動の不変の伴侶と見做すこともできます。静止はただそのような運動の一例に過ぎず、速度が常にゼロである場合だけです。したがって、物体が動いているときには、速度の率や方向の変化を説明するための外部の影響を探さないで済む場合があります。物体が同じ速度と同一の方向で動いている限り、力の援助を求めることはありません。

二つの視点の違いは、惑星の運動理論を通じてよくわかります。コペルニクスはポーランド人で、西普鲁士トーン出身（1473年生まれ、1543年に没）でした。彼は太陽を中心に近似円軌道上で回る惑星を含め地球も含めて考えると、それほどシンプルな説明になることを示しました。その後、ドイツの数学者ケプラーは1609年に実際にはこれらの軌道がほぼ楕円であることを証明し、これは後の章で詳しく説明する特別な種類の卵形曲線です。

すぐに問題となったのは、これらの力は何によって惑星がこのような運動を保たれているのかという点でした。旧来の誤った見解では、ケプラーは実際の速度そのものが必要な力をもって保存されるべきであると考えていました。したがって、添付図（図4）のように接線方向の力を求めました。しかしニュートンの法則によれば、一部の力以外は惑星は既存の速度で一直線に動くため、太陽から完全に離れてしまうことになります。そこでニュートンは、運動を楕円軌道へ曲げる力が必要であることを探求しました。彼は次のような図（図5）において、その力を太陽に向かう方向であると示しました。実際には、その力は距離の逆平方に従って働く太陽の引力です、これは先ほど述べた通りです。

ギリシャ人によって機械学の科学は、レバーを使用して得られる機械的利得に関する理論から発展し始めました。また、物体の重さに関連する様々な問題も考慮されました。16世紀末から17世紀にかけて、前段落で述べたように、この分野は主に落下する物体の理論を説明するためにも、そして主には惑星運動の科学的な理論を提供するために基礎づけられました。しかし、それ以来、動力学はより大きな任務を引き受け、他の分野とは異なる最終的な科学として自らを位置付け始めました。その主張は、物体の感覚で捉えられる様々な性質は、空間に存在するものの位置変化に対する我々の特殊な評価方法であるという点まで到達しています。

例えば、ウェストミンスター寺院を見ると、それは数世紀間灰色と動かないまま立っています。しかし、現代科学の理論によれば、その灰色が建物の不動性を強調するように感じさせるものは、実際には建築物の外表面を形成し、エーテルに振動を伝える最終的な分子の急速な運動の方法であるだけです。また、手で石を触ると、その冷たさと均一な温度を感じますが、これは建物の静かな安らぎを象徴しています。しかし、この感覚は手から石へ、または石から手への熱の移動に対する我々の認識を示すだけです。そして、現代科学によれば、熱は物体の分子の振動であるとされます。最後に、オルガンが演奏し始めると、音は耳の鼓膜に空気の運動が衝突して生じる結果です。

従って、現象を動力学的な説明で与える試みは、そのようなものや体がこの場所にあり現在は彼方の場所にあるという一般的形式の声明によってそれらを説明することである。これにより我々は現代科学の大切な基本的な概念へと到達する、すなわちすべての感覚は空間における事物の構造の変化との比較の結果である。従って運動の法則、つまり事物の構造の変化の法則は物理学の最終的な法則であることがわかる。

数学を自然哲学の研究に適用するとき、科学は系統的に普通の思考が偶発的に行うことを行います。椅子について話すと、通常は何らかの形で見たり感じたりしたものを指します；しかし私たちの言葉の大半は、視覚や感覚とは独立して何かが存在することを前提としています。一方で、数学的な物理学ではその反対の道筋を歩みます。椅子は特定の誰かや特別な感知モードに関連付けずに考えられます。結果として、思考の中で椅子は空間に位置する分子の集合や電子群、運動中の以太の一部などと見なされます。しかし重要なのは、科学が椅子を空間で動くものとして考え、それらの動きが互いに影響することにまで落とし込むことです。そして、そのような状況に含まれる Various 元素や要素は、空間での物体の位置を決定するための長さ、角度の大きさ、面積、体積のような幾何学的な要素です。もちろん、これらの幾何学的要素の変化があることで運動と変化という事実が導入されなければなりません。つまり速度、角速度、加速度などがあります。したがって、数学的な物理学は、これらの幾何学的要素やその変化率の測度間の関係を表す変動する数値間の関連性について扱います。しかし常に、数学的な法則は変動するものに取り組み、法則を実験によって時々検証したり、特別な予測のために使用するときのみ具体的な数値が代入されます。

この抽象的な方法で数学物理学の研究を通じて世界を概念化した際の興味深い点は、そのような抽象的世界の出来事が「説明」するのに十分であるということです。音を聞くとき、空気中の分子は特定の方式で振動しています：振動を与えると、または呼ばれる空気波がある場合、通常の人々は音を聞きます；そして、空気波がない場合、音はありません。同様に、他の感覚も異なる人々がそれをどのように表現しようと物理的な原因や起源、あるいは並行する出来事（これらはそれぞれの人が好む表現法により異なります）に基礎を置いています。私たちの思考さえも脳の構造と運動に対応しているように見えます：脳を傷つけると、思考も傷つくのです。

この物理的世界の出来事は、特殊な感覚や思考、感情を無視した数学的な法則に従って順次進行します。

現在、数学物理学の世界と私たちの感情、感覚、思考との関係は概ねこのようである。そのような議論が頻繁に行われ、多くの論争と多くの墨が使用されたことは事実である。しかし、一つだけ言及したい点がある。我々が見るように、その全体的な状況は、私た各自的な感覚や感情を説明するための「説明的」な外部世界を描写しようとする試みから生じている。しかしながら、その世界は特定の感覚や個々人によって根本的に依存しているわけではない。

このような世界は単なる大きな fairy tale に過ぎないのか？しかし、fairy tale は幻想的で恣意的なものである。もし実際にはそのような世界が存在するのであれば、それは正確な説明を受けるべきであり、その各部分と相互の関係が正確に定義されるべきである。

現在、科学的世界観はこのテストに相当程度従い、抽象的な数学的概念の装置によってその出来事を探求し予測できる。ここでの初期仮定に対する帰納的検証のようすが確かに見受けられるように思われる。しかし、帰納的証明は結論的ではないことを認めなければならず、もし私た各自的な感覚に依存しない独立した存在を持つ世界の概念自体が誤りであるとするのであれば、そのような数学的概念を用いてその特性を描写しようとする試みが驚くべき成功を収めることの説明が必要である。

他の運動法則について詳細に説明することは、私たちの議論を大きく逸脱するため、省略します。この章の残りは、数学的物理学および純粋数学において基礎的な概念であるベクトル量とベクトル加法の平行四辺形法則の解説に専念することに devoted されます。

運動の本質は、物体が点$A$から点$C$へ移動したことです。この移動は完全に決定されるためには、その大きさ（すなわち距離$AC$）と方向という2つの要素が必要です。このような、大きさと方向によって完全に定義される何かをベクトルと呼びます。例えば、速度はその定義のために大きさと方向の両方を指定する必要があります。それは特定の時間単位あたりの何マイルか、そしてどの方向にあるかで指定されます。速度の大きさと方向が独立して決定されることを示す例としては、船長が異なる下部組織に指示を与える場合があります。彼はエンジニアにノット数を伝え、舵手に航法方位を与えます。

また、速度の変化率、つまり時間単位あたりの速度の増加（加速）もベクトル量です：これは「加速度」と呼ばれます。同様に、動力学的な意味での力もまたベクトル量です。実際には、速度と加速度のベクトル性から、力のベクトル性は動力学的原理によって即座に導けます。しかし、この点については深く掘り下げることはありません。ここでは必要なのは、力が特定の大きさと方向で物体に作用することだけです。

すべてのベクトルはグラフィカルに直線で表すことができます。必要なことは以下の2点です：(i)長さとベクトルの大きさとの単位を対応させるスケールを設定する---例えば、速度の場合には1インチが10マイル/時間に対応し、力の場合には1インチが10トン重に対応すること---そして(ii)図上の直線の方向はベクトルの方向に対応させる。そうすると、適切な長さのインチ数で描かれた直線は任意に設定した大きさのスケール上で必要なベクトルを表します。このベクトルのグラフィカルな表現は非常に重要です。その助けにより、同じ種類のベクトルが異なる方向にある場合でも有名な「平行四辺形則」を述べることができます。

ベクトル$\overrightarrow{AC}$を図6に示すように、物体が点$A$から点$C$へと移動した位置の変化を表すものとして考える：これを輸送ベクトルという。上記で説明したように、物理現象を単なる位置の変化に帰着させることが正しい場合、他のすべての種類の物理ベクトルはいずれかの方法でこの単一のタイプに帰着できる。現在の点$A$から点$C$への最終的な輸送は、点$A$から点$B$へと移動した後、点$B$から点$C$へと移動することでも同様に行われるし、平行四辺形$ABCD$を完成させることで、点$A$から点$D$へと移動した後、点$D$から点$C$へと移動することでも同様に行われる。これらの輸送は順次適用されることにより加算されるという。これは単に輸送ベクトルの加算について定義するためのものである。さらに注意すべきは、並行線を同じ方向で引いた線として考えると、輸送ベクトル$\overrightarrow{BC}$と$\overrightarrow{AD}$はそれぞれの初期位置$B$と$A$に適用された同一の輸送ベクトルと考えることができる。この概念に基づき、任意の位置、例えば点$B$においても、ベクトル$\overrightarrow{AD}$が物体に適用されることを話すことができる。したがって、ベクトル$\overrightarrow{AC}$はベクトル$\overrightarrow{AB}$と$\overrightarrow{AD}$を任意の順序で適用することにより考えることができるということができる。ここで、輸送ベクトルの加算に対する平行四辺形法則がある：つまり、ベクトルが$\overrightarrow{AB}$と$\overrightarrow{AD}$である場合、平行四辺形$ABCD$を完成させ、その二つのベクトルの和は対角線$\overrightarrow{AC}$となる。

すべてこれを見ると、非常に人工的であるように思えるかもしれません。しかし自然そのものがこのような概念を私達に提供しています。例えば、蒸気船が方向$AD$（\cf \Fig[fig.]{6}）を移動しているとしますし、その甲板上で人が歩行しています。蒸気船が静止していた場合、1分で彼は点$B$に到着しますが、その間、甲板上の彼の出発点$A$は点$D$へ移動し、彼の甲板上の経路は$AB$から$DC$へと移動します。つまり実際には、彼の輸送は海面を$A$から$C$へと行われたのです。しかし、これは船体が静止していない場合に、その輸送を2つの輸送の合計として分析されています。すなわち、$A$から$B$への相対的な輸送と、$A$から$D$への蒸気船自体の輸送です。

時間という要素を取り入れると、この人の輸送$AC$の図はその速度を表します。例えば、$AC$が何足の移動を表していたとしても、現在は1分あたりの何足の移動を表しており、つまり彼の速度を表しています。すると$AB$と$AD$はそれぞれ2つの速度を表し、すなわち蒸気船に対する彼の速度と、蒸気船そのものの速度です。これらの「合計」が彼の完全な速度を作ります。明らかに、輸送に関する図や定義は、図が1単位時間あたりの輸送を表していると考えることで、速度に関する図や定義へと変換されます。\PageSep{56}また、速度に関する図や定義は、図が1単位時間あたりに加わる速度を表していると考えることで、加速度に関する図や定義へと変換されます。（\Figure[3in]{7}）

したがって、ベクトル速度の加算とベクトル加速度の加算は、平行四辺形法則に基づく加算を意味します。

また、運動の法則によれば、与えられた質量を持つ物体に生じる加速度を生成するベクトルによって完全に表現される力が存在します。従って、力を加えるとは、その合算効果を平行四辺形の法則に従って計算することを意味します。

科学の基本的なベクトル、すなわち輸送ベクトル、速度ベクトル、および力ベクトルに対しては、同じ種類の任意の2つのベクトルの加算は、平行四辺形の法則に従って生成される「結果ベクトル」を生み出すことになります。

最も単純な平行四辺形は長方形であり、純粋数学においては直角に交わるベクトル$AB$と$AD$に対する単一のベクトル$AC$との関係（図7参照）が頻繁に現れます。$x$、$y$、および$r$単位をそれぞれ$AB$、$AD$、および$AC$の長さとするとき、角度$m$単位を角$BAC$の大きさとするとき、$x$、$y$、$r$、および$m$間の関係は純粋数学において常に繰り返される主題であり、これらの関係は数学物理学の基本的なベクトルへの適用に必要な結果の種類を提供します。この図は、純粋数学の結果が自然の事実に適用されるために渡る主要な橋梁です。

\chapter{第5章：Chapter V: The Symbolism of Mathematics}

我達现在再回到纯粹数学，并且更仔细地考虑构成これ门科学的观念体系。我達首先关注的是これ门科学的符号体系，而我達从最简单且广为人知的符号——算术符号开始。

我々は、アラビック記法で $0$,~$1$, $2$,~\dots, $9$, $10$, $11$,~\dots\Add{,} $100$, $101$,~\dots\ と表される整数について十分な概念を持っていることを仮定しよう。この記法はアラブ人を通じてヨーロッパに導入されたが、彼らはおそらくヒンドゥーの源流からそれを得たものである。最初に系統的に説明された作品は、インドの数学家中の1人が著した「Bhaskara（1114年生）」によるものである。しかし実際の数字符号は私たちの紀元7世紀まで遡ることができるし、おそらく本来チベットで発明されたものかもしれない。しかし現在の我々にとって、記法の歴史は詳細な点に過ぎない。この数字システムが示す興味深い点は、良い表記法の巨大な重要性である。不要な作業から脳を解放することで、良い表記法はその力をより高度な問題に集中させ、実質的に人間の知能を増加させる効果がある。アラビック記法が導入される以前には乗算も難しかったし、整数の除算は最高の数学的才能を要求した。おそらく現代世界でギリシャの数学家が驚くべきことに、強制教育の影響下に西欧の大半の人々が大きな数に対する除算操作ができるようになった事実は、彼にとって絶対的な不可能性に等しいものだったろう。この記法を十進小数まで拡張したのは17世紀になってからである。我々現代の容易な計算力は、十進小数への適切な表記が徐々に発見される結果として、ほとんど不思議な結果である。

数学は多くの人々にとって難解で、神秘的な科学と見なされることが多い。その理由は、数多く用いられる記号によるものである。もちろん、理解できないシンボルはもっと不思議に思えるだろう。また、一部しか理解できず慣れていないシンボルは難しい。同じように、どの職業や仕事にも専門的な技術用語があり、それらは未訓練の者には難解である。しかし、これはその性質自体が難しいからではなく、むしろそれを簡単にするために導入されたものである。数学においても、真剣に数学的概念に取り組んでいる場合、シンボルは必ずしも簡略化するものではないが、それは大きな簡便さをもたらす。実用的なだけでなく、非常に興味深いものでもある。なぜなら、それは概念の分析であり、相互関係をほとんど図的に表現しているからである。もし誰かがシンボルの有用性に疑問を持った場合、以下の代数の基本法則をすべてのシンボルを使わず全文で書き出すことを試みてみることをお勧めする（これらは代数の基本法則を表す）：

\begin{gather*}
x + y = y + x, \Tag{(1)} \\
(x + y) + z = x + (y + z), \Tag{(2)} \\
x × y = y × x, \Tag{(3)} \\
(x × y) × z = x × (y × z), \Tag{(4)} \\
x × (y + z) = (x × y) + (x × z). \Tag{(5)}
\end{gather*}

ここで、式(\ref{1})と式(\ref{2})は加法の交換律と結合法則を表し、式(\ref{3})と式(\ref{4})は乗法の交換律と結合法則を表し、式(\ref{5})は加法と乗法の分配律を表している。例えば、シンボルを使わなければ、式(\ref{1})は「任意の数に二番目の数を足すことは、最初の与えられた数に二番目の数を足すことと同じ結果になる」という意味となる。

この例は、シンボルの助けによって論理を推移させることが、眼でほとんど機械的にできるということを示しています。それ以外の場合には、脳の高い機能が呼び出される必要があることを示しています。

「我々は自分たちが何をしているかを考える習慣を育てることを奨励する」というコピー本や有名な人々のスピーチでは繰り返される箴言は極めて誤ったものだと断言できます。しかし、その逆が正解です。文明は、考えずにできる重要な操作の数を増やすことで進歩します。

シンボルに非常に重要なものを持たせる一つの性質は、それが一目で理解でき、速やかに書き下せることです。シンボルをよりコンパクトに配置する方法としては、直接接続する以外にはありません。よって、良いシンボル表現において重要なシンボルの接続は重要な意味を持つべきです。これはアラビア数字による数値表記の一つの長所でもあります；十つのシンボル $0$,~$1$, $2$, $3$, $4$, $5$, $6$, $7$, $8$,~$9$ と単純な接続によって、どのような数も表現できます。また、代数学において、変動する二つの数 $x$ および $y$ の接続 $xy$ を何を表すかを選択しなければなりません。現在の手元にある最も重要なアイデアは加算と乗算です。数学的には、$xy$ を $x × y$ とするようにシンボル表現をコンパクトにすることを選びました。したがって、上記の法則 \Eq{(3)},~\Eq{(4)}, および \Eq{(5)} は一般的には以下のように書き表されます：
\[
xy = yx,\quad
(xy)z = x(yz),\quad
x(y + z) = xy + xz,
\]
これにより、大きなコンパクト性の利益を得ることができます。同じシンボル表現の規則は定数と変動する数の接続にも適用されます；$3x$ は $3 × x$, $30x$ は $30 × x$ のように書き表します。

明確に述べると、変数に具体的な数値を代入する際には、×の記号を適切に復元することが必要で、アラビア数字の表記と矛盾しないように注意しなければなりません。例えば、x に 2 を y に 3 を代入するとき、xy の代わりに 2 × 3 と書き、23 ではなくします。23 は 20 + 3 を意味します。

また、見かけ上は控えめな記号が科学の発展にとっていかに重要であるかを指摘するのは興味深いことです。そのような記号は重要なアイデアの強調表示を可能にし、しばしば非常に細かいアイデアを表現しやすくします。そしてそれが存在することで、そのアイデアと複雑な思想の流れとの関係性を容易に示すことができます。

例えば、最も控えめな記号である 0 を考えてみましょう。これは数零（ゼロ）を表します。古代世界の多くの数学家中のほとんどは、数零という概念が非常に困惑させたでしょう。なぜなら、それは非常に細かいアイデアで、明らかに直感的ではないからです。量の零についての意味に関する議論は哲学的な著作に多く見られます。実際には、ゼロは他の基数数と比べても、より難しく或いは細かい概念ではありません。1 または 2、または 3 とは何を意味するのか？これらは我々が日常的に使用しているアイデアですが、その基本的な要素について明確な分析ができる人は多くはありません。ゼロの重要な点は、日常生活での操作には必要とされないことです。誰も魚を買うときに零個を購入することはありません。それは文明度の高い基数数であり、その使用は文化的思考の需要によってしか強制されるものではありません。記号 0 は数零を表し、多くの重要な役割を果たします。

The 符号は、その数値表記と密接に関連し、その一部である。実際、その表記法において、桁の値はその位置によって決まります。例えば、数字$25$, $51$, $3512$, $5213$における桁$5$を考えてみましょう。最初の数値$25$では$5$は五を表し、第二の数値$51$では$5$は五十を表し、第三の数値$3512$では$5$は五百を表し、第四の数値$5213$では$5$は五千を表します。したがって、数字五十一をシンボリックな形で書き表すとき、桁$1$は右から左へ数えて第二位に桁$5$を押し上げ、その結果、五十を表します。しかし、独自の五十をシンボリックに表現する必要がある場合、桁$1$がこの役割を果たすことはできません；私たちは単なる合計に何も加えないが、桁$5$を第二位に押し上げるための桁$0$が必要です。これはゼロという符号によって果たされます。おそらくこの目的のためにゼロを導入した人々は、その数字が持つ意味について明確な概念を持っていなかったでしょう。彼らは単に、その桁の位置が合計に何も寄与していないことを表すマークが必要だったのです。ゼロという概念は、$1$, $2$, \dots, $9$といった符号が基数を表すものであるという意味と同様なものであることを理解する必要から徐々に形作られた可能性があります。これは、実用的な便利性によって導入されたシンボリズムによって数学に微細な概念が導入される唯一の例ではないでしょう。

Thus,最初の0の使用は耕作記号を可能にした---それは軽微な役割ではなかった。当時、実用的な人々がこの目的のために0を導入したとき、空想的な考えを嫌う人々は、それを数零と同一視する愚かな習慣を非難した。しかし彼らは間違っていた。このような人々は常にその他の人が調理した食事を咀嚼することから役割を離れたときに間違っている。

この符号0の次の使用は最初の一見して異常なほど単純であり、初学者がその重要性を理解するのに困難がある。簡単な例から始めましょう。第II章で、変数数$x$と$y$の関連性を表す方程式$x + y = 1$について述べました。これは無限に多くの方法で表現できます；例えば、$x = 1 - y$, $y = 1 - x$, $2x + 3y - 1 = x + 2y$などがあります。しかし重要な表現は
\[
x + y - 1 = 0.
\]
同様に、方程式$x = 1$を重要としめす方法は$x - 1 = 0$であり、方程式$3x - 2 = 2x^{2}$を表す重要な方法は$2x^{2} - 3x + 2 = 0$です。ポイントは、すべての変数を表す符号、例えば$x$と$y$、そしていくつかの定数（零以外）を表す符号、例えば上例で見られる1や2が左辺に書かれることであり、左辺全体が零に等しいことである。この最初の手口はトマス・ハリートによって行われたとされています。彼は1560年にオックスフォード生まれで、1621年に亡くなりました。\index{ハリート, トマス}しかし、この単純な符号化の重要性は何でしょうか？それは現代の代数的形の概念の発展を可能にしました。\index{形式, 代数的|EtSeq}\emph{代数的形}という現代的な概念の成長を可能にしたのです。

これは我々が繰り返し取り組むべきアイデアであり、現代の数学のどの部分もこれを頻繁に考えなければ適切には理解できないほど重要です。形式という概念は非常に一般化されており、抽象的な言葉で説明するのは難しいです。この段階では、例を考慮することの方が良いでしょう。例えば、方程式 $2x - 3 = 0$, $x - 1 = 0$, $5x - 6 = 0$ はすべて同じ形式の方程式であり、未知数 $x$ を含みますが、$x$ が自乗されていないので、$x^{2}$, $x^{3}$ 等は現れません。また、$3x^{2} - 2x + 1 = 0$, $x^{2} - 3x + 2 = 0$, $x^{2} - 4 = 0$ もすべて同じ形式の方程式であり、未知数 $x$ を含み、$x \times x$、すなわち $x^{2}$ が現れます。これらの方程式は二次方程式と呼ばれます。同様に、$x^{3}$ が現れる三次方程式も別の形式を生じます。上記の三つの二次方程式の中で、最後の方程式 $x^{2} - 4 = 0$ は前二つの方程式とは微小な違いがあります。これは最後の方程式において $x$（$x^{2}$ と異なる）が現れていないことによるものです。しかし、この差異はこれら三つの方程式がすべて二次方程式であるという大きな事実に比べれば非常に重要ではありません。

次に、二つの変数間の関係を表す方程式の形があります；
例えば、$x + y - 1 = 0$, $2x + 3y - 8 = 0$などです。これらは「線形」と呼ばれる方程式の形式の例です。「線形」の名前の理由は、二変数間の関係をグラフィカルに表現する方法として説明されるように、このような方程式は常に直線で表されます。また、他の二つの変数に対する形式もあります——例えば、二次形式、三次形式などがありますが、ここでは強調したいのは、これらの形式の研究は、等式を右辺に$0$という記号を使用する標準的な書き方によって促進され、実際には可能になることです。

$0$は形式の研究に関連して別の機能も果たします。どのような数値$x$であれ、$0 × x = 0$で、$x + 0 = x$です。これらの性質を利用して、形式的な微細な違いを吸収できます。例えば、二次方程式$x^{2} - 3x + 2 = 0$と$x^{2} - 4 = 0$の上の差異は、後者の方程式を形として$x^{2} + (0 × x) - 4 = 0$と書き換えることで消えます。なぜなら、上記の法則によれば、$x^{2} + (0 × x) - 4 = x^{2} + 0 - 4 = x^{2} - 4$です。したがって、方程式$x^{2} - 4 = 0$は二次方程式の特定のクラスを代表し、$x^{2} - 3x + 2 = 0$と同じ一般的な形式に属します。

これらの三つの理由から、数零を表すシンボル$0$は現代数学において欠かせないものです。この記号がなければ不可能だった種類の研究が可能になりました。

数学のシンボルは、科学全体を支配する一般の概念の結果です。現在私たちは、変数と代数的形式という2つの一般的な概念に直面しています。これらの概念の融合により、数学には独特のシンボルが導入され、その性格は奇妙ですが効果的です。方程式 $x + y - 1 = 0$ は $x$ と $y$ の一意的な関連を表し、また $3x + 2y - 5 = 0$ は別の一意的な関連を表します。これら両者は線形関係の形式を持っています。

しかし、変数 $x$ と $y$ の任意の線形関係をどのように表現するでしょうか？ここでは、任意の線形関係をシンボル化したいと考えます。これには、特定の関係 $3x + 2y - 5 = 0$ に現れる数値を文字に変えることで達成されます。結果として、$ax + by - c = 0$ を得ます。

ここで $a$, $b$, $c$ は $x$ と $y$ のように任意の変数を表しますが、その使用方法には違いがあります。私たちは $x$ と $y$ の関係の一般的な性質を研究する一方で、$a$, $b$, $c$ の値は固定されています。これらの値を決定することはありませんが、どのような値であれ、私たちは $x$ と $y$ の全可能性のある値に対する関係を研究します。

しかし、この関係の性質を得た後で、$a$, $b$, $c$ の実際の値は定められていないことに注意すると、証明した性質は任意のそのような関係に属することになります。つまり、$a$, $b$, $c$ を変えることで、$ax + by - c = 0$ が $x$ と $y$ の変動する線形関係を表すという考えに至ります。

$x$ と $y$ に対しては、三つの変数 $a$, $b$, $c$ は定数と呼ばれています。このように使用される変数は時としてパラメータとも呼ばれます。

現在、数学家用慣れており、変数を「定数」として扱うものと、「変数」として扱うものを説明する手間を省くために、アルファベットの末尾の文字で「変数」の変数を表し、アルファベットの頭文字で「定数」やパラメータを表します。この2つのシステムはアルファベットの中間付近で自然に合流します。時折説明が必要な単語がいくつかありますが、実際には慣習と普通の感覚が十分であり、見た目ほど緩慢な Procedure によって生じる混乱は意外と少ないです。

このパラメータの層々による具体的数値の継続的な排除の結果として、数学家用いる算術量は非常に少なくなっています。多くの数学家用一切の数値計算を嫌い、特にその技術に優れていません。算術の領域は「変数」という概念と「代数的形式」の概念が支配するところまでです。

\chapter{第6章：Chapter VI: Generalizations of Number}

数学の大きな特徴の一つは、整数から始まった概念群が発明されたことです。これらの概念は数を拡張または一般化したものと呼びましょう。まず最初に Fraction（分数）という概念があります。我々が所有する最も古い算術に関する論文は、紀元前1700年から紀元前1100年にかけてエジプトの司祭であるAhmesによって書かれました\index{Ahmes} おそらくもっと古い著作のコピーです。この論文は主に分数の性質について扱っています。したがって、この概念は数学の歴史において非常に早期に発展したことがわかります。実際、その主題は非常に自然なもので、田地を3等分し、そのうち2つを取り出すという操作はしばしば行われたでしょう。そのため、遠い文明の人々が三分の二という概念を熟知していたのは驚くことではありません。PageSep{72} そこで最初に数の一般化として分数の概念を置きます。ギリシャ人はこの主題を比という形で考えていたため、ギリシャ人にとって2フィートの線と3フィートの線との比は2対3であると言ったでしょう。我々の代数学記号の影響を受けた場合、一方の線が他方の三分の二であると言った方が自然であり、三分の二是数値的な乗数と考えることになります。

ギリシャ人は比の理論や分数に関連して、大きな発見を行いました。この発見は哲学的な考えだけでなく、数学的な思考にも多くの影響を与えました。彼らは「無比値」の比の存在を発見しました。実際、彼らの幾何学的調査の中で、任意長さの線から他の線が存在することを証明しました。これらの線の長さは元の長さに対する整数ペアの比で表すことができません——言い換えれば、これらは元の長さの正確な分数ではない長さです。

例えば、正方形の対角線は同じ正方形の辺の任意の分数で表現することはできません。現代の記法では、対角線の長さは辺の長さの $\sqrt{2}$ 倍です。しかし、$\sqrt{2}$ という値を正確に表す分数はありません。$\sqrt{2}$ の値に近づくことはできますが、絶対的にその値には到達しません。例えば、$\dfrac{49}{25}$ は $2$ より少々小さい一方で、$\dfrac{9}{4}$ は $2$ より大きいので、$\sqrt{2}$ は $\dfrac{7}{5}$ と $\dfrac{3}{2}$ の間にあることがわかります。$\sqrt{2}$ の値に近づく最も systematic な方法として、小数分数の列を生成するには、平方根を通常の方法で求めるのが最良です。したがって、その列は $1$、$\dfrac{14}{10}$、$\dfrac{141}{100}$、$\dfrac{1414}{1000}$ などとなります。

このような比はギリシャ人によって「無比値」と呼ばれています。それらはギリシャ時代から現在まで、大きな哲学的な議論を引き起こし、その困難性は最近まで解消されていませんでした。

無比比例を分数と並べ、整数、分数、無比比例の全体を一つの数のクラスとして考えます。これを「実数」と呼びます。常に実数を大きさの順序で配置し、ゼロから上に進むにつれて無限大になるように考えます。実数は線上的な点で便利に表現できます。任意の線$OX$を選び、$O$から始まり、$OX$方向に無限に延びる線として定義します。その上で、適当な点$A$を選んで$OA$を1単位長さとします。次に、$AB$, $BC$, $CD$などの各辺をそれぞれ$OA$と同じ長さで分けていきます。このとき、点$O$は数$0$、$A$は数$1$、$B$は数$2$を表します。実際、任意の点が表す数はその点から$O$までの距離（$OA$を単位長さとして）の測定値です。$O$と$A$間の点は真分数や$1$未満の無比比例を、$OA$の中点は$\dfrac{1}{2}$、$AB$の中点は$\dfrac{3}{2}$、$BC$の中点は$\dfrac{5}{2}$を表します。このようにして$OX$上の任意の点が一つの実数を表し、逆に任意の実数も$OX$上の特定の一点で表現されます。

$OX$上から$O$を起点として、$O$から$X$に向かって規則的に配置された点の系列は、大きさの順序で昇順に並べられた実数を表します。

すべてこれだけは単純に思えるが、この段階で興味深いアイデアを抽出することができる。整数を表す点の系列を考える。つまり、点$O$, $A$, $B$, $C$, $D$, 等である。ここでは最初の点$O$があり、次に続く明確な点$A$があるし、各点（例えば$A$や$B$）は一方には明確な直近の前駆点、他方には明確な直近の後継点を持つが、点$O$には前駆点がない。また、この系列は無限に延び、終了点がない。このような順序は整数の型の順序と呼ばれ、その本質は行の中の1番目の点以外の各点が両隣に近接した先駆点と後継点を持つことにある。

次に、整数と分数を一緒に考え、無理数に対応する点を省略する。このとき得られる系列の順序は全く異なるものである。最初の項$O$があるが、どの項も明確な直近の前駆点や後継点がない。これは容易に確認できる。例えば、任意の2つの分数の間に常に別の値の中間の分数を見つけることができる。その一つは分数を足し合わせてその結果を半分にする方法である。例えば、

\[
\frac{2}{3} \text{と}\frac{3}{4} の間に }\frac{17}{24} = \frac{1}{2} \left( \frac{2}{3} + \frac{3}{4} \right) \text{ が存在し、}
\]

\[
\frac{2}{3} \text{と}\frac{17}{24} の間に }\frac{33}{48} = \frac{1}{2} \left( \frac{2}{3} + \frac{17}{24} \right) \text{ が存在し、}
\]

このような無限に続く。この性質から、この系列は「密接」であると呼ばれる。系列には終点がないが、線$OX$を進めるにつれて無限に増大する。最初の見方では、分数から整数を含む系列とすべての実数（整数、分数、無理数を含む）から得られる系列は同じ型であるように思えるかもしれない。しかし、重要な違いがあり、それが今後開発される。

分数の系列に無理数が含まれていないことで、特定のクラスには端点がない状態が生じる。例えば、無理数$\sqrt{2}$を考える。実数の系列では、これは全ての平方が$2$未満の数と$2$より大きい数を分ける位置にある。しかし、分数だけの系列に限定し、無理数を考慮しない場合、$\sqrt{2}$を導入できないので、どの分数もそのような端点を持つ性質を持たない。したがって、分数の系列には$\sqrt{2}$の位置に類似の欠陥がある。

このように、分数の系列に欠けている端点は小さな問題のように思えるかもしれない。しかし、数列の一部が全体の数列を覆わない場合でも極限や最大値の存在が保証されないことは大きな欠点であることを、どの数学家も知っている。無理数を使用することで、系列に欠けた部分を埋め、完全な系列を得ることができるため、この困難を避けることができる。

二つの数列にはさらに根本的な違いがあります。整数の数列のように、分数の項も最初の項があり、各項は直ちな後続項と（最初の項以外では）直じな前駆項を持てます。このように並べ替えられることが示すことができます。分数や整数の系列において、すべての項を分数形式で書き表します。1を$\frac{1}{1}$、2を$\frac{2}{1}$とし、他の整数も同様に表します（0は除きます）。また、一時的に等しい値を持つが約分されていない分数を区別するようにしますので、例えば$\frac{2}{3}$、$\frac{4}{6}$、$\frac{6}{9}$、$\frac{8}{12}$などはすべて区別されます。

次に、各項の分子と分母の和で分数をクラス分けします。この和を分数の指標と呼びます。例えば、$\frac{4}{3}$、$\frac{3}{4}$、$\frac{2}{5}$の指標は7です。各クラスには特定の指標を持つすべての分数が含まれます、つまりそのクラスの指標とも呼ばれることができます。これらのクラスを指標の大きさ順に並べ替えます。最初のクラスは指標2で、その唯一のメンバーは$\frac{1}{1}$です；第二のクラスは指標3で、そのメンバーは$\frac{1}{2}$と$\frac{2}{1}$です；第三のクラスは指標4で、そのメンバーは$\frac{1}{3}$、$\frac{2}{2}$、$\frac{3}{1}$です；第四のクラスは指標5で、そのメンバーは$\frac{1}{4}$、$\frac{2}{3}$、$\frac{3}{2}$、$\frac{4}{1}$です。そして続くようにします。

どのクラスにも属する分数（まだ約分されていない分数も含む）の数は指標より一つ少ないことが容易にわかります。また、各クラスのメンバーを指標が$n$であるクラスの場合、分子が1から順に増えていく順序で並べることができます。つまり、クラスの指標が$n$の場合、メンバーは次のようになります。
\[
\frac{1}{n - 1},\quad
\frac{2}{n - 2},\quad
\frac{3}{n - 3},\
\dots,\quad
\frac{n - 1}{1}.
\]
最初の四つのクラスのメンバーはこの順序で挙げられています。したがって、すべての分数は整数の並びのように並べ替えられました。その並びは次のようになります。
\begin{gather*}
\frac{1}{1},\
\frac{1}{2},\
\frac{2}{1},\
\frac{1}{3},\
\left[\frac{2}{2}\right],\
\frac{3}{1},\
\frac{1}{4},\
\frac{2}{3},\
\frac{3}{2},\
\frac{4}{1},\  \dots, \\
\frac{n - 1}{1},\
\frac{1}{n - 1},\
\frac{2}{n - 2},\
\frac{3}{n - 3},\
\dots,\
\frac{n - 1}{1},\
\frac{1}{n},
\end{gather*}
そして続くようにします。

現在、同じ値を持つ分数の重複をすべて取り除くためには、最初の出現以降でその分数が現れた場合はそれを消去すればよいです。上記に書かれた少数の初期項においては、$\frac{2}{2}$（この分数は上に四角括弧で囲まれています）のみが既約形ではありません。これは以前$\frac{1}{1}$として出現していますから、これを消去しなければなりません。しかし、系列は依然として以下の性質を保っています：(\textit{a})最初の項がある、(\textit{b})各項には隣接する項がある、(\textit{c})系列は無限に続く。

このように実数全体の系列をこのような方法で並べることは不可能であることが証明できます。この奇妙な事実は、ドイツの数学者ゲオルク・カントールによって発見されました；彼はまだ生きており、これは数学的概念の哲学において極めて重要です。ここでは実際には連続性と無限に関する大きな問題に触れています。

数の拡張のもう一つは、操作やステップという概念の導入から来ます。
\index{Steps|EtSeq}\%
この概念は、異なる視点からやや異なる意味で適切です。まず特定の場合から始めましょう。$2 + 3 = 5$という文を考えます。$2$に$3$を加えると$5$が得られます。$3$を加える操作を考えましょう：これを$+3$で表します。また$4 - 3 = 1$も考えます。$3$を引く操作を考えましょう：これを$-3$で表します。したがって、実数そのものではなく、それらを加えたり引いたりする\emph{操作}を考えます：$\sqrt{2}$ではなく、$+\sqrt{2}$と$-\sqrt{2}$を考えます、つまり$\sqrt{2}$を加える操作と$\sqrt{2}$を引く操作です。これらの操作は当然、数を加えるのとは異なる意味での加算で足すことができます。二つの操作の和は、その二つを順次適用したときと同じ効果を持つ単一の操作となります。二つの操作をどのように順に適用するかについては、答えは関係ありません。例えば
\[
2 + 3 + 1 = 2 + 1 + 3;
\]
すなわち、ステップ$+3$と$+1$の加算は可換です。

数学家中には、意味を追跡する人には困惑する習慣がありますが、実践では非常に便利なのが、似たような意味を持つ異なる文脈で同じ記号を使用することです。彼らの目には、どのような可能性のある意味に関わらず、その記号の使用に対する形式的な法則が常に同じであることが唯一の必須条件です。

この習慣に従って、操作の追加は$+$と数値の追加でも表されます。したがって、次の式を書くことができます：
\[
(+3) + (+1) = +4;
\]
ここで、左辺の中間の$+$は操作$+3$と$+1$の追加を示します。

しかし、シンボルにおいては、直接意味を追跡する稀な例外以外では、あまり厳密にしなくても良いです。そこで、行の最初の$+$と括弧は省略し、二つの$+$の連続は書かないことにします。したがって、上の式は次のようになります：
\[
3 + 1 = 4,
\]
これは単純な数値追加を意味するか、または前の書き方で完全に表現される操作の追加を意味するか、最後には操作$+1$を数$3$に適用して得られる数$4$を表すか、どの解釈も常に正しいです。しかし、特定の条件下では常に可能な唯一の解釈は操作の解釈です。他の解釈はしばしば意味不明な結果を生み出します。

これは即座に一つの質問へと導きます。読者の心の中で必ずしも押し寄せてきたことでしょう：このような詳細な説明は何のために必要なのでしょうか？この場で実用的な人間という友人がきっと立ち上がり、頭の中のこれらの馬鹿げた蜘蛛の巣を一掃するよう主張することでしょう。答えは、数学家用く disproportion しているのは普遍性であるということです。これは変数や形式と並んで重要性において数学的手法を支配する概念として価値があります。\index{Generality in Mathematics}\index{Form, Algebraic}\index{Variable, The}これら三つの概念、すなわち変数、形式、そして普遍性は、数学の全体的な主題を統治する一種の数学的三位一体を形成します。これらすべては実質的に抽象科学の本質から派生しています。

この概念の導入によって一般性がどのように獲得されるかを見てみましょう。方程式 $x + 1 = 3$ を考えてみます；その解は $x = 2$ です。ここでは、シンボルを単なる数として解釈できるため、「操作」への依頼は全く必要ありません。しかし、もし $x$ が単なる数である場合、方程式 $x + 3 = 1$ は意味のないものになります。なぜなら、$x$ は 1 個のものを 3 個取り除いた後で残るものの数であり、そのような手続きは不可能だからです。この点で、代数学的な式の概念が登場します；これは別の観点から見た一般化に過ぎません。したがって、同じ形式 $x + 1 = 3$ の一般方程式を考えます。この方程式は $x + a = b$ であり、その解は $x = b - a$ です。ここでの困難は激しくなります；この形式は $b > a$ の限りで数値的解釈が可能ですが、未定形の $a$ と $b$ がどのような定数でもよいとは断定できません。言い換えれば、「定数」 $a$ と $b$ の変動性に制限を設けなければならず、これが私たちの議論を通じて鎖のように引きずらなければなりません。このような条件下では、実際的な数学的調査は不可能です。すべての方程式は最終的に無数の制限の下に埋もれてしまいます。しかし、もし現在のシンボルを「操作」と解釈すれば、すべての制限が魔法のように消え去ります。方程式 $x + 1 = 3$ は $x = +2$ となり、方程式 $x + 3 = 1$ は $x = -2$ となり、方程式 $x + a = b$ は $x = b - a$ となり、これは加法または減法の操作となります。$b - a$ が加法の操作を表すか減法の操作を表すかは決める必要はありません；シンボルの操作規則はどちらの場合も同じです。

この著作の計画には、初等代数に関する詳細な章を書くことは含まれていません。私たちの目的は単に科学の基礎的な概念を明確にするところまでです。そのため、正負の数がどのように乗算や他の組み合わせを行うかについての詳細なルールをさらに説明することはありません。

上記で述べたように、正負の数は操作であり、「ステップ」とも呼ばれます。したがって $+3$ は $2$ から $5$ へのステップであり、$-3$ は $5$ から $2$ への逆方向のステップです。

線分 $OX$ を先に説明したように分割し、その点が数を表すようにします。すると $+2$ は $O$ から $B$ へのステップ、または $A$ から $C$ へのステップ、もしくは $OX'$ 上で分割を逆方向に取ると $C'$ から $A'$ へのステップ、または $D'$ から $B'$ へのステップです。同様に $-2$ は $O$ から $B'$ へのステップ、または $B'$ から $D'$ へのステップ、または $B$ から $O$ へのステップ、または $C$ から $A$ へのステップです。

ステップから $O$ に達する点をそのステップの代表とすることができます。したがって $A$ は $+1$ の代表、$B$ は $+2$ の代表、$A'$ は $-1$ の代表、$B'$ は $-2$ の代表であり、類推します。これまでのところ、単なる「未署名」実数においては $O$ の一方側にある点のみが数を表していましたが、現在ではステップに関しては $O$ の両側にわたるすべての点がステップを表すようになっています。これは正負の数によって導入された優れた一般性の視覚的な表現です、つまり操作やステップという概念です。これらの「符号付き」数はベクトル（ラテン語で \Foreign{veho} から派生した「運ぶ」という意味）の特別な例でもあります。なぜなら、粒子が $O$ から $A$ へ、または $A$ から $B$ へ移動するという考え方が可能だからです。

前にも述べたように、実用的な人間が正の数と負の数の導入に含まれる微妙さを不快に感じるだろうという提案は、その優れた個人を中傷していた。実際、ここでは彼の最大の勝利の場面に立っている。もし真実を告げなければならなかったら、実用的な人間自身が実際にプラスとマイナスの符号を使用した最初の人だった。これらの符号の起源は明確ではないが、おそらくドイツの倉庫で商品の箱にchalked（粉筆で書いた）マークとして使用されたものから発展したと考えるのが最も妥当だろう。それらの最初の記録は1489年レipzigで出版された本に現れる。これらの符号が数学において初めて使用されたのは、1544年にNuremburgで出版されたStifelというドイツの数学者によって行われたようだ。しかし、最近になってからこそ、ドイツ人を実用的な国の人々と見なされるようになった。古いエピграммによれば、海はイギリス人の帝国、陸地はフランス人の帝国、そして雲はドイツ人の帝国だとされている。確かにこれらの符号が生み出したアイデアは人類の福祉にとって極めて重要であるため、それは海や陸からではなく、雲からドイツ人が取り寄せたものだと考えるべきだろう。

-positiveおよび-negativeの数の応用可能性は非常に明白です。
一方の方向に沿った長さを正の数で表すなら、反対方向の長さは負の数で表されます。
一つの方向に対する速度が正であれば、反対方向の速度は負となります。
時計の針と逆向き（反時計回り）での円盤の回転が正であれば、時計回りの回転は負となります。
銀行での余剰金が正であれば、引落しが負となります。
ガラス体の electrification が正であれば、樹脂体の electrification は負となります。
実際には後者の場合、正 electrification および負 electrification という用語自体が他の用語をほとんど駆逐してしまいました。
無数の例を挙げることができます。
正の数と負の数の概念は数学的な精巧さの中で最も成功したものです。

\chapter{第7章：Chapter VII: Imaginary Numbers}

もし前章で扱われた数学的概念が人気を博したのであれば、本章で扱われる概念はほぼ同じ程度の一般の注目を集めました。しかし、その成功は異なる形を持ち、フランス語でいうところの「succès de scandale」となりました。実用的な人々だけでなく、文学者や哲学者も、数学家が神秘的な存在であると認めながらも虚想された数（Imaginary Numbers）に没頭することについて困惑しました。

この点では、ある種の知性は技術用語の適用性について自己および他人を悩ませることがあります。無比整な数は正しく数と呼ばれるべきでしょうか？正負の数は本当に数であるのでしょうか？虚数は虚想されたものであり、数であるのでしょうか？---このような無意味な質問が代表例です。

科学においては、技術用語は子供に与えられるクリスチャン名のように、任意で適当に選ばれる名前です。これらの名前の正しさや誤りについて疑問を呈することは問題ありません。これらは賢明であるか不賢明であるかの問題であり、しばしば覚えやすくするためや関連する重要なアイデアを提示するために配置されることがあります。しかし、その基本的な原理は、Humpty DumptyがAliceに対して言及したように明らかに述べられています。「私はそれらの用語に特別な料金を支払い、私が好きなように意味を持たせます」と。

そこで私たちは虚数が虚想されたものであるかどうか、またはそれが数であるかどうかについて心配する必要はありません。その表現は特定の数学的概念の任意の名前であり、現在それを明確にする努力を払います。

The概念的起源与正负数的概念完全相似。これ种观念同样源于三个伟大的数学思想：变量、代数形式和一般化。正负数是从考虑方程式如$x + 1 = 3$，$x + 3 = 1$以及一般形式$x + a = b$中产生的。类似地，虚数的起源也归因于方程式如$x^{2} + 1 = 3$，$x^{2} + 3 = 1$和$x^{2} + a = b$。同样的过程被经历。方程式$x^{2} + 1 = 3$变为$x^{2} = 2$，これ有两个解，即$x = +\sqrt{2}$或$x = -\sqrt{2}$。存在これ些替代解的陈述通常写为$x = ±\sqrt{2}$。到目前为止一切都很顺利，就像之前的情况一样。但此时出现了类似的困难。に対して方程式$x^{2} + 3 = 1$与える$x^{2} = -2$，而没有正数或负数可以在其自身乘以时产生一个负平方。因此，如果我達的符号表示普通的正负数，则方程式$x^{2} = -2$无解，实际上これ个方程式是无意义的。于是最终考虑一般形式$x^{2} + a = b$，我達发现当且仅当$b \geq a$时，其解为$x = ±\sqrt{(b - a)}$。因此我達不能无限制地说“常数”$a$和$b$可以是任何数字；也就是说，“常数”$a$和$b$并不是它達应该成为的独立无限制的“变量”。随着工作的进行，これ将导致一系列限制和约束累积起来。

同じ前のタスクが私たちに待ち構えています：方程式 \(x^2 + a = b\) の解 \(\pm\sqrt{b - a}\) が常に意味を持つように、我々の記号に新しい解釈を与える必要があります。言い換えれば、\(\sqrt{a}\) が \(a\) が正でも負でも常に意味を持つように解釈する必要があります。

当然のことながら、その解釈は加法、減法、乗法、除法の通常の形式的な法則を満たすものでなければなりません。また、有理数と無理数を使用して得られた一般性に干渉しないことも求められます。実際には、その解釈は有理数と無理数を含む特別な場合として包含する必要があります。

\(a\) が負のとき、それを \(-c^2\) と書き表すことができますから \(c^2\) は正です。すると、

\[
\sqrt{a} = \sqrt{(-c^2)} = \sqrt{\{(-1) \times c^2\}} = \sqrt{-1} \sqrt{c^2} = c\sqrt{-1}.
\]

したがって、我々の記号を解釈する方法を見つけることができれば、\(\sqrt{-1}\) に意味を与えることができれば目的は達成されます。つまり、\(\sqrt{-1}\) はすべての虚数量の頭と先頭として見られています。

\(\sqrt{-1}\) の解釈を見つける仕事は、\(-1\) の解釈よりも遥かに難しい作業です。実際、比較的簡単な問題はすぐに直感的に解決されましたが、最初はすら最大の数学者たちにとっても問題が存在し得る可能性があるということがほとんど思い浮かばなかったのです。方程式 \(x^2 = -3\) が現れたときには、それらは単に無意味とされて側線に置かれました。

しかし、18世紀には、さらにはそれ以前にも、これらの無意味な記号に意味を付与する解釈が割り当てられることのいかに便利であるかが徐々に認識されるようになりました。これらの記号に対する形式的な論理は、これらが通常の代数変換法則に従うことを仮定して行われました。そして、これらの記号を使用することが正当化されれば、興味深い結果が得られることを見出しました。多くの数学者はその手続きの論理について当時まだ明確ではなかったため、適切な操作によって意味のない記号が虚幻的な方法で命題の正当な証明を生むという考えが広まりました。これは全く誤りです。正しく定義されていない記号はただの紙に書かれた形の blot に過ぎません。 blot の連続によって何らかの主張が証明されるのは、悪いペンの存在や粗雑な筆者の存在を示すだけです。

この時代には、「虚数」という言葉が $\sqrt{(-1)}$ に使われるようになりました。これらの数学者は実際には、$\sqrt{(-1)}$ とその加減乗除に対する解釈が存在し、通常の代数法則（たとえば $x + y = y + x$ 等）を満足させるなら、特定の結果が導かれるという一連の仮定的な命題を証明しました。そのような命題は次の形で空白になっています：「$\sqrt{(-1)}$ とその加減乗除に対する解釈が存在し、通常の代数法則（たとえば $x + y = y + x$ 等）を満足させるなら、このような結果が導かれる。」数学者は常にその大きな「仮定」が結果の前にすべきだったことを理解していなかったのは自然でした。

如予想されるように、その解釈は負の数に対するものよりもずっと複雑なものでした。読者の注意を引くために、いくつかの慎重な前駆的な説明が必要です。既に、点を二つの数で表す方法について触れました。正の数と負の数を利用することで、平面内の任意の点の位置は一対のこのような数で表現できます。したがって、直線 $XOX'$ および $YOY'$ を直角に交差させたものを「軸」として取り、すべての測定からここから始めるものとします。$OX$ 方向および $OY$ 方向の長さは正であり、$OX'$ 方向および $OY'$ 方向の長さは負です。

例えば、順に書かれた一対の数 $(+3, +1)$ が存在し、この一対には最初の数（上記例では $+3$）と二番目の数（上記例では $+1$）があります。$(+3, +1)$ は $O$ から $XOX'$ 方向に $+3$ の長さを測定し、$YOY'$ 方向に $+1$ の長さを測定する指示を与えます。したがって、$(+3, +1)$ では $O$ から $X$ 方向に $3$ 単位の長さを測定し、$O$ から $Y$ 方向に $1$ 単位の長さを測定します。同様に、$(-3, +1)$ では $O$ から $X'$ 方向に $3$ 単位の長さを測定し、$O$ から $Y$ 方向に $1$ 単位の長さを測定します。また、$(-3, -1)$ ではそれぞれ $OX'$ 方向と $OY'$ 方向に二つの長さを測定し、$(+3, -1)$ ではそれぞれ $OX$ 方向と $OY'$ 方向に二つの長さを測定します。このような一対の数を「順序対」と呼びましょう。すると、$1$ および $3$ の二つの数から八つの順序対が生成されます：

\begin{gather*}
(+1, +3),\ (-1, +3),\ (-1, -3),\ (+1, -3), \\
(+3, +1),\ (-3, +1),\ (-3, -1),\ (+3, -1).
\end{gather*}

これらの八つの「順序対」は、他のどのものよりも異なる測定プロセスを指示します。

最後の四組の順序-payi-が表す測定プロセスは、図で示されています。$OM$と$ON$の長さの合計は$(+3, +1)$に対応し、$OM'$と$ON$の長さの合計は$(-3, +1)$に対応し、$OM'$と$ON'$の長さの合計は$(-3, -1)$に対応し、$OM$と$ON'$の長さの合計は$(+3, -1)$に対応します。これらの各矩形を完成させることで、点$P$が順序-payi-$(+3, +1)$によって完全に決定され、逆もまた真であり、同様に点$P'$が$(-3, +1)$、点$P''$が$(-3, -1)$、点$P'''$が$(+3, -1)$をそれぞれ決定します。より一般的には、前の図（Fig. 8）において、点$P$は順序-payi-$(x, y)$に対応し、図中の$x$と$y$はともに正であると仮定されます。点$P'$は$(x', y)$に対応し、図中の$x'$は負と仮定され、点$P''$は$(x' y')$に対応し、点$P'''$は$(x, y')$に対応します。したがって、任意の正または負の数$x$および$y$を持つ順序-payi-$(x, y)$とその対応する点は互いに決定しあいます。この時点では、いくつかの名前を導入することが便利です。順序-payi-$(x, y)$において最初の数字$x$は対応する点の「 абсцисса」と呼ばれ、第二の数字$y$は点の「 ордината」呼ばれます。そして、この二つの数は点の「 座標」と呼ばれます。点の位置を「 座標」で決定するという考え方は、「 虚数」理論が形成されるときには新しいものではありませんでした。これはフランスの大数学者兼哲学者デカルトによって提唱され、1637年リーデンで出版された彼の『Discours』に現れています。順序-payi-を独立したものの概念として捉える考え方は後から生まれたものであり、可能な限り抽象的に虚数を解釈する努力の結果です。

さらにこの順序対の概念についての例として、図9（\Fig[fig.]{9}）の点$M$は順序対$(+3, 0)$、点$N$は順序対$(0, +1)$、点$M'$は順序対$(-3, 0)$、点$N'$は順序対$(0, -1)$、点$O$は順序対$(0, 0)$であることに注意できます。

順序対$(x, y)$を表現するもう一つの方法として、この順序対を点$P$ではなく点$P$を表す点線$OP$（\Chg{cf.}{\cf}図8参照、\Fig[fig.]{8}）と考えることができます。したがって、順序対は特定の長さと方向を持つ「原点」$O$から引かれた線を表します。線$OP$は$O$から$P$へのベクトル線またはステップと呼びます。したがって、この章では以前に与えた正の数と負の数に対する解釈を拡張していることがわかります。ベクトルによる表現法は、順序対の加法や乗法の意味を考える際に非常に有用です。

私たちはこの問題に取り組み、順序対 $(x, y)$ と $(x', y')$ の和の意味をどのように定義するのが便利かを尋ねます。その解釈は、(a) 結果がまた別の順序対であることを、(b) 操作が可換であること、すなわち $(x, y) + (x', y') = (x', y') + (x, y)$ であり、(c) 操作が結合法則を満たすこと、すなわち
\[
\{(x, y) + (x', y')\} + (u, v) = (x, y) + \{(x', y') + (u, v)\},
\]
(d) 減算の結果が一意であること、つまり方程式
\[
(x, y) + (a, b) = (c, d)
\]
を満たす未知の順序対 $(x, y)$ を求めるとき、その答えは一意であり、それを
\[
(x, y) = (c, d) - (a, b)
\]
で表すことができる、という条件を満たします。これらの要件は、$(x, y) + (x', y')$ を順序対 $(x + x', y + y')$ と定義することで満たされます。従って定義により
\[
(x, y) + (x', y') = (x + x', y + y')
\]
とします。

ここで、等式の左辺の $+$ の意味は我々が現在定義する新しい意味であり、右辺の二つの $+$ は前章で定義した正負数の加法（操作）の意味です。この双方向の使用からは実質的な混乱は生じません。

加算の例として
\begin{align*}
(+3, +1) + (+2, +6) \&= (+5, +7), \\
(+3, -1) + (-2, -6) \&= (+1, -7), \\
(+3, +1) + (-3, -1) \&= (0, 0).
\end{align*}

現在、減算の意味が定まりました。それは
\[
(x, y) - (u, v) = (x - u, y - v)
\]
であり、したがって
\[
(+3, +2) - (+1, +1) = (+2, +1),
\]
また
\[
(+1, -2) - (+2, -4) = (-1, +2),
\]
さらに
\[
(-1, -2) - (+2, +3) = (-3, -5).
\]

容易に確認できるように、
\[
(x, y) - (u, v) = (x, y) + (-u, -v)
\]
であり、また
\[
(x, y) - (x, y) = (0, 0)
\]
です。したがって $(0, 0)$ は零順序対と見なされるべきで、例えば
\[
(x, y) + (0, 0) = (x, y).
\]

図形的な表現による順序対の加算は驚くほど簡単です。
\Figure{10}

Let $OP$ represent $(x, y)$ so that $OM = x$ and $PM = y$; let $OQ$ represent $(x_{1}, y_{1})$ so that $OM_{1} = x_{1}$ and $QM_{1} = y_{1}$. Complete the parallelogram $OPRQ$ by the dotted lines $PR$ and~$QR$, then the diagonal $OR$ is the ordered couple $(x + x_{1}, y + y_{1})$. For draw $PS$ parallel to~$OX$; then evidently the triangles $OQM_{1}$ and $PRS$ are in all respects equal. Hence
\[ MM' = PS = x_{1} \]
and
\[ RS = QM_{1} \]
and therefore
\begin{gather*}
OM' = OM + MM' = x + x\\textit{{1}, \\
RM' = SM' + RS = y + y}{1}.
\end{gather*}

Thus $OR$ represents the ordered couple as required. This figure can also be drawn with $OP$ and $OQ$ in other quadrants.

ここでは直行四角形の法則に再び帰着していることが明らかである。
\index{直行四角形の法則}\%
これは第六章で運動の法則として velocitie と force に対して適用されることを述べたものである。記憶に留めておくとよいが、$OP$ と $OQ$ が二つの velocity を表す場合、粒子は $OR$ の velocity と同じ velocity で動いていると定義される。つまり $OR$ は $OP$ と $OQ$ の二つの velocity の結果であるという。また、body の一点に作用する force も lines によって表現できるし、同じ直行四角形の法則が成り立つ。すなわち、$OP$ と $OQ$ の二つの force の結果は対角線 $OR$ によって表される力である。したがって、順序対を velocity や force を表現するものとして見做すことができる。我々が先に述べた順序対の加算の規則は、force および velocity の加算の基本法則を表している。数学の最も魅力的な特性の一つは、異なる部分でのアイデアと結果が驚くべき方法で整合的に繋がることである。この章と前章での議論では、最も抽象的な純粋数学的考慮事項に従って進行したが、結局のところ自然法則の中でも最も根本的な法則へと導かれた。これらはエンジニアが機械を設計するときや、海軍建築家が船舶の安定性を計算するときに心に置くべき法則である。最も理論的な思考状態においても、我々は最も実用的な応用に近い存在であるという主張は矛盾しない。

\chapter{第8章：Chapter VIII: Imaginary Numbers (Continued)}

順序対の乗法の定義は、その加法の定義と同様の考慮事項に基づいています。乗法の解釈は次のようであるべきです。

(α) 結果は別の順序対となるべきです，

(β) 操作は可換性を持つので、
\[
(x, y) × (x', y') = (x', y') × (x, y)
\]

(γ) 操作は結合法則を満たし、すなわち
\[
\{(x, y) × (x', y')\} × (u, v) = (x, y) × \{(x', y') × (u, v)\}
\]

(δ) 除算の結果が一意的であることを保証する必要があります（ゼロ対 $(0, 0)$ の場合を除く）ので、未知の順序対 $(x, y)$ を求め、方程式
\[
(x, y) × (a, b) = (c, d)
\]
を満たすとき、一意的な答えが存在し、これを次のように表すことができます：
\[
(x, y) = (c, d) ÷ (a, b), \quad \text{または} \quad
(x, y) = \frac{(c, d)}{(a, b)}
\]

(ε) さらに、加法と乗法を含む法則、すなわち分配律が満たされるべきです：
\begin{multline*}
(x,y) × \{(a, b) + (c, d)\} \\
= \{(x, y) × (a, b)\} + \{(x, y) × (c, d)\}.
\end{multline*}

これらの条件（α）、（β）、（γ）、（δ）、（ε）は、最初に複雑に見えるかもしれませんが、単純な幾何学的解釈を持つ解釈によって満たすことができます。

定義により、
\[
(x, y) × (x', y') = \{(xx' - yy'), (xy' + x'y)\}
\]

Thisは、二つの順序対間に記述された符号~$×$の意味を定義しています。この定義から明らかに、乗算の結果は別の順序対となることがわかります。また、方程式\Eq{(A)}の右辺の値は、同時に$x$と$x'$、および$y$と$y'$を入れ替えることで変化しないこともわかります。したがって条件\Eq{(\alpha)}と\Eq{(\beta)}は明らかに満たされます。方程式\Eq{(\gamma)}、\Eq{(\delta)}、\Eq{(\epsilon)}の満足性を証明するのも、与えられた幾何学的解釈がある場合、同等に簡単です。その幾何学的解釈についてすぐに説明しますが、それを行う前に、我々がこの詳細な展開で達成した目的を見つめ直すのは興味深いでしょう。

形式$x^{2} = -3$の方程式に遭遇し、これは正または負の実数で表すことができない解を持ちませんでした。その後、方程式$x^{2} = -1$を解釈できるなら、つまり$\sqrt{(-1)} × \sqrt{(-1)} = -1$となるような$\sqrt{(-1)}$を定義できれば、すべての困難が消えます。

今、特別な順序対\index{Zero}$(0, 0)$、$(1, 0)$、および$(0, 1)$を考えましょう。

既に証明済みですが
\[
(x, y) + (0, 0) = (x, y).
\]

さらに現在は
\[
(x, y) × (0, 0) = (0, 0).
\]

したがって、加算と乗算の両方において、順序対$(0, 0)$は初等的な算術と代数におけるゼロの役割を果たします。上記の式を$x + 0 = x$、および$x × 0 = 0$と比較してください。

また、$(1, 0)$も考慮します：これは初等的な算術と代数における$1$の役割を果たします。これらの初等的科学では、$1$の特別な特性はすべての値の$x$に対して$x × 1 = x$であることです。現在の乗算法則によれば
\[
(x, y) × (1, 0) = \{(x - 0), (y + 0)\} = (x, y).
\]

したがって、$(1, 0)$は単位順序対です。

Finally consider $(0, 1)$: このものは $\sqrt{(-1)}$ の意味を表します。したがって、$\sqrt{(-1)} × \sqrt{(-1)} = -1$ という特性を持つべきです。現在、順序対の乗法法則により
\[
(0, 1) × (0, 1) = \{(0 - 1), (0 + 0)\} = (-1, 0).
\]

$(1, 0)$ は単位対であり、$(-1, 0)$ は負の単位対です。したがって、$(0, 1)$ には望まれる特性があります。しかし、$-1$ の平方根は $±\sqrt{(-1)}$ と2つあります。次に $(0, -1)$ を考慮します；ここで再度、$(-1)^2 = 1$ を念頭に置いて
\[
(0, -1) × (0, -1) = (-1, 0).
\]

したがって、$(0, -1)$ は $\sqrt{(-1)}$ のもう一つの平方根です。従って、順序対 $(0, 1)$ と $(0, -1)$ は $±\sqrt{(-1)}$ の順序対による解釈となります。しかし、どちらがどの符号に対応するのでしょうか？$(0, 1)$ が $+\sqrt{(-1)}$ に対応し、$(0, -1)$ が $-\sqrt{(-1)}$ に対応するのか、それとも $(0, 1)$ が $-\sqrt{(-1)}$ に対応し、$(0, -1)$ が $+\sqrt{(-1)}$ に対応するのでしょうか？答えは、採用するシンボリズムを選択することに完全な自由があるということです。

順序対は三種類に分類できます：(i) 両方の成分がゼロでない「複素虚数」型 $(x, y)$；(ii) 第二成分がゼロの「実数」型 $(x, 0)$；(iii) 第一成分がゼロの「純虚数」型 $(0, y)$。これらの種類間の関係を考慮します。まず、「複素虚数」対 $(x, y)$ と「実数」対 $(a, 0)$ の積を考えます：
\[
(a, 0) × (x, y) = (ax, ay).
\]

これは単に順序対 $(x, y)$ の各成分を正または負の実数 $a$ で乗算するだけです。

次に、「複素虚数」対 $(x, y)$ と「純虚数」対 $(0, b)$ の積を考えます：
\[
(0, b) × (x, y) = (-by, bx).
\]

ここでの効果はより複雑で、幾何学的解釈を考慮する前にさらに3つの特別なケースを挙げてから理解するのが最も適切です。

第三に、「実」の偶対$(a, 0)$を「虚」の偶対$(0, b)$で掛けると、
\[
(a, 0) × (0, b) = (0, ab).
\]

第四に、「実」の偶対$(a, 0)$と$(a', 0)$を掛けると、
\[
(a, 0) × (a', 0) = (aa', 0).
\]

第五に、「虚」の偶対$(0, b)$と$(0, b')$を掛けると、
\[
(0, b) × (0, b') = (-bb', 0).
\]

次に幾何学的な解釈に移ります。まずいくつかの特別なケースから始めます。
図11（Fig. 11）において偶対$(1, 3)$と$(2, 0)$を考え、方程式
\[
(2, 0) × (1, 3) = (2, 6)
\]
を考察します。

図11の図において矢印$OP$は偶対$(1, 3)$を表し、矢印$ON$は偶対$(2, 0)$を表し、矢印$OQ$は偶対$(2, 6)$を表します。したがって、幾何学的に偶対$(2, 0) × (1, 3)$の積を求めることは、矢印$OQ$の長さが矢印$OP$と$ON$の長さの積であることを確認し、さらに矢印$OP$を矢印$Q$まで延ばすことによって求めることができます。また、偶対$(0, 2) × (1, 3)$を考えると、
\[
(0, 2) × (1, 3) = (-6, 2)
\]
となります。

矢印$ON_{1}$は偶対$(0, 2)$を表し、矢印$OR$は偶対$(-6, 2)$を表します。したがって、新しい積を表す矢印$OR$は矢印$OQ$に対して直角であり、同じ長さを持ちます。注意すべきは、矢印$OQ$の長さに関する法則が前の場合と同様である点で、その長さは二つの矢印が掛け算されるときの長さの積であるという点です；しかし今度は矢印$ON_{1}$が「縦軸」$OY$に沿ってあり、かつ矢印$ON$が「横軸」$OX$に沿っている代わりに、矢印$OP$の方向が直角に回転している点です。

これまでの例では、ベクトル$\vec{OP}$がベクトル$ON$と$ON_{1}$によって変更されることを観察してきました。考え方が逆転し、ベクトル$ON$と$ON_{1}$がベクトル$\vec{OP}$によって変更されることを考えることで、方向の一般法則を得ることができます。長さに関する法則は変わらず、結果ベクトルの長さは二つのベクトルの積の長さとなります。

拡大された$ON$（すなわち$OQ$）の新しい方向は、$OX$から$OY$に向かって反時計回りに角度$XOP$だけ回転することで得られます：この特定の場合においては、この回転が$OQ$を$\vec{OP}$と同じ直線上に配置することになります。

また、ベクトル$ON_{1}$と$\vec{OP}$の積を考えます。拡大された$ON_{1}$（すなわち$OR$）の新しい方向は、反時計回りに角度$XOP$だけ回転することで得られます：すなわち、角度$N_{1}OR$は角度$XOP$と等しいです。

現在、ベクトルの乗算の幾何的な表現に関する一般法則を述べることができます：
図3（縦12cm）

二つのベクトル$\vec{OP}$と$\vec{OQ}$の積はベクトル$\vec{OR}$で、その長さは$\vec{OP}$と$\vec{OQ}$の長さの積であり、方向$\vec{OR}$は角度$XOR$が角度$XOP$と角度$XOQ$の和に等しいように定まります。

したがって、ベクトル$\vec{OP}$を角度$XOP$（すなわち$\angle QOR = \angle XOP$）だけ回転させることでベクトル$\vec{OQ}$を得るか、またはベクトル$\vec{OQ}$を角度$XOQ$（すなわち$\angle POR = \angle XOQ$）だけ回転させることでベクトル$\vec{OP}$を得ることができます。

私たちはこの一般の法則を証明しません。それによって、本書の目的から逸脱するより複雑な数学的手法に導かれる恐れがあります。しかし、乗算に対する結合律 [式(\gamma)] が満足されていることはすぐにわかります。

まず、結果ベクトルの長さは通常の実数の乗算の過程で得られます；したがって、この過程に対して結合律は成り立ちます。

また、結果ベクトルの方向は角度の単純な加算から得られ、この過程に対しても結合律は成り立ちます。

以上が乗算についてです。現在、加法と乗法を考慮することで、平面内のベクトルに対する代数や「計算」がどのように構築されるかを素早く示しました。これは任意の2つの平面ベクトルが加算または減算され、また乗算または除算も可能です。

これらの過程の技術的な詳細については言及しませんでした。それは数学的な詳細に深く入り込む可能性があるためです；しかし、私たちは一般的な手続き方法を示しました。このような代数的シンボルを解釈するとき、「虚数」や「複素数」と呼ばれます。これらの用語は細部であり、私たちが考えるべきことが多すぎるために、これらが非常にうまく選ばれたものであるかどうかを調べる余裕はありません。

調査の結果、方程式 $x + 3 = 2$ や $(x + 3)^{2} = -2$ のような式は常にベクトルの形で解釈され、それらに対する解が見つかることが可能です。このような解釈を求める際には、$3$ は $(3, 0)$、$-2$ は $(-2, 0)$、そして $x$ は「未知数」のペア $(u, v)$ に相当することに注意することが良いでしょう。したがって、これらの方程式はそれぞれ $(u, v) + (3, 0) = (2, 0)$、および $\{(u, v) + (3, 0)\}^{2} = (-2, 0)$ となります。

我達现在完全解决了最初一考虑代数元素便立即注意到的初始困难。从解题中发展出来的科学在思想上比我達开始时所涉及的内容复杂得多。事实上，我達创造了一门全新的且完全不同类型的科学，它不仅能够満たす旧科学发明的所有目的，还能解决更多的问题。
しかし、自分たちの努力に対する自慢話をする前に、学生の心に今度は疑念が生じるべきである。読者自身が問うべき質問は：このような新しい解釈の創造はどこまで続くのか？確かに、我々は二次方程式 $x^{2} - 2x + 4 = 0$ を常に解くことができるような代数を解釈する方法を見つけた。しかし、他の無限に多くの方程式が存在し、例えば $x^{3} - 2x + 4 = 0$, $x^{4} + x^{3} + 2 = 0$ などがある。新しい方程式が現れたときには新たな科学を創造する必要があるのか？

もしこれが真なら、我々の以前の調査は、一部の人々にとっては興味深いものであっても、実際には非常に微細な重要性しかないだろう。しかし、現代解析が可能となった偉大な事実は、ベクトルの計算法によって、すべての公式が適切に解釈され、方程式中の「未知数」は必ずしもベクトルを示すことが示されたことである。この科学は、基本的な概念に関して自立して完成したと言える。その形態化は蒸気機関が精巧化される頃とほぼ同じ時期であり、思想が物事の勝利を成し遂げる強大な武器として残り続けるだろう。その機械の例が博物館で展示され、やや早い時代の頭盔や胸当てと共に並ぶときでもある。

\chapter{第9章：Chapter IX: Coordinate Geometry}

坐標幾何の手法や概念は、既に前章で使用されてきました。現在はそれらを自身のためにより精密に検討する時です；これにより他の到達した概念に対する理解も深まるでしょう。現章および次章では、正負実数という概念に戻り、最後二章で導入された虚数については無視します。

我々は常に、平面において二つの軸$XOX'$と$YOY'$を用いて、どのような点$P$も正のまたは負の数$x$および$y$という一対の数で位置を定めることができるという考えを使用してきました。ここで（\cf \Fig[fig.]{13}）$x$は長さ$OM$であり、$y$は長さ$PM$です。この概念は見かけ上単純ですが、座標幾何学の主要なアイデアです。その発見は数学的思想の歴史上重要な時代を示すものです。これは（既に述べた通り）哲学者デカルトによって生み出されました\index{Descartes}{}、彼はベッドで横になっている朝にこの重要な数学者の方法を見つけることに到達しました。

哲学者は数学的な知識が深ければ深いほど、その科学に最良のアイデアをもたらすことが多いです。一方で、それらの哲学者のうちほとんど例外なく、僅かなまたは急いで得られた数学に関する言及は全く価値がないものであり、 trivia か誤りであることが多くあります。事実、数学における最終的な概念は非常に単純で、子供でも理解できるほど簡単であり、哲学的な思考の範囲内に留まっています。おそらくその極めて単純さが誤りを生む原因かもしれません。我々はこのような単純な抽象的概念について考える習慣がなく、既成の思考の道からそれた瞬間には長年の訓練が必要でさえあります。

座標幾何の発見と、ほぼ同時期にProjective几何学的も発見したことは、知識の歴史において繰り返し確認されている事実を示しています。つまり、最も広く知られているトピックの中から、いくつかの最伟大的な発見がなされることがしばしばあります。

17世紀になると、幾何学はすでに紀元前2000年以上にわたって研究されていました。ギリシャ人による幾何学の隆盛期を始めるのが適切であれば、その起源はユーリドクスにまでさかのぼります。アレクサンドリア大学で教鞭を執ったユーリドクスは紀元前330年頃生まれで、彼は長い系列の前人作業を系統化し、発展させました。その一部は天才的な人物でした。

その後、世代を超えて数学者たちはこの分野の改善に努め続けました。しかし、その進歩を妨げる致命的な障害は、その研究が狭い集団の類似した出身と視点を持つ人々によって制限されていたことではありませんでした。むしろ、17世紀にはエジプト人やギリシャ人、アラブ人やドイツ人の頭脳を通じて幾何学が広く普及していました。

それでも、そのような多様な頭脳によって長い年月にわたって研究が行われた後でも、最も重要な秘密はまだ発見されていませんでした。

誰も初等幾何学の要素を学んだことなく、導き方となる方法の欠如を感じたことはないだろう。各命題は新たな独創性の展示によって証明される必要がある。そして、この特徴を持つ科学は、科学的な思考にとって重要な要素である方法性を欠いている。

座標幾何学の特別な点は、それが初めて方法を導入したという点にある。数学的科学の遠く離れた推論は、理論上重要なものではない。科学が完全なものになるまでには、その本質は、任意の範囲内のどのトピックについても容易に情報を得るための大規模で関連する方法の展示であるべきだ。

科学の成長は主に量ではなく、アイデアによるものであり、アイデアが増えるほど、書き出す価値のある推論は少なくなる。不幸にも、数学は教科書によって数多くの補助的な命題が繰り返し記述され、それらの重要性はより一般的な真理の特殊な例として吸収されて失われている。そして、既に強調したように、一般性は数学の魂である。

また、座標幾何学はもう一つの特徴を示す。それはすでに指摘された通り、数学的科学が発展するにつれて、互いに嵌め合わ合い、共通のアイデアを持つことだ。言い換えれば、数学の各分野は永続的な一般化プロセスを経験し、その過程で融合していく。この理由は、科学の本質からくるものであり、それが普遍性を持っていること、つまりすべての存在するものの性質に関連する一般的な真理を扱うことに由来する。

座標幾何学の関心点は、空間の科学として始まった幾何学と、数の科学として起源を持つ代数学を結びつける点にある。

現在、二つの科学の主なアイデアを思い返し、デカルトの座標法によってどのように関連付けるかを見ていきましょう。まず代数について考えてみましょう。虚数については触れず、正または負の符号を持つ実数のみを考えます。基本的なアイデアは、任意の数である変数数というものです。これは文字で表され、具体的な数値では表現されません。次に、変数間の関係性を考慮します。例えば、$x$ と $y$ が二つの変数の場合、これらは方程式 $x + y = 1$ または $x - y = 1$ によって関連付けられることも考えられます。あるいは無数の他の方法でも関連付けられることがあります。

これにより、代数形の概念を適用することができます。実際には、何か興味深い型の関係性を考えます。つまり、変数数という最初の概念から、数値の変数関係性という次に進んだ概念へと昇華します。したがって、$x + y = 1$ の関係を一般化して $ax + by = c$ の関係にします。ここで $a$、$b$ そして $c$ は文字であり、任意の数を表し、実際にはこれらも変数です。しかし、これらの変数が変数関係性を決定するものです。関係性が確定されると、それは変数数 $x$ と $y$ の間に関連付けられます。

$a$、$b$、そして $c$ などのような、関係性を決定するために使用される変数は「定数」またはパラメータと呼ばれます。この文脈での「定数」という用語の使用が最初に見えて不自然に思えるかもしれませんが、実際には非常に自然です。なぜなら、数学的な調査は $a$、$b$、$c$ が既に決定された後の $x$ と $y$ の関係性についての関係を調べることに関与しているからです。したがって、$x$ および $y$ 相対的に、「定数」である $a$、$b$、そして $c$ は実質的に定数となります。したがって、$ax + by = c$ は特定の代数形の一般的な例を表し、あるクラスに属する変数関係性を表します。

以下を忠実に翻訳しました：

また、$x^{2} + y^{2} = 1$ を一般化して $ax^{2} + by^{2} = c$ とします。
さらに、$ax^{2} + 2hxy + by^{2} = c$ とします。
そして、さらに $ax^{2} + 2hxy + by^{2} + 2gx + 2fy = c$ とします。

またここでも、様々な代数的形式によって示される変動的な関係が導き出されます。

次に幾何学へと移りましょう。その科学の名前は、すぐに私たちの心を図形や図表、三角形や長方形、正方形や円などの考えに連想させます。これらの図形の単純な性質の研究が初等幾何学の主題であり、適切に初心者に対して教えられるべき内容です。しかし一瞬考えてみれば、これは真の概念とは言えません。子供が scissors で切り出した三角形や正方形のような形状から始める幾何学的な論理は正しいかもしれません。しかし、三角形とは何か？それは三本の直線の端点で囲まれた図形です。

空間の境界を線分によって定義することは非常に複雑な概念であり、真の単純な一般的概念が主題の骨格となることを示すものではないのです。我々はもっと単純で一般的なものが必要です。この誤った初期概念への執着---これは初歩的な考えを創出するためには自然で良い概念である---が、数世紀にわたる科学の研究における比較的不活性化の原因でした。座標幾何学とその発明者 Descartes は、幾何学的思想のために真の単純な対象を開示した価値があるのです。

直線の一部ではなく、両方向に無限に延びる全体の直線について考えましょう。これは私たちが幾何学的調査から始めるべき一般的な概念です。ギリシャ人は現在すべての現代的な幾何学的思想において基本となるこの概念を見つけることなく、その有用性を発見しなかったようです。ユークリッドは一直線を常に二つの確定した点間で引かれたものとして考え、それをこのセグメントを超えて延長する場合に注意を喚起します。彼は直線を一度与えられた全体的な存在とは考えません。この明確な定義と制限は、感覚によって直接的には見えない無限を排除することにより、ギリシャ人の多くの活動において非常に特徴的でした。これはギリシャ建築とゴシック建築、またギリシャ宗教と現代の宗教との違いに表れています。ゴシック cathedral の尖塔と現代幾何学における無限大の直線の重要性は、現代世界の変化を象徴しています。

直線を全体として考えると、現代の幾何学はその基本的な概念から始まる。しかし、それ以外の種類の線も私たちに思い浮かび、それが点ごとに一定の特性を示す完全な曲線という概念へと到達する。直線はすべての点で一直性を持つことを示すように、曲線もそのような特性を特定の点で示す。例えば、円はすべての点で中心から与えられた距離にある特性を示し、また楕円は二つの固定点（焦点）からの任意一点に対する二つの距離の和が定数である特徴を持つ卵形の曲線である。つまり、二つの焦点が同じ点に重なったとき、即ち焦点が一致したときに円は楕円の特殊な場合となることが明白である。そのとき、二つの距離の和はただの直径の2倍になる。

古代の人々は楕円と円の性質を知っていたし、当然のことながらこれらを全体として考えていた。例えば、ユークリッドは円の単なる部分（すなわち、セグメント）から始まることなく、常に円が描かれた全体を考えていた。幾何学において真の基本的な線が円ではなく直線であることが不幸であり、そのためユークリッドが直線を不完全に扱ったことの影響は小さくならなかった。

この曲線の一般的な概念は、その上でどのような性質を示すかにより異なる点においてある一定の性質を有することであり、幾何学では「 locus（ locus）」という用語で表現される。locusとは、ある性質を持つすべての点からなる曲線（または平面以外の場合は表面）である。

各点間の関係にある任意の点が持つ性質に対して、対応する locus が存在し、その locus は所有する性質を持つ全ての点で構成される。locus全体の性質を調査する際には、locus上の任意の点や点群を考える。したがって、幾何学では再び変数の基本的な概念に遭遇する。

さらに、直線、円、楕円などの分類を行う際にも、同じ形式の条件を持つ locus の分類という考え方に立ち返る。

従って、代数学では変数の数、変数間の関係、そして代数的形式に基づくタイプへの分類に関心を持ち、同様に幾何学では、ある条件を満たす点群からなる locus とその locus のタイプによる分類に関心を持つ。

現在、座標幾何の本質は代数的な関係と幾何学的場所との識別にあります。平面上の点はアルгебра的には$x$と$y$という2つの座標で表され、 locus上の任意の点が満たす条件はそれに応じて$x$と$y$間の関係を表します。最後に、一般の代数的形式で表現可能な関係には、その形式が同じ幾何学的な条件を持つ一般的な型のlocusが対応します。これにより、両科学間でのアイデアや結果の完全な交換が可能になりました。一方の科学は他方を照らし、自身も著しく強化されます。歴史的な冒険と発見の瞬間に人々が感じた感情を思い浮かべるのも不可能です——コロンブス\index{コロンブス}が初めて西海岸を見たとき、ピサロ\index{ピサロ}が太平洋を凝視したとき、フレンチル\index{フレンチル}が自分のキテの糸から電気火花が来たとき、ガリレオ\index{ガリレオ}が初めて望遠鏡で天体を見たとき。これらの瞬間は学習者にとっても与えられ、抽象的な思考の領域において最も高い位置に置かれるべき朝は、デカルト\index{デカルト}がベッドで座標幾何の方法を発明したときです。

一次元の代数形式 $ax + by = c$ に対応する locus は直線であり、逆に任意の直線にはこの形式の式が対応します。幾何学的な simplest の locus が代数的な simplest の形式と一致することは幸運です。実際、幾何学と代数の simplicity の一般的な対応関係こそがこの分野全体の力を与えています。これは幾何学と代数との間のつながりが偶然的ではなく、深く根ざし、本質的なものであることを示しています。

locus に対応する式は locus ``の''（または ``に''）式と呼ばれます。直線の式のいくつかの例を挙げることでこの主題を説明します。

$y - x = 0$ を考慮すると、一般形の $a$,~$b$, および~$c$ はそれぞれ $-1$,~$1$, そして $0$ に置き換えられています。この直線は図中的 ``原点''~$O$ を通っており、角度~$XOY$ の二等分線になっています。これは図中の直線~$L'OL$ です。この直線が原点~$O$ を通ることを確認するには、$x = 0$ および $y = 0$ を同時に代入することで方程式が満たされることを見ることができます。実際、$O$ の座標は $0$ と $0$ であるためです。同様の方法で一般化すると、原点を通る任意の直線の式は $ax + by = 0$ の形式であることがわかります。

方程式 $y + x = 0$ の locus も原点を通過し、角度~$X'OY$ の二等分線になっています。これは図中の直線~$L_{1}OL_{1}'$ です。

考慮する $y - x = 1$：この方程式に対応する軌跡は原点を通らない。したがって、その軌跡が軸に交わる点を求めよう。$x$ 軸との交点は座標 $(x, 0)$ の点である。$y = 0$ を方程式に代入すると $x = -1$ となるので、この点の座標は $(A) (-1, 0)$ である。同様に、$y + x = 1$ は図中的直線 $A_{1}B$ の方程式であり、その軌跡は $L_{1}OL_{1}'$ と平行である。一般の定理として証明できるのは、形が $ax + by = 0$ と $ax + by = c$ の直線は平行であるということである。

次に遭遇する軌跡群は十分重要で、個別の章で扱う価値がある。しかし、それらの前に、主題の主なアイデアについてもう少し詳しく説明しよう。

任意の点$P$の位置は、原点$O$を適当に選び、直角に交わる軸$OX$と$OY$を定め、その後その点の座標$x$と$y$、すなわち$OM$と$PM$を記録することによって決定される（図13参照）。また、前章で見たように、点$P$はベクトル$OP$によっても決定できる。ベクトルの概念には定まった方向だけでなく定まった長さも含まれている。抽象的な数学的な観点から見れば、任意の原点の概念は人工的で不自然に思えるかもしれませんが、同様に自由に描かれた軸$OX$と$OY$も同じように感じられるかもしれません。しかし、宇宙の現象に対する数学の適用に関しては、ここでは直接簡潔な形で、私たちが感覚的に認識する世界をどのように見ているかという最も基本的な事実を表しています。私たちはそれぞれ、「ここ」と呼ぶ原点に対して自分の経験を参照します：私たちは空間内の特定の部分に位置し、その周りに宇宙全体をグループ化することによって自己の存在を表現します。「ここ」に対する私たちの位置は、私たちが現象を観察する基本的な事実です。すべての現象を平等な視点で観察できるような生物を考えることができます。しかし私たちはそうではありません。私たちの足元にある猫よりも、ケープ hornでの地震やミキシル道上での世界の破滅の方がより多くの注意を集めます。私たちが他の人々と共有する知識を作り出す際に、自分自身の個別の「ここ」に対する厳格な自己中心性を放棄しなければならないこともあります。私たちは「ほぼここ」という言葉で「ここ」を置き換えることになります；したがって、最も近い町の市庁舎からマイルを測定するか、または国の首都から距離を測定します。地球を測定する際には、科学者は地球の中心を原点とし、天文学者は甚至に極度の利他的性から太陽内部を原点とすることもあります。最後の原点がいかに遠くても、あるいは近い固定星の中間の便利な点まで進んでも、空間の無限大に比べれば、私たちが宇宙を探査する最初の手順は「ほぼここ」を原点として決定することであるという事実は変わりません。

再度、座標$OM$と$MP$（すなわち$x$と$y$）がベクトル$OP$に対して持つ関係は有名な平行四辺形の法則の一例であり、これは図\[fig.]{8}\に示すように、平行四辺形$OMPN$を完成させることで容易に確認できる。物理科学における「ベクトル」つまり方向付いた量という概念は根本的なアイデアである。どのような運動する物体も特定の方向に向けて一定の速度を持っているため、その速度は方向付いた量、すなわちベクトルである。また力も特定の大きさを持ち、特定の方向を持つ。したがって、解析幾何において「原点」、「座標」、「ベクトル」という概念を導入するとき、物理世界の基本的事実に対応する抽象的概念を研究していることを意味する。

\chapter{第10章：Chapter X: Conic Sections}

当ギリシャの幾何学者が、直線と円で構成される図形のより明白で興味深い性質を尽きるまで探求した後、彼らは他の曲線に焦点を移しました。そして、彼らの几乎不败的本能により、重要なことを考え出すことに成功した彼らは主に放物線、双曲線、楕円線などの放物面で切り取られる円錐表面の曲線に注力しました。この学科の発明者として信用されるべき人物はメネラウスです（紀元前375年生紀元前325年没）；彼はプラトンの弟子であり、アレクサンドロス大王の教師の一人でした。また、彼の他の教師の一人は哲学者アリストテレスでした。

アレクサンドロスは、良い教育の明確な例であると言えるでしょう。彼の他の教師の一人はアリストテレスでした。メネラウスがアレクサンドロスに少し退屈だと感じていたのか、彼は証明を短くするようにと頼みました。メネラウスはその要求に対して次のように返答しました：「国には秘密の道や王様の道があるが、幾何学にはすべての人々のために一つだけの道しかない」。この回答はアレクサンドロスがすぐに理解した意味においては十分に真実でした。しかし、メネラウスが証明を短くすることは不可能だと考えていたら、彼は深刻な誤りを犯していたのです。そして、多くの現代の数学者は、ギリシャ人が放物線の性質を証明する方法を学ぶことを強いられたらひどく退屈することでしょう。

科学における概念の導入によって得られる力の増加を示す例としては、証明が短くなる過程と同時に概念の豊かさが成長する過程を観察することでよく理解できます。ある種の数学者は、主題のアイデアに耽るよりも早く重要な問題の証明に進むことを常に少し不満に感じることがあります。しかし、科学の歴史は彼に対して完全に対立しています。科学には王様の道がありますが、最初にそれを歩いたのは天才であり、国王ではありません。

円錐曲線が数学的に最初に現れた方法は次の通りでした：円錐（図\[fig.]{15}\参照）を考え、その頂点（または点）を$V$とし、底面が円形の$STU$上に立っています。例えば、電球に対して設置されたコンパクトな影はこのような表面の例です。今、$V$を通る「生成」する直線が全て延長されると、二つの円錐が得られます。そして$PQR$は$V$と反対側にある別の円形の交差点となります。円錐$CVC'$の軸はこれらの円の中心を通り、それぞれの平面に垂直であり、これらは互いに平行です。図において紙の平面の奥に存在する曲線の部分は点線で示され、平面またはその前方にある部分は継続線で示されています。今、この二つの円錐が軸$CVC'$に対して垂直ではない平面上で切られるとします。あるいは、必ずしも垂直である必要はありません。その結果、以下の3つの状況が生じます：---

（1）平面が一つの半円錐上に完全に存在する閉じた楕円形の曲線、例えば$ABA'B'$のように切れる場合があります。この場合、平面はもう片方の半円錐と全く交わらないことがあります。このような曲線は楕円と呼ばれ、それは楕円形の曲線です。円錐のこのような断面の特別な例として、平面が軸$CVC'$に対して垂直である場合、断面、例えば$STU$や$PQR$は円となります。したがって、円は楕円の特別な例と言えます。

(2) 平面は、円錐の生成線の一つに沿って接平面を引くことができる。例えば図において、曲線 $D_{1}A_{1}D_{1}'$ の平面は生成線 $VS$ に沿って接する接平面と平行であり、この曲線は依然として半円錐の一つに制限されるが、閉じた卵形の曲線ではなくなり、生成線が頂点から延長されると無限に続く。このような円錐截痕をパラボラという。

(3) 平面が二つの半円錐を両方とも切る場合、完全な曲線は二つの離れた部分、または「枝」と呼ばれる二つの部分からなります。図において、二つの枝 $G_{2}A_{2}G_{2}'$ と $L_{2}A_{2}'L_{2}'$ が組み合わさって曲線を形成しています。どちらの枝も閉じていないことがわかり、半円錐が頂点から延長されるにつれて無限に広がります。このような円錐截痕を双曲線という。

従って、三種類の円錐截痕があります：楕円、パラボラ、および双曲線。パラボラは、一定の条件を満たすより特殊な種であり、楕円と双曲線の間の極限的な場合であることが容易に理解できます。アポロニウス・ペルガエウス（紀元前260年頃生没）によって書かれた系統的な円錐截痕に関する論文が存在し、16世紀まで標準的な著作として受け入れられました。

これらの曲線の性質を調査することは、ギリシャの幾何学者にとっていかに不快で困難な作業であったかがすぐに明らかとなる。これらの曲線は平面曲線であるが、その調査には立体図形を透视的に描くことが含まれている。したがって、上記の図においても補助線をほとんど引いていないにもかかわらず、図は十分に複雑になっている。これらの曲線は平面曲線であり、これらを定義するためには平面内だけで済ませることができそうに思える。同時に、「立体」での定義では一貫した方法論が一つ存在し、それは円錐を平面で切るという方法で、楕円、双曲線、そして放物線の三種類を生み出す。同様に、「平面」での定義でも一貫した方法論が三種類に分かれることであるべきだ。これらの曲線は図形 \FigNum{16}, \FigNum{17}, および \FigNum{18} の曲線の形状と一致している。図形において点 $A$ と $A'$ は頂点と呼ばれ、線分 $AA'$ は主軸と呼ばれる。

パラボラ（\cf \Fig[fig.]{17}）には一つだけ頂点があることが注意されるべきである。アポロニウスは、楕円と双曲線（図形 \FigNum{16} および \FigNum{18}）において $PM^{2}$ に対する $AM·MA'$ の比が一定であり、放物線（\cf \Fig[fig.]{17}）においては $PM^{2}$ に対する $AM$ の比が一定であることを証明した。彼はこの事実に基づいて多くの仕事を進めている。明らかに我々は平面内だけで定義する一貫性のある方法論へと進んでいるが、まだ完全には到達していない。

図 \FigNum{16} および \FigNum{18} では、二点 $S$ と $S'$ を見ることができます。また、図 \Fig[diagram]{17} では一点 $S$ を見ることができます。これらは曲線の焦点であり、最も重要な点です。アポロニウスは、楕円の場合、点 $P$ が曲線上を動くとき、$SP + S'P$ の和が一定で、かつその値は $AA'$ と等しいことを知っていました。同様に、双曲線の場合、$S'P - SP$ の差が一定で、$P$ が一方の枝上にあるときにその値は $AA'$ であり、$P'$ が他方の枝上にあるときに $SP' - S'P'$ の差も一定で、その値も $AA'$ であることが知られています。しかし、パラボラには対応する点が見当たりませんでした。

ついに 500 年後、最後の大ギリシャ幾何学者パッポス・アレクサンドリヌスは、この思考の最終的な秘密を発見しました。図 \FigNum{16} および \FigNum{18} では、二つの直線 $XN$ と $X'N'$ を、また図 \Fig[diagram]{17} では単一の直線 $XN$ を見ることができます。これらはそれぞれの曲線の直交点であり、楕円と双曲線には各々二つあり、パラボラには一つあります。各直交点はその近接する焦点に対応しています。

特異性の性質として、どのような曲線かに関わらず、焦点 $S$ およびその対応する直交点 $XN$ の特性は、任意の点 $P$ が曲線上にあるとき、$SP : PN \left( \ie\ \dfrac{SP}{PN} \right)$ の比が一定であることです。ここで、最終的に平面から離れないで求めた望む性質を見つけました。この性質は全ての三つの曲線に対して一貫して述べられます。楕円の場合、その比は 1 より小さい、パラボラの場合、その比は 1 に等しく、双曲線の場合、その比は 1 より大きいです。

パッパスが調査を終えたとき、彼はおそらく、小規模な拡張を除いてはその主題は実質的に尽きていると感じたことでしょう；そして、科学の歴史が千年以上先まで予見できていたら、彼の信念は確認されたことでしょう。しかし実際には、この数学の分野に関連する真に生産的なアイデアはまだ触れられておらず、誰もその極めて重要な自然への応用を予測していませんでした。知識と研究を明らかや有用なものだけに限定しようとする人々に対するより印象的な警告を与えるものはありません。コンキックスセクションが十八世紀間も抽象的な科学としてのみ研究され、他のどの有用性よりも数学家の好奇心を満たすことに他ならないという考えがなかったのに、長い期間の抽象的研究の末にそれが自然の最も重要な法則の理解のために必要な鍵であることが発見されたからです。

同時に、全く異なる天文学の研究は進行していました。偉大なギリシャ天文学者ポタメデウス（没年168年）は、アレクサンドリア大学でその主題に関する標準的な論文を発表し、太陽と惑星の恒星間の観測可能な運動を地球が静止しているという概念によって説明しました。その後の約一千三百年の間に、天文観測の数と精度は増加し、結果としてポタメデウスの仮説に基づく惑星の運動の説明はますます複雑になりました。コペルニクス（1473年生没1543年）は、太陽が静止していると仮定し、地球と惑星がそれを回るという考えでこれらの天体の運動をより単純な方法で説明できると指摘しました。しかし、彼はこれらの運動を基本的な円運動に修正される一部の任意の補正によって修正された本質的に円形であると考えていました。ケプラーが1571年ドイツのシュタウフェンブルクで生まれたときには、この状況でした。二つの科学がありました：双曲線などの接線の幾何学と天文学。これらは遠い古代から研究されていましたが、互いに関連しているという疑念はありませんでした。ケプラーは天文学者でしたが、優れた幾何学者でもあり、双曲線に関するアイデアを時代遅れよりも進んでいたと考えることができます。科学的研究における成功は一つの専門分野に独占的に没頭する必要性という考えは誤りであり、新しいアイデアは知識の不思議な組み合わせから生まれやすいです—必ずしも広範な知識ではなく、異なる思考線路の方法とアイデアに対する深い理解からです。チャールズ・ダーウィンが進化法則の概念に至った際には、人口に関する有名な『エッセイ』（マラース著）を読むことで助けを得たことを思い出すでしょう—少なくとも当時は異なる主題についての著作でした。

ケプラーは1609年に第1条と第2条、そして10年後の1619年に第3条を提唱した太陽系の運動に関する三つの法則を述べました。

(1) 行星の軌道は楕円であり、太陽はその焦点にあります。

(2) 行星が軌道上を動くとき、太陽から惑星までの半径ベクトルが等しい時間で等しい面積を描きます。

(3) 各々の惑星の周期の二乗比は、その主軸の三乗比に比例します。

これらの法則は、より根本的な思想の発展の一歩に過ぎませんでした。ニュートン（1642年生まれ、1727年に没）は普遍的重力の概念を思いつき、任意の二つの物体間には質量の積に比例し、その他の物体からの距離の二乗分の逆数に比例する引力があるという考えを持ちました。この広範な一般法則と、彼が最終的に完成させた運動法則を組み合わせることで、すべての天文現象、つまりケプラーの法則も含むもの説明できることが証明されました。これにより現代物理学の基礎となりました。ニュートンはさらに証明しました：彗星は非常に長い楕円軌道や双曲線軌道、または近似して双曲線である放物線を描く可能性があるということです。再帰する彗星（ハレーの彗星など）は当然、楕円軌道を描きます。しかし、重力法則の証明における決定的なステップは、惑星の運動と放物線理論との関係性であるケプラーの法則の検証でした。これが重力法則の最初の概念を提示するヒントにもなりました。

从17世纪开始，曲线的抽象理论与座標几何和射影几何的双重复兴一同发展。射影几何中，基本概念围绕を通じて一个共同点（“束”的顶点）的一组线（或称为“束”）进行考虑。

図2.5in (Fig.[fig.]{19}) 任意の4つの固定点 $A$, $B$, $C$, $D$ を円錐曲線上にとるとき、変動点 $P$ 上の線 $PA$, $PB$, $PC$, および $PD$ の束は特別な性質を持ち、それが交叉比の恒定性である。交叉比は射影幾何において基本的な概念であり、ここでその定義を詳しく説明する必要はない。

射影幾何では、これは実際には曲線の定義、または等価な性質である。我々が数世紀にわたって研究の中でどの程度古い原始的概念から逸脱したかは、明らかになるだろう。現在では、我々はギリシャ人が比較的小さな重要性を持つ一部の性質を握ったことを知っている。しかし、その曲線自体は与えられた注意に価値があった。この“截面”的概念の無関心性は、通常の数学用語において、その名前から「セクション」の言葉を省略することで表わされている。しばしば、これらは単に「コンキス」と呼ばれるようになり、「円錐の截面」という呼称が省かれた。

最後に、直交座標幾何学について前章で離れたところに戻ります。
\index{位勢}%
当時、一般の代数形 \(ax + by = c\) に対応する位勢の種類は何であるかを尋ねましたし、その結果は平面内の直線のクラスであることが分かりました。また、すべての直線がこの形式の式を持つことを見た一方で、この形式の任意の式も必ず直線に対応することも確認しました。
今度は次の一般的な代数形に進みます。これは \(x^2\) および \(xy\) と \(y^2\) の項を導入することで得られます。したがって、新しい一般形は以下のように書かなければなりません：
\[
ax^2 + 2hxy + by^2 + 2gx + 2fy + c = 0
\]
この式は何を表していますか？答えは、（これを位勢として表す場合）、常に放物線の一種であることが分かりますし、さらに、すべての放物線の式も必ずこの形に変換できることが言えます。この方程式によって与えられる放物線の特定の種類を区別する方法は非常に簡単です。これは \(ab - h^2\) の値に完全に依存します。ここで、\(a\), \(b\), および \(h\) は上記のように「定数」として書かれています。もし \(ab - h^2\) が正の数であれば、曲線は楕円です；もし \(ab - h^2 = 0\) であれば、曲線は放物線であり；そしてもし \(ab - h^2\) が負の数であれば、曲線は双曲線です。

例えば、$a = b = 1$, $h = g = f = 0$, $c = -4$と置く。すると方程式$x^{2} + y^{2} - 4 = 0$を得る。
この方程式は中心が原点で半径が2単位の円の式であることが簡単に証明できる。次に、$ab - h^{2}$は$1 \times 1 - 0^{2}$となり、つまり1であり、従って正である。したがってこの円は楕円の一例であり、その通りである。一般化すると、任意の円の方程式は$a(x^{2} + y^{2}) + 2gx + 2fy + c = 0$という形で表すことができる。よって$ab - h^{2}$は$a^{2} - 0$となり、つまり$a^{2}$であり、必ず正である。従ってすべての円は楕円の条件を満たす。
放物線の方程式の一様な形は
\[
(dx + ey)^{2} + 2gx + 2fy + c = 0,
\]
であり、二次項（平方項）は完全平方に書き換えることができる。展開すると
\[
d^{2} x^{2} + 2dexy + e^{2} y^{2} + 2gx + 2fy + c;
\]
よって比較して$a = d^{2}$, $h = de$, $b = e^{2}$となり、従って$ab - h^{2} = d^{2} e^{2} - (de)^{2} = 0$である。したがって必要条件は自動的に満たされる。
$a = b = g = f = 0$, $h = 1$, $c = -4$と置く方程式$2xy - 4 = 0$は双曲線を表す。条件$ab - h^{2}$は$0 - 1^{2}$となり、つまり-1であり、負である。

一般の方程式がどのような軌跡を表すかによって定まる制限は必要である。なぜなら、一般的な方程式の特定の場合には実際的な軌跡を表さない場合があるからである。例えば$x^{2} + y^{2} + 1 = 0$は実数の$x$と$y$に対して満たされない。通常はこの軌跡が虚点で構成されるものだと述べるが、幾何学における虚点という概念は非常に複雑なものであり、ここでは詳しく説明しない。

Some exceptional cases are included in the general form of the equation which may not be immediately recognized as conic sections. By properly choosing the constants, the equation can be made to represent two straight lines. Now two intersecting straight lines may fairly be said to come under the Greek idea of a conic section. For, by referring to the picture of the double cone above, it will be seen that some planes through the vertex \( V \) will cut the cone in a pair of straight lines intersecting at \( V \). The case of two parallel straight lines can be included by considering a circular cylinder as a particular case of a cone. Then a plane which cuts it and is parallel to its axis will cut it in two parallel straight lines. Anyhow, whether or no the ancient Greek would have allowed these special cases to be called conic sections, they are certainly included among the curves represented by the general algebraic form of the second degree. This fact is worth noting; for it is characteristic of modern mathematics to include among general forms all sorts of particular cases which would formerly have received special treatment. This is due to its pursuit of generality.

\chapter{第11章：Chapter XI: Functions}

数学における関数という用語は、日常生活でも採用されています。例えば、「彼の気分は消化状態による関数である」という表現は、この数学的な意味で正確に使われています。これは、彼の気分がどのように消化状態によって決まっているかを知るルールがあることを意味します。つまり、「関数」の概念は単純で、変数の数値に対してどのように適用されるかを見るだけです。

まずいくつか具体的な例を考えましょう：ある列車が時速20マイルで走行している場合、任意の時間$t$（時間）後の距離$s$（マイル）は$s = 20 × t$で与えられます。そして$s$は$t$の関数と呼ばれ、$20 × t$は$s$と同一の関数です。また、ジョンがトーマスより1歳年上である場合、トーマスの年齢$x$（年）が任意のとき、ジョンの年齢$y$（年）は$y = x + 1$で与えられます。そして$y$は関数$x + 1$と同一です。

これらの例では、$t$と$x$を関数に現れる「引数」と呼びます。したがって$t$は関数$20 × t$の引数であり、$x$は関数$x + 1$の引数です。もし$s = 20 × t$で、$y = x + 1$である場合、$s$と$y$はそれぞれ関数$20 × t$と関数$x + 1$の「値」と呼ばれます。

現在、一般の場合について述べるが、数学における関数を
変動する二つの数の間の相関として定義できる。それぞれを「関数の引数」と「関数値」呼ばれるものとするとき、関数の引数にどのような値を代入しても、関数値は一意に（すなわち独特に）決定される。逆は必ずしも成り立たない。つまり、関数値が決定されても、その引数が一意に決定されるわけではない。
引数 $x$ の他の関数として $y = x^{2}$, $y = 2x^{2} + 3x + 1$, $y = x$, $y = \log x$, $y = \sin x$ がある。このグループの最後二つの関数は、少しの代数学と三角法を理解している人にはすぐに識別できるだろう。その説明に時間を費やす価値はなく、例示のためだけに引用されているものである。

これまでまで、関数の一般的な定義は述べましたが、特定の関数の一連のみを挙げてきました。しかし、数学は一般的な方法で進行するため、任意の関数の一般的な概念を表すために符号化します。これは、$x$ の関数として任意の関数を表すために $F(x)$, $f(x)$, $g(x)$, $\phi(x)$, などと記述することで達成されます。ここで、引数$x$は括弧に入れられ、$F$, $f$, $g$, $\phi$ などのような文字が括弧の前に置かれることで関数を表します。

この表記法には欠点があります。つまり、これは明らかに単一の文字が変動する数字を表す慣習と衝突します；ここでは、括弧の前に置かれた $F$, $f$, $g$, $\phi$ などが変動する関数を表しています。意味が明確でない例を挙げると、直感や文脈に頼らなければなりません。混乱を避けるための一法は、関数のためにギリシャ文字（例えば上の $\phi$）を使用することです；また、関数の文字として $f$ と $F$ （関数の頭文字）を使用し、他の変動する関数が必要な場合は隣接する文字（例：$g$）を用いることです。

これらの説明と注意事項に基づき、$y = f(x)$ と書きます。これは $y$ が引数 $x$ の未定義の関数の値であることを表します；ここで $f(x)$ は $x + 1$, $x^2 - 2x + 1$, $\sin x$, $\log x$, または単に $x$ 自身を表すことができます。重要な点は、$x$ が与えられたときには $y$ もそれに従って確定的に定まるということです。この考えの一般性について十分に理解することが重要です。したがって $y = f(x)$ の場合、例えば $f(x)$ を $x$ が整数であるとき $f(x) = 0$ で、それ以外の値を持つとき $f(x) = 1$ とすることができます。この意味での $f(x)$ を用いて $y = f(x)$ と置くと、$x$ の値が整数であるときは $y = 0$ であり、それ以外の場合は $y = 1$ となります。したがって $f(1) = 0$, $f(2) = 0$, $f(\frac{2}{3}) = 1$, $f(\sqrt{2}) = 1$ などです。この $f(x)$ の意味の選択は、引数 $x$ の関数として一般的な定義に従って完全に良い関数を提供します。

結局のところ関数は、二つの変数間の一種の相関関係であるため、他の相関関係同様座標幾何学的なグラフで表されます。例えば第II章のFig. 2 は、引数 $v$ と関数の値 $p$ の関数 $\frac{1}{v}$ のグラフです。この場合、グラフは正の $v$ のみを描いており、その物理的な応用に意味を持つ唯一の値です。また第IX章のFig. 14 では、線分 $AB$ 全体（両方向に無限延長）が引数 $x$ と関数の値 $y$ の関数 $x + 1$ のグラフであり、同じ図において無限な線分 $A_1B$ は関数 $1 - x$ のグラフであり、線分 $LOL'$ は引数 $x$ と関数の値 $y$ の関数 $x$ のグラフです。

これらの関数は単純な代数的公式で表され、グラフによって表現できる。しかし、一部の関数については、詳細な説明なしではそのグラフによる表現は非常に誤解を招くものであり、甚至不可能である場合もある。例えば、引数$x$が整数以外のすべての値に対して値が$1$となるが、$x = 0, x = 1, x = 2,\ldots$のような整数では値が$0$となる関数を考える。この関数のグラフ上の表現は、$XOX'$軸から$1$単位離れた平行線である直線$ABA'$に相当する。しかし、$x$の値が$0, 1, 2, 3, 4,\ldots$に対応する点$B, C_1, C_2, C_3, C_4,\ldots$は省略され、代わりに軸$OX$上の点$O, B_1, B_2, B_3, B_4,\ldots$が取り入れられる。グラフによる表現が不便であるだけでなく不可能な関数も存在し、そのような関数は高次の数学において重要であり、ここではそれらについてさらに詳しく触れることはない。

関数の最も重要な分類は連続関数と不連続関数との区別にあり\index{連続関数|EtSeq}\index{不連続関数|EtSeq}。引数が徐々に変化するときその値も徐々に変化する関数を連続関数、突然の跳躍によって値が変わる関数を不連続関数と呼ぶ。図\[fig.]{14}\ の第9章で示したグラフである$x + 1$と$1 - x$は連続関数であり、第2章で示された$\dfrac{1}{v}$も正の$v$の値を考える限りでは連続関数である。しかし、この章の図\[fig.]{20}\ に示した関数は不連続関数であり、その引数$x = 1, x = 2,\ldots$の値において突然の跳躍が生じるからである。

自然に与えられた関数の例を考えましょう。これにより連続性と不連続性の実質を頭に入れることが出来ます。列車がロンドンのエストン駅からラノーウェスト・レイルウェイの終点であるラノーウェスト駅までの鉄道線路上を旅しているとしましょう。その線上にブレッチリーとラビジの駅があります。$t$ を列車がエストン駅を出発してから経過した時間を表す数値、$s$ を列車が通過した距離を表す数値とします。すると$s$は$t$の関数であり、つまり変化する引数$t$に対する変化する値に対応しています。\PageSep{151}

もし列車の運行情報を知っているなら、特定の$t$の値が与えられればすぐに$s$も分かります。しかし、奇跡とは別に、$s$が$t$の連続的な関数であることを確信的に仮定することができます。列車がエストンからブレッチリーまで継ぎ接ぎなく通過し、その後わずかな時間も無くラビジに現れることはあり得ません。その考え方は私たちの計算には適さないほど幻想的です：それは《アラビアン・ナイト》以外では見つけることのできない可能性を含みます；すらそれすらは運動の不連続性について想像することを恐れるのです。彼らは私たちに信じさせること以上の非常識な速さを提示することさえ恐れています。しかし、非常に不思議な速度は自然の動きの大きな法則である連続性と矛盾しません。光は約$190,000$マイル毎秒で動くため、太陽からの到達時間は7分から8分程度ですが、その速さにもかかわらず、その移動距離は常に時間の連続的な関数です。

私たちには、物体の速度が時間の連続的な関数であることが直感的に明らかではない。任意の時刻 \( t \) における列車を考える：その列車は一定の速度 \( v \) マイル每時間で動いている。\( v \) は列車が駅で静止しているときゼロになり、後退しているときは負の値を取る。大きな重い列車について考えてみると、\( v \) の突然の変化は許されないことが明らかだ。列車は午前11時45分から正午まで40マイル每時間で走行し、その後瞬時に50マイル每時間に加速することはあり得ない。私たちは即座に、速度の変化は徐々な過程であると認めざるを得ない。

しかし、十分な大きさの突然の衝撃はどうだろうか？ 例えば、2つの列車が衝突する場合を考える；または小さな物体として、人がフットボールを蹴る場合を考える。私たちの感覚には、フットボールが突然動き出すように見える。したがって、速度に関しては、私たちの感覚はその時間の不連続的な関数であるという考えに反対しない。列車が一瞬でブレッチリーからラグビーへと移動するという考えとは異なり、実際には運動の法則と質量の概念が正しいなら、自然には不連続な速度は存在しない。私たちの感覚上では不連続的な速度変化に見えるものは、実際には私たちはそれを認識できないほど速い徐々な変化であると考えるべきだ。

しかし、自然界で不連続関数が現れることが絶対にないという一般化は早計である。ロンドンとパリ間の海抜平均高さが距離から連続的な関数であることを信じて夜にシェイクスピアのドーバー崖を歩いた男は、科学的結論に対する注意が必要なことについて考え直す前に死んでしまうだろう。

したがって、自然には不連続な関数が現れる可能性があるという結論は慎重に検討すべきである。

連続でない関数を見つけることは非常に簡単です。代数式の最も単純な形でも同様です。例えば、関数 $y = \frac{1}{x}$ を考えてみましょう。これはすでに $p = \frac{1}{v}$ の形で $v$ が正の値に限定されているときを考えたことがあります。しかし今度は $x$ が任意の値、つまり正でも負でもあるとします。関数のグラフは \Fig[fig.]{21} に示されています。

$x$ が非常に大きな負の値から始まり、逐次的に減少する一連の負の値を経て $0$ まで変化し、その後増加する一連の正の値を通ります。したがって、動点 $M$ が $XOX'$ 上で $x$ を表すと、$M$ は軸 $XOX'$ の左端から始まり、順次 $M_{1}$, $M_{2}$, $M_{3}$, $M_{4}$, 等を通過します。これらの点に対応する関数の点は $P_{1}$, $P_{2}$, $P_{3}$, $P_{4}$, 等です。容易に理解できるように、$x = 0$、すなわち原点 $O$ において関数が不連続であることがわかります。原点の左側（負の側）では関数の値は無限大で負となりますが、右側（正の側）では無限大で正になります。したがって、$M_{2} M_{3}$ の長さをいかに小さくしても、関数の値において有限な跳躍があります。実際、この例には $M_{2}$ と $M_{3}$ の間の距離をどれほど小さくしても（原点を含む限り）、関数の値がそれらの間でより大きな跳躍を示すという特徴があります。このグラフは、この章の \Fig[fig.]{20} でも明らかですが、多くの関数では不連続性は孤立した点にのみ発生し、引数の値を制限することで残りの点で連続な関数を得ることができます。したがって、\Fig[fig.]{21} から $y = \frac{1}{x}$ の場合、正の値のみを取り入れて原点を除外すると、連続な関数を得ることができます。同様に、負の値のみを取り入れて原点を除外すると、同じ関数も連続になります。また、\Fig[fig.]{20} で示されたグラフは $B$ と $C_{1}$ の間、そして $C_{1}$ と $C_{2}$ の間、さらに $C_{2}$ と $C_{3}$ の間など、常に端点を除いて連続です。しかし、すべての点で不連続性が発生する関数も簡単に見つけることができます。例えば、$x$ が任意の分数のとき $f(x) = 1$ であり、$x$ が任意の無理数のとき $f(x) = 2$ のような関数を考えます。この関数はすべての点で不連続です。

最後に、上記で与えられた連続性の定義についてもう少し詳しく見ていきましょう。
関数が値を徐々に変化させるためには徐々に引数が変化する必要があると言いましたし、突然の跳躍によって値が変わるときには不連続であるとも言いました。これはまさに当時の数学的先駆者たちが満足した定義であり、現代の数学家たちはそれに満足していません。この定義を理解するために時間を費やす価値があります；なぜなら、その現代的な反論を理解するときには、現代数学の精神についても多くの理解を得られるからです。
古い数学と新しい数学の違いは、曖昧な半分の言説的な用語「徐々に」が正確な表現において容認されない点にあるのです。現代の数学は、科学を基礎とする少数の単純な数や量、変数に関する概念のみを使用する声明と定義と議論だけを認めます。
二つの数のうち一つが他の数より大きいか小さいか、あるいは一つが他の数の何倍かを特定できるかもしれませんが、二つの数間には「徐々に」の関係はありません。そのためこの用語は不適切です。
最初にこれを大げさな形式主義だと考えるかもしれません。しかし、そのような批判に対して二つの答えがあります。まず第一に、19世紀前半にいくつかの大数学者が、特にスウェーデンのアベルとドイツのヴィエトリスによって指摘されたように、古い楽観的な方法で述べられた多くの数学的結果は単なる誤りでした。
マカーギルの「 Bacon 論文」では、数学の確実性と哲学の不確実性を対比し、「テイラーの定理に対する反動が存在しない」という例として引用しています。彼が選んだ最悪の例です。この論文発表当時の英語圏の数学教科書を調査せずに、テイラーの定理はすべての教科書で誤って述べられていたという仮定は比較的安全でしょう。
したがって、正確性のために現代の数学の緊張した精密さが必要です。二番目には研究にとって必要です。それは明瞭な思考を促し、新しいアイデアの組み合わせを試すための大胆さと豊かさを生むからです。初期の声明が曖昧で手抜きである場合、その後のすべての思考段階では普通の感覚が適用範囲を制限し、意味を説明しなければなりません。
しかし創造的な思考においては、普通の感覚は悪い主人です。その唯一の判断基準は新しいアイデアが古いものに見えて見えることです。つまりそれは本来性を抑圧することによってしか作用しないのです。

精密な連続性の定義に向かって進むとき、数との「段階的な」関係が存在しないという主張をより詳しく検討しましょう。一つの数が別の数よりもわずかに大きいとは言えないでしょうか、言い換えれば、二つの数の差が小さいとは言えないでしょうか？その全体のポイントは、抽象的に見れば、任意の仮定的な応用を除いて、大きな数や小さな数というものは存在しないということです。百万マイルは天文学者が固定星を調査する際には小さな距離であり、百万ポンドは年収としては大きな所得です。また、寄付として自分の収入の四分の一は慈善のために与えるのは大きいですが、個人的な使用のために残すのは小さいのです。このような例を無限に積み重ねることで、絶対的な意味での大きさや小ささは数に対して抽象的に適用されないことが示されます。二つの数について言えることは、一つがもう一つよりも大きいまたは小さいということだけです。具体的な状況を指定せずに、どの数が大きいか小さいかを述べることはできません。したがって、関数の値の変化における「小さな」または「段階的な」変化を一切言及することなく連続性を定義することが私たちの仕事となります。

これを達成するために、いくつかのアイデアに名前をつけます。これは、極限や微分法を考える際にも有用です。

引数$x$の関数$f(x)$の「区間」\index{Interval|EtSeq}は、ある二つの引数の値間に存在するすべての値からなります。例えば、$x=1$と$x=2$の間の区間は、$x$が$1$と$2$の間で取れるすべての値からなります。つまり、それは$1$と$2$の間に存在するすべての実数からなります。しかし、区間の境界となる数は整数である必要はありません。引数の「包含」する数$a$があるとき、$a$が区間に属しているということです。例えば、$1$と$2$の間の区間には$\frac{3}{2}$、$\frac{5}{3}$、$\frac{7}{4}$などが含まれます。

一組の数が数$a$に標準$k$での近似値として近似されるとは、$a$とその一組の数との数的差が$k$未満であるときを指す。ここで$k$は「近似標準」と呼ばれる。例えば、数$3, 4, 6, 8$の一組は数$5$に標準$4$での近似として近似される。この場合、選べる最小の標準とは限らず、その一組は標準$3.1$や$3.01$や$3.001$でも$5$に近似される。また、数値$3.1, 3.141, 3.1415, 3.14159$は標準$.032$での$3.13102$への近似と同時に、より小さな標準$.03103$でも近似される。

これらの二つの考え、すなわち区間と数$a$に対する標準$k$での近似という概念は単純で、その唯一の難しさはそれらが非常に平凡に思えることである。しかし、次の考え、すなわち「数の近傍」という概念と組み合わされると、現代の数学的論理の基礎となる。何かを主張するとき、「引数$x$の値$a$における関数$f(x)$について何かが真である」と言うとは何を意味しているのか？これが我々が現在明確にするべき基本的な概念である。

ある関数$f(x)$の値が「$a$の近傍」において特性を持つとは、ある区間が存在し、(i)この区間に数$a$は端点として含まれず、(ii)その区間内の引数が$a$でないすべての値に対して、これらの関数の値が該特性を持っているときを指す。関数$f(x)$の引数$x=a$における値$f(a)$はこの特性を持つかどうかは確定していない。これは$a$の近傍に関する主張によって決定されるものではない。

例えば、特定の関数$x^{2}$を取ります。この関数について、$2$の近傍では値が$5$未満です。なぜなら、$1$から$2.1$までの区間は次の条件を満たすためです。(i) $2$は端点として含まれていませんし、(ii) この区間に属する$x$の値に対して$x^{2}$が$5$未満であるようなものです。

現在、前節の考えを組み合わせると、関数$f(x)$が点$a$の近傍で$c$に標準$k$での意味で近似することとは、ある区間が存在し、(i) $a$はその端点ではなく含まれており、(ii) その中の$x$が$a$でないすべての値について、$f(x)$の値が$c$から$k$未満の差異を持つことを意味します。例えば、$2$の近傍では、関数$\sqrt{x}$は標準$.0001$での意味で$1.41425$に近似します。これは、$\sqrt{1.99996164}=1.4142$かつ$\sqrt{2.00024449}=1.4143$であるためです。したがって、$x$が$1.99996164$から$2.00024449$の区間にあり、$a=2$はその端点でない場合、関数$\sqrt{x}$の値はすべて$1.4142$と$1.4143$の間であり、したがってこれらは全て$c=1.41425$から$.0001$未満の差異を持つことになります。この場合、もし希望すれば、より小さな近似の標準を固定できます。例えば、$.000051$や$.0000501$です。

また別の例として、$2$の近傍では関数$x^2$は標準$.5$での意味で$4$に近似します。これは$(1.9)^2=3.61$かつ$(2.1)^2=4.41$であるため、必要とする区間$1.9$から$2.1$が見つかりました。この例は、関数$f(x)$の値が点$a$での近傍における声明とは異なることを示しています。真偽を保証する区間が必要です。したがって、$2^2=4$という単なる事実だけでは、$2$の近傍で関数$x^2$が$4$に等しいという声明は正当化されません。これは、必要な性質を持つ区間を生成できないためです。

また、$2^2=4$という単なる事実は、$2$の近傍で関数$x^2$が標準$.5$での意味で$4$に近似しているという声明も正当化しない。しかし実際には、この声明はすでに証明されました。

もし前段の考えを理解すれば、現代数学の基礎も理解できるでしょう。系列に関する章で類似の概念に戻り、また微分法に関する章でも再び扱います。

「連続関数」というものを定義しましょう。「関数」$f(x)$がその引数の値$a$において「連続」であるとは、$a$の近傍においてその値が$f(a)$（すなわち$a$における値）に標準的などの程度でも近づくことを指します。

これは、任意の標準$k$を選んでも、$a$の近傍において$f(x)$が$f(a)$に標準$k$以内で近づくことを意味します。例えば、$x^2$はその引数$x$の値$2$において連続です。なぜなら、どの程度$k$を選んでも、(i) $2$を端点として含まない区間(ii) その内にある引数における$x^2$の値がすべて$k$以内で$4$（すなわち$2^2$）に近づくような区間を見つけることができます。例えば、標準$k = 0.1$を選ぶとします。このとき$(1.999)^2 = 3.996001$であり、$(2.01)^2 = 4.0401$で、これらはどちらも$4$から$0.1$未満の差異があります。したがって、区間$1.999$から$2.01$までの$x^2$の値は標準$k = 0.1$以内で$4$に近づきます。同じように、任意の他の標準を試したいときにも同様の区間を作ることができます。

鉄道列車の例を取りましょう。信号ボックスを通過するとき、任意の速度（例えば1マイル每時間の百万分の一）を選んでも、その瞬間から前後一定の時間区間に、その区内のすべての瞬間において列車の速度が信号ボックスを通過した時の速度から百万分の一マイル每時間以内に異なるという区間を見つけることができます。そして、他の任意の速度を代入しても同じことが言えます。

\chapter{第12章：Chapter XII: Periodicity in Nature}

自然の全生涯は、周期的な出来事の存在によって支配されています。すなわち、類似性が高く互換的に繰り返される一連の出来事が存在することです。言葉を引張ることなく同じ出来事の再現と呼び得るようなものです。地球の回転は、連続する日を生み出します。確かに、どの程度抽象的に「日」と定義しても、各日は前の日に異なるものであることは否定できません。しかし、「日」の意味を十分に抽象化すれば、二つの日の性質的な違いは実用的な興味から遠ざかり、それぞれの日は地球の1回の回転の現象の再現として考えることができます。

また、地球が太陽周りを巡る経路は年間の季節の再現をもたらし、自然全体のすべての操作に別の周期性を課します。月相によるもう一つの基本的でない周期性もあります。現代の文明社会では、人工的な光によってこれらの月相は微々たるものですが、古代時代、日が灼熱し晴れ渡った気候の中で人間の生活は月明かりの存在に大きく影響されたようです。そのため、週や月による分類とその宗教的関連性はシリアや美索不達ミアからヨーロッパ諸民族へと広まりました。しかし、多くの国々では、月相に基づく独立した儀礼が見られます。

しかしながら、月の周期性は主に潮汐を通じて地球の歴史に大きな影響を与え、その明暗による周期性よりも重要な役割を果たしました。

私たちの身体的生活は本質的に周期的である。
心臓の拍動と呼吸の再現がその主導を握っている。
周期性の前提は、実際上生命自体の概念に根本的なものである。
我々は進展する出来事に対して、「これ以前にも起こっていた」と言える状況を想像できない。
経験が行動指針となるためには、過去の歴史と類似した新しい状況に常に直面することもなければ、時間量を測定する手段もなくなる。
ある出来事が他のよりも前後関係にあるという単純な認識を超えて、私たちは次のように言えることができる：
三つの出来事 $A$,~$B$,~$C$ がこの順序で起こったので、$A$ は $B$ より前に、そして $B$ は $C$ より前に起こった。
さらに、$A$ の発生から $B$ の発生までの時間の長さが、$B$ の発生から $C$ の発生までの時間の長さの二倍であることも言える。
現在の時間量は、自然な再現が介入した数を観察することに本質的に依存している。
私たちは、$A$ と $B$ の間の時間を、何日、何ヶ月、または何年など、我々が参照したいタイプの再現によって、そのような数だけと言える。
実際には、文明が始まった段階では、これらの三つの時間測定方法は本当に異なるものだった。
科学は文明や半文明国家にとって最初の重要な任務の一つであり、これらを一貫した計測に統合することである。
この任務の全貌を理解することが必要だ。単にどの数日（例えば $365.25$ など）が某一年に相当するかだけでなく、さらに前もって次の年々にも同じ数日の長さが適用されることを決定しなければならない。
周期性が存在する世界を想像できるが、その周期性は互いに整合的ではない場合もある。ある年の間に $200$ 日で、他の年の間には $350$ 日であるなどということがあり得る。
より重要な周期性の広範な一貫性の決定は自然科学の最初の段階だった。この一貫性は抽象的な直感的法則ではなく、経験によって保証された自然の観察事実によるものだ。
実際には、これは必要法則ではないし、正確には真でもない。すべての場合において、いくつかの逸脱が存在する。一部の例では、これらの逸脱は容易に観察され、すぐに明らかになる。他の場合では、最も精密な観測と天文的な正確性が必要となる。
一般的には、生物に関連した再現（例えば心臓の拍動）は他の再現と比較して急速に変化する。大きな安定性を有し、互いに非常に正確に一致する重要な再現は、地球全体の運動や類似の天体の運動によるものである。

我々はしたがって、これらの天文的な再現が等しい時間間隔を表していると仮定する。

しかし、精密な天文観測によって検出されるその不一致はどう対処すればよいのか？明らかに私たちは、これらの現象のいずれかが等しい時間を表すという随意的な仮定に至る。例えば、すべての日が等長であるとするか、またはすべての年が等長であるとするか。

しかし、そうではない：何かしらの仮定が必要だが、天文学者が時間の尺度を決定する全過程において基礎となる仮定は、運動の法則が完全に遵守されているという点にある。

この方法がどのように機能するかを説明する前に、そのような上位の一貫性が人間の体の特性的な再現と一致していた場合、時計の調整は医師たちに委ねられたことだろうという興味深い観察を行うのは面白い。

運動の法則がどのようにして関係するかを考える際に、時間の測定方法が異なる場合、同じ物体に対して異なる速度変動を得ることになります。例えば、1時間は1日の中の24分の1と定義し、列車が2時間にわたって平均速度で20マイル毎時で走行するという状況を考えてみましょう。次に、非常に不一致な時間測定法を考え、最初の1時間が次の1時間の2倍長であるとする場合があります。すると、この別の時間尺度では、列車の走行時間が二つの部分に分かれ、各部分で同じ距離、すなわち20マイルを走行しますが、最初の部分の持続時間は次の部分の持続時間の2倍です。したがって、列車の速度は一貫して均一ではなく、平均的には第二部分での速度は第一部分での速度の2倍となります。つまり、列車が一貫して均一に走行しているかどうかは、採用する時間基準によって異なります。

地球上の日常生活におけるすべての通常の目的において、様々な天文的再発現は絶対的に一致していると見なすことができます。さらにこれらの一致性を仮定し、それによって物体の速度や速度変化も仮定すると、上記で考慮した運動の法則がほとんど正確に確認されます。しかし、いくつかの天文的な現象に関しては、完全には正確ではありません。しかし、惑星と恒星の回転および運動のわずかに異なる速度を仮定することで、法則が正確に確認されることを見つけることができます。この仮定に基づき、時間の尺度を定義し、特定の一つの周期的現象とは一貫していないように定義します。しかし、重要な事実は残ります。つまり、多くのことが基盤となっている均一な時間の流れは、周期的なイベントの観測に依存しているという事実です。

現象は表面的に見れば偶発的で例外的なものや、一方では一定の持続性を保つものであっても、遠隔な周期性の影響により生じる可能性がある。例えば、共振の原理を見てみよう。共振は、二つの関連する状況が同じ周期を持つときに生じる現象である。

共振は動力学法則であり、すべての物体が自己の振動は特有の一定の時間で行われることがある。つまり、小さな振幅での単振り子はその形状や重心分布、長さによって特有の一定の時間を常に振動する。より複雑な物体は多くの振動方法を持つかもしれないが、それぞれの振動モードには独自の「周期」がある。これらの体の振動の周期は「自由」周期と呼ばれる。したがって単振り子は一つの周期しか持たない一方で、吊り橋は多数の周期を持つ。

弦楽器のような音響機器を得るためには、すべての振動周期が最も長い周期の単純な倍数である必要がある；つまり、$t$秒が最も長い周期であれば、他の周期は$\frac{1}{2}t$, $\frac{1}{3}t$などであり、これらの小さな周期のいずれかが欠けていることもあります。

今、物体の振動を自己以外の周期的な原因で激しく引き起こすと仮定する。その周期が物体の周期に非常に近ければ、その振動モードは非常に激しく引き起こされる；これは原因の大きさに関係なくである。この現象は「共振」と呼ばれる。

共振の一般的な理由は容易に理解できる。例えば、岩を倒す人たちは常に有利な瞬間に推しやすいように、岩の揺れに合わせて推す。もし推しが調子が合わなければ、一部は揺れを増幅する一方で他の部分は制御する。しかし、調子が合うと、時間とともにすべての推しが有利になる。

「共振」という言葉は音響に関する考慮から来ているが、その現象は音響以外にも広く及ぶ。光の吸収や放出の法則、無線電信受信機の「チューニング」、惑星間の運動への影響、吊り橋上の軍隊が一列に歩くときの危険性、特定の速度で機器のリズミカルな拍動による船の過度な振動など、共振は広範囲にわたる現象を生む。この周期性の一致は恒常的な二つの周期的イベントの関連があるときには安定した現象を生み出すが、偶然の一時的な関連では激しい突然の爆発を引き起こす可能性もある。

再度、上記で述べた特性的かつ恒常的な振動周期は、私たちの感覚に恒定した刺激として見受けられる根本的原因です。数時間にわたって恒定した光の中で作業をしたり、恒定して変化しない音を聴いたりします。しかし、現代科学が正しいとすれば、この恒定性は自然の中にその相応するものは存在しません。恒定した光は、振動エーテルからの無数の周期的波による目への衝突によって引き起こされ、恒定した音は同様の波により引き起こされます。ここでは光の理論や音の理論を説明することではありません。私たちが述べたように、数学が自然の調査に適切な道具となるためには、物事の根本的な周期性を表現することができることが必要であることが明らかになりました。この理解を得た上で、次に考慮すべき重要な数理概念、すなわち周期関数の重要性を理解することができます。

\chapter{第13章：Chapter XIII: Trigonometry}

初等三角法は、自然の周期性に対する一般的な考慮から発展したわけではありません。この点では、その歴史は双曲線などの二次曲線の起源と類似しています。実際、二つの科学の歴史を比較すると、非常に教益的な類似点と対比が見られます。三角法も双曲線同様、ギリシャ人によって発明されました。その発明者はヒッパルコスで、紀元前160年頃に生まれたギリシャの天文学者です。彼はロデ島で観測を行いました。彼の天文への貢献は非常に大きく、彼の手から科学的な分野が真の科学として成立し、正しい発展方法が示されました。おそらく三角法の発明も、彼が研究した主要な科学に対する重要なサービスの一つでした。

次に三角法を拡張した人物は、アレクサンドリアの伟大的な天文学者ポタメデウス（プトレマイオス）です。すでに述べた通りです。ここで直ちに双曲線と三角法の大きな対比が明らかになります。三角法の起源は実用的であり、それは天文学的研究のために必要だったから発明されました。一方で、二次曲線の起源は純粋理論的でした。その最初の研究の唯一の理由は、関連する概念の抽象的な興味でした。特徴的に、二次曲線は三角法よりも約150年前に、ギリシャ思想の最も良い時期に発明されました。しかし、三角法が数学の理論と応用にとってどれほど重要であるかは、一般科学から実用的適用によって得られた無数の有効なアイデアの一例です。

私たちは三角法が何であるかを理解し、その生成は科学的な天文学研究によってなされるべき理由について明確にする試みを行います。\index{天文學}

まず最初に：天文学者が行うことができる測定は何でしょうか？それは時間の測定と角度の測定です。天文学者は（現代の天文学者にとって馴染みのある）望遠鏡を調整することができます。その望遠鏡は東西南北を指す固定軸を中心に回転するように設定できます。結果として、望遠鏡は南に向けてしか指向できず、方向がより仰角に近ければ高い位置へ向かうか、天頂を超えて北に向かうことができます。これが軌道観測器で、正確な星座の南または北への通過時間の測定のために大きな役割を果たします。しかし、この観測器は角度も測定します。なぜなら、二つの星の通過間隔が記録された後、地球の均一な回転の仮定により、その期間における地球がどの程度回転したかの角度が得られるからです。また、他の装置を使用して、二つの星間の角度を直接測定することも可能です。例えば、観測者の目である$E$と、星を見た方向である$EA$と$EB$を考えます。角度$AEB$を測定するための機器は容易に作ることができます。

したがって、天文学者が天体の調査を行う際には、実質的には瞬間的な各星や惑星の相対的方向を角度として測定していることになります。また、土地測量における類似問題においても、角度は主な測定対象となります。長さの直接的な測定は非常に正確に行うことが稀であり、河川、建物、森林、山脈、そして一般的な地面の不規則性が邪魔をします。一国全体の測量は、サリー平野のような選ばれた場所で最大限の精密さで行われた数回の長さの直接的な測定に依存します。主な作業は角度の測定です。例えば、$A$、$B$、および$C$は調査区域における目立つ点であり、教会塔の尖端を示すものとします。これらの点は互いに見えます。すると、$A$から角$BAC$を測定し、$B$から角$ABC$を測定し、$C$から角$BCA$を測定することは非常に簡単です。理論的には、三角形の三つの角度の合計が二直角に等しいという有名な幾何学的命題により、二つの角度が分かれば残りの角度は推定できます。しかし実際には、すべての三つの角度を測定することが好ましく、観測上の小さな誤差をチェックすることができます。

地図作成過程では、この方法で一国全体が三角形で覆われます。このプロセスは接続測量と呼ばれ、測量における基本的なプロセスです。\index{接続測量}

すべての三角形の角度が分かれば、その形状は分かることになります---つまり、サイズとは区別して形状。ここで幾何学的類似性の大原則に出会います。この概念は、私たちの実用的な応用において非常に馴染みがあります。スケール図を描くというアイデアがよく知られています。例えば、プランのスケールが1インチに対し1ヤードであれば、プラン上の3インチは元の長さの3ヤードを意味します。また、プラン上に描写された形状は元の形状と同じであり、つまり元の直角はプラン上で直角として現れます。同様に、地図も国を描いたプランであるため、地図上の距離比は実際に示される場所間の距離比であり、地図上の方向は実際の方向です。例えば、地図上での一つの場所が他の場所に対して北北西にある場合、それは実際でも同じです。つまり、地図上で角度は実際と同様です。

地図（またはプラン）と元の形状における対応する線分の長さ比が固定された比率であるとき、そして対応する角の大きさが等しいとき、二つの図形は類似していると言えます。この固定比率は地図（またはプラン）と元の形状における対応する線分の長さ比を「地図のスケール」と呼びます。地図やプランには常にスケールが表示されるべきです。

角度がそれぞれ等しい二つの三角形は類似していることが既に述べられました。したがって、三角形$ABC$と$DEF$において、角$A$と$D$、角$B$と$E$、角$C$と$F$がそれぞれ等しければ、線分$DE$は線分$AB$に対する比と同様に線分$EF$は線分$BC$に対する比、そして線分$FD$は線分$CA$に対する比となります。しかし、他の図形においては角度が等しいだけで類似性が保証されるわけではありません。例えば、正方形と長方形の例を考えてみましょう。正方形$ABCD$と長方形$ABEF$を考えると、すべての対応する角は等しくなりますが、線分$AB$は正方形と長方形で同じ長さですが、線分$BC$は線分$BE$の半分程度の大きさです。したがって、正方形$ABCD$は長方形$ABEF$に類似しているとは言えません。

三角形には他の直線図形とは異なる独特の性質があり、これが相似性理論における基本的な図形となっています。そのため、測量では三角法が基本的なプロセスとなり、「trigonometry」（三角法）という用語もこの二つのギリシャ語から派生しています。

円についても同様に述べられます。任意の2つの円は相似な図形であり、その違いはスケールのみです。円周の長さ比は半径の長さ比と比例します。さらに、2つの円が同じ中心を持つ場合（図26参照）、角度$AOP$によって切り取られた弧$AP$と$A_{1}P_{1}$もその半径に対する比で比例します。したがって、弧$AP$の長さと半径$OP$の長さの比，すなわち$\dfrac{\text{弧 } AP}{\text{半径 } OP}$は$OP$の長さによらず一定の数値であり、$\dfrac{\text{弧 } A_{1}P_{1}}{\text{半径 } OP_{1}}$と等しいです。この「弧を半径で割った比」が角度の大きさを測る理論的な方法であり、任意の長さ単位や任意に決められた角度分割法によらないため、より正確な測定方法となります。

角$AOP$の大きさを表す分数$\dfrac{AP}{OA}$を考えます。次に$PM$を$OA$に対して垂直に引きます。すると、ギリシャの数学では線分$PM$を弧$AP$の正弦、線分$OM$を弧$AP$の余弦と呼びました。彼らはこれらの線分間の関係が我々が先ほど述べた相似性理論に基づいて重要であることをよく理解していましたが、その性質をこの理論から直接定義することはしなかったし、現代的な関数の概念や代数学的概念も持っていました。

したがって、我々にはより強力なアイデアを反映させることが自然です。

現代の数学では、弧$AP$を考慮する代わりに、分数$\frac{AP}{OP}$を考えます。これは$OP$の長さによらない定数です。また、線分$PM$と$OM$を考慮する代わりに、分数$\frac{PM}{OP}$と$\frac{OM}{OP}$を考えます。これらもまた$OP$の長さによらない定数であり、つまり図形のスケールによらない定数です。

次に、数$\frac{PA}{OP}$を角度$AOP$の大きさとして定義し、分数$\frac{PM}{OP}$をこの角度に対する正弦（\emph{sine}）、$\frac{OM}{OP}$を余弦（\emph{cosine}）と定義します。これらの分数形式は印刷に不便なので、$\frac{AP}{OP}$を$u$、$\frac{PM}{OP}$を$v$、$\frac{OM}{OP}$を$w$と置きます。すると、$u$、$v$、$w$は数であり、任意の角度$AOP$について変動する定数です。

しかしながら、これらの数値には関連性があり、$u$（すなわち角度$AOP$）が与えられたとき、$v$と$w$の大きさは明確に決定されます。したがって$v$と$w$は引数$u$の関数です。$v$を正弦（\emph{sine}）、$w$を余弦（\emph{cosine}）と呼びます。

一般の関数表記$y = f(x)$にこれらの特殊な場合を適用するため、現代の数学では、特別な「正弦」関数を示すために$\sin$を用い、特別な「余弦」関数を示すために$\cos$を使用します。したがって、上記の意味で$u$, $v$, $w$について、

\[
v = \sin u,\quad w = \cos u,
\]

と表せます。特殊な関数の場合には$f(x)$の括弧は省略されます。

これらの関数$\sin$と$\cos$の意味は、数値$u$と$v$、および$u$と$w$のペアを結びつけることです。つまり、角度$AOP$が「$AP$除以$OP$」の大きさが$u$であるような角を構成し、そのとき$v$は「$PM$除以$OP$」の数値であり、$w$は「$OM$除以$OP$」の数値となります。

明らかに、弧APが円周の四分の一を超えるほど大きな数uを用いる場合、困難な状況に遭遇する可能性があります。すると点M（図26および図27参照）はOとA'の間ではなくOとAの間に位置することがあります。またPはAOA'の線の下部にあり、図26のように上部にあるわけではありません。この困難を避けるためには、座標幾何学の概念と規則を利用し、正弦および余弦の完全な定義を行います。

角度の一方の半辺OAを軸OXとして取り、その負の部分OX'を得るために軸を反対方向に延長します。もう一方の軸YOY'は它に対して垂直に引きます。点PがOから距離r離れた場合、その座標をxとyとします。これらの座標は、平面の最初の「四分象限」においてともに正です（例：図27におけるPの座標xとy）。他の四分象限では、一方または両方の座標が負になります（例：図27におけるP'の座標x'とy、P''の座標x'とy'、P'''の座標xとy'。ここでx'およびy'は負の数です）。

正角POAは弧APをrで割ったものであり、その正弦はy/r、余弦はx/rです。正角AOP'は弧ABP'をrで割ったものであり、その正弦もy/r、余弦はx'/rです。正角AOP''は弧ABA'P''をrで割ったものであり、その正弦はy'/r、余弦もx'/rです。正角AOP'''は弧ABA'B'P'''をrで割ったものであり、その正弦はy'/r、余弦はx/rです。

しかし、まだ十分に進んでいない。
円の全体的な周長に対する半径の比を、数~$u$ として選ぶとしよう。すべての円が相似であることに従ってこの比は全ての円で同じであり、数学ではその値を~$2\pi$ の形で表し、$\pi$ はギリシャ文字の~\Foreign{p} であり、ギリシャ語では ``パイ'' と呼ばれる。そして $\pi$ は無比数であることが証明され、そのためその値は分数や終端小数・循環小数で表現することはできない。その値を少数桁にすると~$3.14159$ であり、多くの用途では十分な近似値として~$\dfrac{22}{7}$ を使うことができる。数学的には $\pi$ の正確な計算は必要とする精度まで容易に行うことができ、$\sqrt{2}$ も同様である。実際には $\pi$ の値を小数点以下 $707$ 桁まで求めているが、そのような計算の詳細化はただの趣味であり、実用的または理論的な興味を持つものではない。$\pi$ の正確な決定は有名な ``円の正方形作図'' という問題の一部である。この問題のもう一つの部分は、純粋幾何学の理論的手法を使って円周と等しい長さの直線を描写することである。両方の部分は現在では不可能であることが知られており、不可解な問題は実用的または理論的な興味を持つものではなくなり、より広い概念に吸収された。

この円周率$\pi$に関する余談後、私たちは任意の値$u$に対応する角度を生成できる一般的な角度の定義に戻ります。点$Q$が点$A$から$OX$上を開始し、図27（\cf同図）で示すように円周上で正方向（図では反時計回り）に何回でも回転し、最終的に任意の点、たとえば$P$や$P'$や$P''$や$P'''$に停止します。このとき、曲線的な円周経路を全行程した距離が円の半径$r$で割った値が、任意の大きさの正角度の一般化された定義となります。点$Q$が停止する点の座標$x$，$y$を用いて、この一般化された角度の符号は$\dfrac{y}{r}$であり、余弦値は$\dfrac{x}{r}$です。これらの定義により、関数関係$v = \sin u$と$w = \cos u$がついにすべての正実数の$u$に対して定義されます。$u$が負の値の場合には、点$Q$の回転方向を反対（時計回り）にしますが、一般的な手続き方法が説明された現在、この点については詳細に論じる価値はありません。

これらの正弦と余弦の関数は、次のように定義することで、三角形に関連する問題を扱うことができます。しかし、現在では周期性というより広い概念に三角法を結びつける位置にあります。この周期性の重要性については前章で説明しました。\index{Periodicity}周期性は容易に理解できます。動点$Q$が点$A$から円周を回転し、位置$P$（図27参照）にあるとします。この位置$P$は角度$\dfrac{\text{弧長 } AP}{r}$, $2\pi + \dfrac{\text{弧長 } AP}{r}$, $4\pi + \dfrac{\text{弧長 } AP}{r}$, $6\pi + \dfrac{\text{弧長 } AP}{r}$, そして無限に続くことを表します。これらのすべての角度は同じ正弦と余弦を持ち、それは$\dfrac{y}{r}$と$\dfrac{x}{r}$です。したがって、$u$が任意の値を持つ場合、引数$u$, $2\pi + u$, $4\pi + u$, $6\pi + u$, そして無限に続くすべての角度は対応する正弦と余弦の同じ値を持ちます。言い換えれば、

\begin{alignat*}{4}
\sin u \&= \sin(2\pi + u) \&\&= \sin(4\pi + u) \&\&= \sin(6\pi + u) \&\&= \text{等}; \\
\cos u \&= \cos(2\pi + u) \&\&= \cos(4\pi + u) \&\&= \cos(6\pi + u) \&\&= \text{等}
\end{alignat*}

この事実は、$\sin u$と$\cos u$が周期$2\pi$を持つ周期関数であることを述べるものです。

The関数$y = \sin x$のグラフ（$v$と$u$を$y$と$x$に再利用しました）は、図\[fig.]{28}に示されています。$x$軸上には任意の長さを選んで$\pi$を表し、$y$軸上には任意の長さを選んで$1$を表します。正弦と余弦の数値は決して$1$を超えることはありません。周期が$2\pi$で再現されることが注意されます。このグラフは最も単純な周期関数を表しており、他のすべての周期関数はこれから構築されます。余弦は基本的には正弦と異なるものではありません。なぜなら$\cos x = \sin(x + \frac{\pi}{2})$が証明できるため、グラフ$\cos x$は図\[fig.]{28}を$\frac{\pi}{2}$を示す$OX$上の点を通じて$OY$軸を描くことで変形されます。

任意の周期$a$を持つ`正弦'関数を図\[Figure{28}\]で構築することは容易です。なぜなら
\[
y = \sin \frac{2\pi x}{a}
\]
と書くだけで、
\[
\sin \frac{2\pi (x + a)}{a} = \sin \left(\frac{2\pi x}{a} + 2\pi\right) = \sin \frac{2\pi x}{a}
\]
となるため、この新しい関数の周期は現在$a$となります。

次に、周期関数について一般的な定義を与えましょう。関数$f(x)$が周期関数であり、その周期が$a$であるとは、(i)任意の$x$に対して$f(x) = f(x + a)$が成り立つこと、および(ii) $a$より小さい任意の数$b$が存在しないことを意味します。すなわち、任意の$x$に対して$f(x) = f(x + b)$となるような$b$が存在しないということです。

第二の項は定義に入れる理由は、$\sin \dfrac{2\pi x}{a}$ が周期 $a$ の周期だけでなく、$2a$ および $3a$ も周期となるためです。これは
\[
\sin \frac{2\pi (x + 3a)}{a}
  = \sin \left(\frac{2\pi x}{a} + 6\pi\right)
  = \sin \frac{2\pi x}{a}
\]
という等式から導けます。したがって、求めたい最小の周期はこの関数の「周期」として呼びます。周期関数に関する抽象的な理論の大半と、その理論を物理学に応用する全体は、Fourierの定理によって支配されています。すなわち、周期 $a$ の周期関数 $f(x)$ が特定の条件を満たす場合（自然現象から提示される関数において実際には常に仮定される条件）、$f(x)$ は次の形で表すことができます：
\begin{multline*}
c${0} + c${1} \sin \left(\frac{2\pi x}{a} + e\\textit{{1}\right)
  + c}{2} \sin \left(\frac{4\pi x}{a} + e\\textit{{2}\right) \\
  + c}{3} \sin \left(\frac{6\pi x}{a} + e\_{3}\right) + \text{等}
\end{multline*}

この式において、$c_{0}$, $c_{1}$, $c_{2}$, $c_{3}$ などと $e_{1}$, $e_{2}$, $e_{3}$ などは特定の関数に合わせて選ばれる定数です。また、どの程度の項を選択するかについても問われます。ここで新たな難題が生じ、ある特定の場合には一定の数の項で十分であることが証明できますが、一般的には任意に近い値を得るためにより多くの項を追加する必要があります。この逐次的な近似は無限級数理論への考慮につながり、数学理論の重要な部分であり、次の章で詳しく検討します。

上の方法で周期的な関数を正弦の和として表すことは、「調和解析」と呼ばれる。例えば、海沿いの任意の地点では潮が定期的に上昇と下降する。したがって、ドーバー海峡近くには地球の回転により一日に二回の潮が生じる。一日の間に上昇と下降する潮は、イングランドチャネルから上流に向かう一連の潮と、スコットランド北端を巡り南下する一連の潮によって複雑化される。また、一部の高潮は他の高潮よりも高い：これは太陽も月と同じく潮を生成する影響があるためである。このようにして月ごとの周期やその他の周期が導入される。

風の例外的な影響を考慮に入れない。予測不能な風による影響を除いて、潮の調和解析の一般問題は、任意の瞬間における潮位に対する「周期」の生成因子の寄与を近似的に表すための項のセットを見つけ出すことである。議論$x$は、どの便利な開始点から時間として計測される。

また、ヴァイオリンの弦の振動運動も同様に調和解析が行われ、それに応じて光波と音波に対応する空気の振動も調和解析が行われる。ここでは数学物理学の基本的な過程の一つに直面している：すなわち、周期性という大きな自然現象を扱う一般的な手法である。

\chapter{第14章：Chapter XIV: Series}

数学において、初学者に最初に提示される系列の平凡さは、大きな主題である系列の部分を最も苦しめています。
系列の重要な分類として算術系列と幾何系列の2つの例が考慮されています；
これらの例は重要であり、それは重要な一般理論の最も簡単な例だからです。しかし、一般的な概念は一切明らかにされず；そして、その結果、何も示さない例が軽薄で平凡なものに削減されます。

一般的な数学的な列の概念は、順序に並べられた一連のものの集合を指すもので、つまり順列という意味です。この意味は一般的な用語の使用において正確に表現されています。

例えば、十九世紀におけるイギリス首相の系列を考えましょう。その系列は十九世紀最初にその職位に就いた人物から始まり、最後にはウィリアム・ピットからローズベルリー卿へと続きます。適切なことに、彼は最初のメンバーの生誕記を著しています。

これらの男性を他の順序で並べ替えることも考えられます；例えば、身長や体重による順序です。これらの提案された順序は首相との関連において trivia であり、自然に思い浮かぶものではありませんが、抽象的に見ればどの他の順序よりも良いものです。項の一つの順序が他より著しく重要または明らかである場合、その順序は「その項の順序」と呼ばれます。整数の順序は常に大きさによる順序を指します。しかし当然、それ以外の並べ方の無限に存在します。

考慮されるものの数が有限の場合、これらを順序で並べ替える方法の数はそれらの並び替えの数と呼ばれ、$n$ 個のものからなる集合の並び替えの数は
\[
n × (n - 1) × (n - 2) × (n - 3) × \dots × 4 × 3 × 2 × 1
\]
であり、つまり最初の $n$ 個の整数の積です。この積は数学において非常に重要であるため、特別なシンボルが用いられ、常に `$n!$` と書かれます。したがって、$2! = 2 × 1 = 2$, $3! = 3 × 2 × 1 = 6$, $4! = 4 × 3 × 2 × 1 = 24$, $5! = 5 × 4 × 3 × 2 × 1 = 120$ です。

$n$ が増えるにつれて、$n!$ の値は非常に速く増加します。したがって、$100!$ は $99!$ より百倍大きいです。

小なる$n$の値の場合、$n!$は$n$つのものの順列の方法数であることが容易に確認できる。例えば、二つのもの$a$と$b$を考えると、これらは$ab$と$ba$という2つの順序を持つことができ、$2! = 2$。

また、三つのもの$a$、$b$、および$c$を考えてみる；これらは$abc$、$acb$、$bac$、$bca$、$cab$、および$cba$という6つの順序を持つことができ、$3! = 6$である。同様に、四つのもの$a$、$b$、$c$、および$d$を順列する24通りの方法も同様である。

無限集合の場合、例えばすべての整数、またはすべての分数、またはすべての実数のような場合、順序型の理論の複雑性に直面することになる。この主題は第VI章（\ChapRef{VI}）で整数、分数、および実数の可能な順序について触れられている。順序型全体の問題は比較的新しい数学の重要な分野であるが、これについてはさらに考慮しない。現在検討しているすべての無限級数は、昇順に並べられた整数と同じ順序型であり、つまり最初の項があり、各項には左右に隣接する2つの項がある（最初の項を除き）。したがって、$m$が0でない任意の整数であれば、常に$m$番目の項がある。項の有限な数（例えば$n$個）を持つ級数は、無限級数と比較して隣接する項に関して同じ特性を持っているが、最後の項である第$n$項を除いて異なる。

数列の各項について考える重要なことは、今後「数列」という用語は先に述べた制限された意味で使用することにして、その連続した項を加算することである。

Thus if $u_{1}$, $u_{2}$, $u_{3}$, \dots, $u_{n}$, \dots are respectively the 1st, 2nd, 3rd, 4th, \dots, nth, \dots terms of a series of numbers, we form successively the series $u_{1}$, $u_{1} + u_{2}$, $u_{1} + u_{2} + u_{3}$, $u_{1} + u_{2} + u_{3} + u_{4}$, and so on; thus the sum of the 1st n terms may be written as
\[
u_{1} + u_{2} + u_{3} + \dots + u_{n}.
\]

If the series has only a finite number of terms, we come at last in this way to the sum of the whole series of terms. But if the series has an infinite number of terms, this process of successively forming the sums of the terms never terminates; and in this sense there is no such thing as the sum of an infinite series.

But why is it important successively to add the terms of a series in this way? The answer is that we are here symbolizing the fundamental mental process of approximation. This is a process which has significance far beyond the regions of mathematics. Our limited intellects cannot deal with complicated material all at once, and our method of arrangement is that of approximation. The statesman in framing his speech puts the dominating issues first and lets the details fall naturally into their subordinate places. There is, of course, the converse artistic method of preparing the imagination by the presentation of subordinate or special details, and then gradually rising to a crisis. In either way the process is one of gradual summation of effects; and this is exactly what is done by the successive summation of the terms of a series. Our ordinary method of stating numbers is such a process of gradual summation, at least in the case of large numbers. Thus $568,213$ presents itself to the mind as
\[
500,000 + 60,000 + 8,000 + 200 + 10 + 3.
\]

小数分数の場合、これはより明白にそうである。例えば $3.14159$ は以下のように表される：
\[
3 + \frac{1}{10} + \frac{4}{100} + \frac{1}{1000} + \frac{5}{10000} + \frac{9}{100000}
\]
また、$3$, $3 + \frac{1}{10}$, $3 + \frac{4}{100}$, $3 + \frac{1}{1000}$, そして $3 + \frac{1}{1000} + \frac{5}{10000}$ は完全な結果 $3.14159$ への逐次的な近似である。右から左へ $568,213$ を逆順に読むと、$3$ 単位から始めて芸術的に読み進め、最終的には $500,000$ の危機を準備する。

通常の数値乗算は級数の総和によって行われる。例えば以下の計算を行う：
\[
\begin{array}{*{6}{@{}c@{}}}
 & & &3&4&2 \\
 & & &6&5&8 \\
\cline{4-6}
\Strut
 & &2&7&3&6 \\
 &1&7&1&0&  \\
2&0&5&2& &  \\
\cline{1-6}
\Strut
2&2&5&0&3&6
\end{array}
\]

この計算では、三行を加えることで級数の系列が形成され、その最初の項は上段である。この系列は最も重要な項を最後に提示する芸術的な方法を採用しているが、これは美に対する感覚ではなく、単位の位置を確実に保つことによる利便性からであり、そのため一部の $0$ を省略できる。

しかし、無限級数の逐次的な項を加えて近似する場合、我々は何かに近づいているのか？問題はその系列が直感的な意味での「和」を持たないということである。つまり、項を加える操作は永遠に完了しないからだ。答えは我々が級数の総和の極限に近づいていることである。そして今、この「極限」とは何であるかを説明しなければならない。

級数の項の総和は、その項の数が十分大きいとき、任意の数の項の和が限界に近づくほど、その総和も限界に近づくという意味で限界に収束します。しかし、この限界に収束するという意味を正確に説明するために、現代の数学では激しい検討を受けるべきです。「十分大きい」とは何か、「ほぼ等しい」とは何か、「近づきたい」とは何か？これらの曖昧な表現は、純粋数学で唯一認められる単純抽象概念で説明する必要があります。

級数の連続した項を $u_{1}$,~$u_{2}$, $u_{3}$, $u_{4}$,~\dots, $u_{n}$, \dots とします。ここで、$u_{n}$ は級数の第 $n$ 項です。また、$s_{n}$ を最初の $n$ 個の項の和と定義します（どの $n$ でも）。したがって、

\begin{gather*}
s${1} = u${1},\quad
s${2} = u${1} + u\\textit{{2},\quad
s}{3} = u${1} + u${2} + u\\textit{{3},\quad \text{および} \\
s}{n} = u${1} + u${2} + u\\textit{{3} + \dots + u}{n}.
\end{gather*}

すると、項 $s_{1}$,~$s_{2}$, $s_{3}$,~\dots, $s_{n}$,~\dots が新しい級数を形成し、この新しい級数の形成は元の級数の総和の過程です。そして、「総和」が「限界」に収束するという言葉は、「この新しい級数の項が限界に収束する」という意味になります。現在、我々は系列の項が限界に収束することをどのように定義するかを説明しなければなりません。

現在、近似の基準（第XII章で定義された）を思い出すと、極限という概念は次のことを意味します：$l$が級数$s_{1}$, $s_{2}$, $s_{3}$, \dots, $s_{n}$, \dots の項の極限であるとは、任意の実数$k$（近似の基準として取る）に対応して、級数の某項$s_{n}$が存在し、その項以降のすべての項（すなわち$s_{n+1}$, $s_{n+2}$, など）はその近似の基準内で$l$に近づいていることを言います。もし別のより小さな基準$k^{1}$を選んだ場合、項$s_{n}$が系列の早期にある可能性があり、その性質を持つ後期の項$s_{m}$を見つけることができます。

この性質が成り立つ場合、明らかに左から右へと級数$s_{1}$, $s_{2}$, $s_{3}$, \dots, $s_{n}$, \dots の項を順に見ていくと、ある時点ですべての項が任意に指定した数より$l$に近いことがわかります。言い換えれば、$l$にどれだけ近づくかを選べば選ぶほど、その程度まで近づきます。この級数の極限の定義は第XI章で与えられた連続関数の定義と密接に関連していることがすぐに分かります。

次に元の系列$u_{1}$, $u_{2}$, $u_{3}$, \dots, $u_{n}$, \dots の項の極限を考えてみましょう。級数$s_{1}$, $s_{2}$, $s_{3}$, \dots, $s_{n}$, \dots の項の極限は元の系列の「無限和」と呼ばれます。しかし、「和」のこの用法は非常に人工的であり、有限個の項の通常の和と同様の性質を仮定する前に特別な調査が必要です。

無限級数の「和」について例を見てみましょう。再帰小数 $.1111\dots$ は、系列 $.1$, $.01$, $.001$, $.0001$, 等の「和」を表す方法です。その系列による合計で得られる系列は $s_{1} = .1$, $s_{2} = .11$, $s_{3} = .111$, $s_{4} = .1111$, 等となります。この系列の項の極限は $\frac{1}{9}$ です；これは単純な除算で容易に確認できます。具体的には
\[
\tfrac{1}{9}
  = .1 + \tfrac{1}{90}
  = .11 + \tfrac{1}{900}
  = .111 + \tfrac{1}{9000} = \text{等}
\]
よって、$\frac{3}{17}$（定義の $k$）が与えられた場合、$.1$ とすべての後続項は $\frac{1}{9}$ から $\frac{3}{17}$ よりも小さい差異を持ちます；$\frac{1}{1000}$（定義の別の $k$ の選択）が与えられた場合、$.111$ とすべての後続項は $\frac{1}{9}$ から $\frac{1}{1000}$ よりも小さい差異を持ちます；その他の任意の $k$ の選択でも同様です。

何も述べた内容が無限級数の「和」をどのように求めることかの少しのアイデアを与えるものではありません。ただ、そのような数が満足すべき条件を述べいただけです。実際、すべての場合で無限級数の「和」を見つける一般的な方法は存在しないことが根本的に問題であり、その理由は、ここで定義された「和」が必ずしも存在しないからです。無限級数に合計を持つ系列は「収束」と呼ばれ、それを持たない系列は「発散」と呼ばれます。

明らかなる発散級数の例は、$1$,~$2$, $3$,~\dots, $n$~\dots である。すなわち、位相順に並べられた整数列である。どんなに無限大までの和$l$を試みても、どんなに近似度$k$を選んでも、級数の十分多くの項を取り出すことで、その和が$l$から$kより大きい差異を持つことができる。

また、別の発散級数の例は $1$,~$1$, $1$,~等である。すなわち、各項が$1$に等しい級数である。この級数の$n$項までの和は$n$であり、$n$が増えるにつれてその和は無限大に向かって成長する。

また、別の発散級数の例は $1$,~$-1$, $1$,~$-1$, $1$,~$-1$, 等である。すなわち、項が交互に$1$と$-1$である級数である。奇数個の項の和は$1$であり、偶数個の項の和は$0$である。したがって、級数 $s_{1}$,~$s_{2}$, $s_{3}$,~\dots, $s_{n}$,~\dots の項は無限大に収束するとは限らないが、増加することもしない。

$u_{1}$,~$u_{2}$,~\dots, $u_{n}$,~\dots が無限大までの和を持つための条件として、$n$が増えるにつれて$u_{n}$が無限に減少するべきであると仮定したくなる。数学がこれほど難解な科学でないなら、そうなるだろう。不幸にもその仮定は正しくない。

例えば、級数
\[
1,\quad
\dfrac{1}{2},\quad
\dfrac{1}{3},\quad
\dfrac{1}{4},\ \dots,\quad
\dfrac{1}{n},\ \dots
\]
は収束しない。これは容易に理解できる；$(n + 1)$番目の項から$n$項を考慮する。これらの$n$項は$\dfrac{1}{n + 1}$, $\dfrac{1}{n + 2}$, $\dfrac{1}{n + 3}$, \dots, $\dfrac{1}{2n}$である：これらは$n$個で、その中で最小値は$\dfrac{1}{2n}$である。したがって、これらの和は$n$倍の$\dfrac{1}{2n}$より大きい、すなわち$\dfrac{1}{2}$より大きい。無限大までの和が存在する場合でも、その和を変更せずに隣接する項を加えることで、級数
\[
1,\quad
\dfrac{1}{2},\quad
\dfrac{1}{3} + \dfrac{1}{4},\quad
\dfrac{1}{5} + \dfrac{1}{6} + \dfrac{1}{7} + \dfrac{1}{8},\quad
\text{等.}
\]
を得る。これは先に述べたことから、第2項以降の項が級数
\[
1,\quad
\dfrac{1}{2},\quad
\dfrac{1}{2},\quad
\dfrac{1}{2},\quad
\text{等.}
\]
のそれよりも大きいという級数である。ここで最初の項以外はすべて$\dfrac{1}{2}$である。しかし、この級数は収束しない。したがって元の級数も収束しない。

この収束性に関する問題は、有限項の和から無限級数の和への推論において注意を払う必要があることを示す。有限項の最も基本的な性質は当然その和を持つことである：しかし、これは必ずしも無限級数には必要とされない基本的な性質でもある。この注意は技術用語「無限級数の和」によって誤解される可能性があることを述べるだけである。無限級数
\[
u_{1},\quad
u_{2},\quad
u_{3},\ \dots,\quad
u_{n}\Add{,}\ \dots
\]
の和は通常、
\[
u_{1} + u_{2} + u_{3} + \dots + u_{n} + \dots\Add{.}
\]

我達现在一般化级数的概念，数学如其所是的方法一样，利用变量来进行この推广。これまで、各明確な項が明確な数である定数列のみを考慮してきました。しかし、等しく、各項を変数$x$を含むある数学的表現に一般化することも可能です。したがって、次の級数を考えることができます：$1$, $x$, $x^{2}$, $x^{3}$, \dots, $x^{n}$, \dots および
\[
x,\quad \frac{x^{2}}{2},\quad \frac{x^{3}}{3},\ \dots,\quad \frac{x^{n}}{n},\ \dots.
\]

このような任意の関数の一般的概念を象徴するために、$x$ の関数 $f_{n}(x)$ を考えます。その形成には変数整数$n$が含まれています。次に、$n$ に順次 $1$, $2$, $3$, \dots を与えると、次の級数を得ることができます：
\[
f_{1}(x),\quad f_{2}(x),\quad f_{3}(x),\ \dots,\quad f_{n}(x),\ \dots.
\]
このような級数は一部の$x$の値では収束し、他の部分では発散することがあります。実際、変数$x$を含む級数がすべての$x$の値に対して収束するとは稀です。少なくとも特定の例においては、これが成り立つと仮定するのは非常に危険です。

例えば、最も単純な例として「幾何学的」\index{幾何級数|EtSeq} 級数を検討しましょう：
\[
1,\quad x,\quad x^{2},\quad x^{3},\ \dots,\quad x^{n},\ \dots.
\]
$n$項の和は次の式で与えられます：
\[
s_{n} = 1 + x + x^{2} + x^{3} + \dots + x^{n}.
\]

今両辺に$x$を乗じ，则有
\[
xs_{n} = x + x^{2} + x^{3} + x^{4} + \dots + x^{n} + x^{n+1}.
\]
次に上式から下式を引くと，
\[
s_{n}(1 - x) = s_{n} - xs_{n} = 1 - x^{n+1},
\]
よって$x$が$1$でない場合，
\[
s_{n} = \frac{1 - x^{n+1}}{1 - x}
  = \frac{1}{1 - x} - \frac{x^{n+1}}{1 - x}.
\]
今、$x$が$1$未満の絶対値を持つ場合、十分に大きな$n$に対して$\dfrac{x^{n+1}}{1 - x}$は任意に選んだ$k$よりも常に絶対値で小さい。したがって、$x$が$1$未満の場合、級数$1$, $x$, $x^{2}$, \dots, $x^{n}$, \dots は収束し、その極限は$\dfrac{1}{1 - x}$である。この主張は以下のように符号化される：
\[
\frac{1}{1 - x} = 1 + x + x^{2} + \dots + x^{n} + \dots,\quad
(-1 < x < 1).
\]
しかし、$x$が$1$を超える絶対値を持つ場合や、$1$に等しい場合は収束しない。つまり、$x$が$-1$と$+1$の間にある場合、級数は収束し、$x$が$-1$または$+1$に等しくなる場合や、$-1$から$+1$の区間にない場合は収束しない。したがって、区間$-1$から$+1$内（端点を除く）のすべての「点」で級数は収束する。

この調査の段階で別の問題が生じる。級数
\[
f_{1}(x) + f_{2}(x) + f_{3}(x) + \dots + f_{n}(x) + \dots
\]
が区間$a$から$b$内のすべての$x$に対して収束する場合、すなわち任意の$a$より大きい$x$と$b$より小さい$x$に対するすべての$x$で収束する場合を考える。また、近似値$k$を満たすために十分な項数を加算することを確認したいとする。そのような項数$n$を常に指定できるか？つまり、$n$またはそれ以上の項を使用して和を計算した場合、区間内のどの$x$の値にもかかわらず、必要な近似度を満たすことができるか？

Sometimes私たちは各$k$の値に対してこれを行うことができ、またできないこともあります。この操作ができるとき、級数は区間全体を通じて一様収束と呼ばれ、できないときは非一様収束と呼ばれます。一様収束または非一様収束が区間に及ぶかどうかは級数の性質に大きな影響を与えます。最も単純な例と数字を使ってこの事実を説明しましょう。

幾何級数を考えましょう：
\[
1 + x + x^{2} + x^{3} + \dots + x^{n} + \dots
\]

これは区間$-1$から$+1$（端点$x = ±1$は除く）において収束します。

しかし、この区間にわたって一様収束とは言えません。なぜなら、$s_{n}(x)$が$n$項の和を表すとき、$\dfrac{1}{1 - x}$に収束する$s_{n}(x)$とその差は$\dfrac{x^{n+1}}{1 - x}$であることを証明しました。ここで$n$を任意の項数として、例えば20項とし、近似の基準$k$を任意に設定します、たとえば$.001$です。すると$x$が$+1$または$-1$に十分近い値を取ることで、$\dfrac{x^{21}}{1 - x}$の数値が$.001$を超えることができます。したがって、全区間において20項では足りず、しかし一部の区間では十分な項数であることがわかります。

同じ議論は、20に代わる他の任意の数を取り、.001に代わる他の任意の近似基準を取ることでも適用できます。したがって、幾何級数 $1 + x + x^2 + x^3 + \dots + x^n + \dots$ は収束間隔 $-1$ から $+1$ の全域において非一様に収束します。しかし、$-1$ から $+1$ の間の両端内にある任意の小さな収束間隔で、幾何級数はその内部で一様に収束します。例えば、$0$ から $+\frac{1}{10}$ の間の収束間隔を取ることができます。この場合、$\frac{x^{n+1}}{1 - x}$ がこれらの境界点での値で数値的に $k$ よりも小さい任意の $n$ 値は、それらの境界点間のすべての $x$ の値に対して有効です。なぜなら、$\frac{x^{n+1}}{1 - x}$ は $x$ が数値的に小さくなるにつれて数値的に減少するからです。

例えば、$k = .001$ とします。すると、$x = \frac{1}{10}$ のとき、次のようになります：

\begin{alignat*}{3}
\&\text{for $n = 1$,}\quad \& \frac{x^{n+1}}{1 - x}
  \&= \frac{\left(\frac{1}{10}\right)^2}{1 - \frac{1}{10}}
  \&\&= \frac{1}{90} = .0111\dots, \\
\%
\&\text{for $n = 2$,}\quad \& \frac{x^{n+1}}{1 - x}
  \&= \frac{\left(\frac{1}{10}\right)^3}{1 - \frac{1}{10}}
  \&\&= \frac{1}{900} = .00111\dots, \\
\%
\&\text{for $n = 3$,}\quad \& \frac{x^{n+1}}{1 - x}
  \&= \frac{\left(\frac{1}{10}\right)^4}{1 - \frac{1}{10}}
  \&\&= \frac{1}{9000} = .000111\dots.
\end{alignat*}

したがって、全体の間隔に対して3項で十分ですが、当然、一部の部分では余分な項が必要かもしれません。

注意していただきたいのは、$-1$ から $+1$ の間全域において収束（ただし一様ではない）する幾何級数 $1 + x + x^2 + \dots + x^n + \dots$ は、各 $x$ の値に対してある数の項 $n$ を見つけることができ、どの $x$ 値でも必要な近似基準を満たすことができますが、$x$ がどちらか一方の端点 $+1$ や $-1$ に近づくにつれて、必要とする $n$ の値は大きくなります。

この重要な区別、すなわち一様収束と非一様収束の間の違いは、1847年にストークスによって発表されました---その後、ジョージ・ストークス卿となりました---そして後に独立して1850年にセイデル、ドイツの数学者によって発表されました。

非一様収束が現れる臨界点は、収束が保たれる区間全体において限界点であるとは限りません。これは幾何級数に特有の性質です。

幾何級数 $1 + x + x^{2} + \dots + x^{n} + \dots$ の場合、その収束する区間に於いて極限は単純な代数的表現 $\dfrac{1}{1 - x}$ で与えられます。しかし、これは常にそうではありません。しばしばある区間内で級数が収束することを証明できますが、その極限についてはそれ以上のことは何も知ることができないことがあります。
この方法は関数を定義する非常に良い手段です；すなわち無限の収束級数の極限としてとらえること、実際には多くの関数が、あるいはべきすべきであるように、この方法で定義されています。

したがって、初等解析において最も重要な級数は
\[
1 + x + \frac{x^{2}}{2!} + \frac{x^{3}}{3!} + \dots + \frac{x^{n}}{n!} + \dots,
\]
です。ここで $n!$ はこの章の最初で定義された意味を持ちます。この級数はすべての $x$ の値に対して絶対収束することが証明され、任意の区間内で一様収束することも示されます。したがって、それは級数として持つべき全ての快適な数学的性質を持っています。これは指数級数と呼ばれ、その無限和を $\exp x$ で表します。従って、定義により
\[
\exp x = 1 + x + \frac{x^{2}}{2!} + \frac{x^{3}}{3!} + \dots + \frac{x^{n}}{n!} + \dots
\]
と定義されます。
$\exp x$ は指数関数と呼ばれます。

初等数学の知識があれば、次のような性質を証明することは比較的容易です。
\[
(\exp x) \times (\exp y) = \exp(x + y).
\]
つまり，
\begin{multline*}
(\exp x) \times (\exp y) \\
  = 1 + (x + y) + \frac{(x + y)^2}{2!} + \frac{(x + y)^3}{3!} + \dots
  + \frac{(x + y)^n}{n!} + \dots.
\end{multline*}
これは性質~(A)で、加法定理の例です。任意の関数 [例えば $f(x)$] を定義した後には、まず $f(x + y)$ を $x$ だけの已知の関数と $y$ だけの已知の関数で表す方法を探します。これを達成できる場合、その結果は加法定理と呼ばれます。加法定理は数学解析において重要な役割を果たします。例えば、正弦の加法定理は
\[
\sin(x + y) = \sin x \cos y + \cos x \sin y,
\]
余弦の加法定理は
\[
\cos(x + y) = \cos x \cos y - \sin x \sin y.
\]
と与えられます。

実際、$\sin x$ および $\cos x$ を定義する最良の方法は、前章で説明した複雑な幾何学的手法ではなく、それぞれ次の級数の極限として
\[
x - \frac{x^3}{3!} + \frac{x^5}{5!} - \frac{x^7}{7!} + \dots,
\]
および
\[
1 - \frac{x^2}{2!} + \frac{x^4}{4!} - \frac{x^6}{6!} + \dots
\]
として定義することです。したがって、次のように定義します。
\begin{align*}
\sin x \&= x - \frac{x^3}{3!} + \frac{x^5}{5!} - \frac{x^7}{7!} + \dots, \\
\cos x \&= 1 - \frac{x^2}{2!} + \frac{x^4}{4!} - \frac{x^6}{6!} + \dots.
\end{align*}
これらの定義は幾何学的な定義と同等であり、両方の級数はすべての $x$ の値に対して収束し、任意の区間において一様に収束することが証明できます。正弦および余弦のこれらの級数は上記の指数関数の級数と一般的な類似性があります。実際、虚数の理論によって第 VII 章および第 VIII 章で説明される方法により、これらは密接に関連しています。

指数関数のグラフは、図29に示されています。この関数は$OY$軸において点$y=1$で交差します。これは$x=0$のとき、級数の各項が最初の項以外全てゼロになるためです。指数関数の重要性は、任意の瞬間における増加率がその瞬間の値に対する一定のパーセンテージである変化する物理量を表現することにあります。例えば、上記のグラフは、均一な出生率、均一な死亡率、そして移住がない状況下的任意の時間における人口の大きさを表しています。ここで$x$は任意の便利な日から計算された時間を示し、$y$は適切なスケールで人口を表します。スケールは$OA$が原点として取り入れられた日における人口を表すように設定する必要があります。しかし、ここでは「増加率」の概念に遭遇しました。これは次の章のトピックです。

指数関数に密接に関連した重要な関数は、$x$を引数として指数関数において$-x^2$に置き換えることで得られます。図30にそのグラフ$y = \exp(-x^2)$が示されています。
この曲線は「正規誤差の曲線」と呼ばれ、帽子のような形をしており、これは統計学理論にとって生身的に重要で、多くの場合における平均結果からの逸脱を教えてくれます。

重要な別の関数は、指数関数と正弦波を組み合わせることによって得られます：
\[
y = \exp(-cx) × \sin \frac{2\pi x}{p}
\]
図31にそのグラフが示されています。

そのグラフは\fig[fig.]{31}に与えられている。点$A$, $B$, $O$, $C$, $D$, $E$, $F$は等間隔$\frac{1}{2}p$で配置され、さらに前後無限に同様の点が引き伸ばされるべきである。
この関数は摩擦や「減衰」力の影響によって振動が衰えることを表している。摩擦以外であれば、これらの振動は周期的なものであり、周期$p$を持つことになる；しかし摩擦の影響により、各振動の範囲は前回よりも恒常的な割合で小さくなる。この「周期性」（正弦や余弦をシンボルに用いる必要がある）と「恒常的割合」（指数関数をシンボルに用いる必要がある）という概念の組み合わせが、この関数の形を決定しており、それは正弦関数と指数関数の積の形となっている。

\chapter{第15章：Chapter XV: The Differential Calculus}

微分積分の発明は数学の歴史上の転換点を示すものです。
科学の進歩は、アイデアの着実な累積期と、その蓄積された新たな思考材料により、ある天才が新しい方法や視点を発見し、一挙に全体をより高度なレベルに移行させる時期という、対照的な時期に分かれています。シェリーはこれらの歴史的進展の異なる時期を Avalanche の形成に例えています。

\begin{verse}
\footnotesize
\index{シェリー（引用）}\%
太陽が覚ます雪崩！その質量は嵐によって三度篩い、一層積み重なって山頂まで集まり、思想もまた積み重なり、大きな真実が解放され、諸国民が反響するまでです。\\
\dotfill
\end{verse}

この比較は一定の妥当性があります。雪崩を覚ます最後の太陽光の爆発は、その規模において他の自然力と比べて優越的であるとは限らないでしょう。科学でも同じことが言えます。最終的なアイデアで一帯の思考全体を変革する幸運な天才は、前段階での思想形成に努めた先駆者の全てを上回るわけではありません。科学の歴史を考えるとき、最終的な進歩を成し遂げた人物のみに絶叫するような賞賛と感銘を向けることは、無知で感謝の欠如です。

当該事例においては、その主題が最終的な形にまで発展するまでには長い歴史があり、それが2人の発明者によって手がけられた。その方法の痕跡すらギリシャの数学家中にも見られ、実際に主題が生み出される直前に、1601年生まれで1665年に没した優れたフランス人数学家フェルマートは、前段階のアイデアを大幅に改善し、主題の大半は彼によって作られたものとなった。また、フェルマートは、同時代の同胞ダーウェーツと共に座標幾何学の共同発明者と主張できる。

実際的には、科学界へ新しい概念を提供したのはダーウェーツであり、しかしフェルマートも独立してそれらに到達していたことは確かなことである。

ニュートンやレインツに対する私たちの敬意は、惜しむべきではない。ニュートンは数学と物理学の研究者であり、レインツは数学と哲学家であった。それぞれが自身の分野において、世界で最も天才的な人物の一人だった。共同発明は不幸な争いを引き起こした機会となった。1666年頃には、ニュートンは主題を「流動法」と呼んでその方法を使用し、『原理』の作成に従事していたが、印刷された著作では特別な代数記号は避けていた。しかし彼は直接的な声明を1693年に発表しなかった。レインツは1684年に初めての声明を公刊した。ニュートンの友人たちは、彼がニュートンによって秘密に見せられた手稿からそれを得たと主張し、またレインツもニュートンが自身から剽窃したと非難した。現在では、両者が独立して発見者としての信用を得るべきであるという疑問はほとんどない。主題は発見の準備ができていた段階に達しており、二つの能力のある人物がそれを独立して発見すること自体には驚くべきことはない。

これらの共同の発見は科学において非常に一般的である。
一般には、発見が以前の思想の流れによって導かれてからでなければならず、その頃には多くの心が重要なアイデアを追いかけていた。イギリス人に関連する発見に限定すれば、ダーウィンとウォールシーによる自然選択法の同時公表、そしてアダムスとフランスの天文学者レヴィエールによるネプチューンの同時発見が思い浮かぶ。これらの争いは、誰にその功績を与えるべきかという点でしばしば、不名誉な民族主義的精神性によって影響を受けている。しかし、数学の歴史から得られる真正に感動的な反省は、多くの時代、多くの国、そして多くの人種の人々間での思想と興味の統一性である。インド人、エジプト人、アッシリア人、ギリシャ人、アラブ人、イタリア人、フランス人、ドイツ人、イギリス人、ロシア人が科学の進歩に欠かせない貢献をしたことは確かだ。特定の国の貢献を狭き門とする嫉妬深い自慢は、より広い視野を示すものではない。

微分積分学の重要性はその本質から生じる。これは関数の増加率の系統的な考慮を行うというアイデアに基づいている。この概念は自然の研究によってすぐに私たちに提示される；速度は移動距離の増加率であり、加速は速度の増加率である。したがって、現象を全体として認識する基礎となる変化の基本的なアイデアは、変化の速さについての調査を即座に提起する。「早く」や「遅く」といった一般的な用語は、変化の速度に対する無意識的な参照から意味を得る。したがって、微分積分学は数学が自然の進行を説明するために取るべき立場の鍵に関心を持っている。

この変化の速度に関する考えは、ニュートンの頭に確かに存在しており、その主題を説明する言語においても体現されていました。
しかし、自然現象から導かれたこの視点が、主題の誕生を準備した前後の数学家中でどれほど頻繁に考えられていたかどうかは疑問です。彼らは曲線に対する接線の引出\index{Tangents}、曲線の長さの求め方、および曲線によって囲まれた面積の計算などのより抽象的な問題に関心を向けました。後二つの問題、すなわち曲線の正規化と曲線の平方は、積分学に属し\index{Integral Calculus}ますが、これは微分学と同じ一般的な主題に絡んでいます。

坐標幾何の導入
\index{Tangents}\%
は二つの視点を融合させる。例えば、(\Chg{cf.}{\cf}\ \Fig[fig.]{32}) 任意の曲線 $AQP$ とその上にある点 $P$ における接線 $PT$ を考える。座標軸を $OX$ および $OY$ とする; 曲線の方程式を $y = f(x)$ とし、$OM = x$、$PM = y$ とする。次に、曲線上の任意の動点 $Q$ の座標を $x_{1}$, $y_{1}$ とするので $y_{1} = f(x_{1})$; 接線 $TP$ 上における同じ横座標 $x_{1}$ にある点 $Q'$ の座標を $x_{1}$ および $y'$ とする。今、直線 $OX$ 上から左に右へと等速に動く点 $N$ を考える; このとき、点 $Q'$ が接線 $TP$ 上で等速に移動するにつれて、その縦座標 $y'$ も等速に増加することが明らかである。実際、直線上のどの点でも $Q'N$ に対する比率は $TN$ と同じであるから、$Q'N$ の増加率に対する $ON$ の増加率の比もまた常に同じであることがわかる。一方で、曲線上のどの点でも $QN$ の増加率（すなわち $f(x_{1})$ の増加率）は直線ではない限り点から点へと変化する。点 $Q$ が点 $P$ を通るとき、$x_{1}$ が一時的に $x$ に一致するので $f(x_{1})$ の増加率は接線 $PT$ 上の $y'$ の増加率と同じである。したがって、関数 $f(x)$ の増加率を一般化して決定する方法があるなら、曲線上の任意の点 $(x, y\Typo{,}{})$ における接線の傾きを求めることができ、それから接線を描くことができる。つまり、曲線に接線を引く問題と、関数の増加率を決定する問題は実質的に同一である。

Cone部と三角法において同様に、この二つの視点のうちより人工的なものが科学が発展した方向である。実質的に重要な側面は比較的遅くまで prominence して初めて明らかになった。これは歴史的事実に基づいた一般的な定説であり、どの科学においても最後に発見されるのはその科学が本当に何についてのものかということである。人々は数世紀にわたって盲目の探求を続け、明確な直感と困惑した好奇心のみを指南としていたが、ついには「大きな真実」が明らかになったという瞬間がある。

Let us take some special cases in order to
familiarize ourselves with the sort of ideas
which we want to make precise. A train is
in motion---how shall we determine its velocity
at some instant, let us say, at noon? We can
take an interval of five minutes which includes
noon, and measure how far the train has gone
in that period. Suppose we find it to be five
miles, we may then conclude that the train
was running at the rate of $60$ miles per hour.
But five miles is a long distance, and we
cannot be sure that just at noon the train
was moving at this pace. At noon it may
have been running $70$ miles per hour, and
afterwards the brake may have been put on.
It will be safer to work with a smaller interval,
say one minute, which includes noon, and to
measure the space traversed during that
period. But for some purposes greater
accuracy may be required, and one minute
may be too long. In practice, the necessary
inaccuracy of our measurements makes it
useless to take too small a period for measurement.
But in theory the smaller the period
the better, and we are tempted to say that
for ideal accuracy an infinitely small period
is required. The older mathematicians, in
particular Leibniz, were not only tempted,
but yielded to the temptation, and did say
it. Even now it is a useful fashion of speech,
provided that we know how to interpret it
into the language of common sense. It is
curious that, in his exposition of the foundations
of the calculus, Newton, the natural
scientist, is much more philosophical than
Leibniz, the philosopher, and on the other
hand, Leibniz provided the admirable notation
which has been so essential for the progress
of the subject.

今、純粋数学の範囲内で別の例を挙げよう。任意の値 \( x \) における関数 \( x^2 \) の増加率を求めることを考える。まだ増加率という言葉の意味を正確に定義していないが、この特定の場合に関連してその意味を理解しようと試みる。

\( x \) が \( x + h \) に増加すると、関数 \( x^2 \) も \( (x + h)^2 \) に増加する。したがって、全体の増加量は \( (x + h)^2 - x^2 \) であり、これは引数の増加 \( h \) によるものである。

従って、\( x \) から \( x + h \) の区間において、引数の単位増加に対する平均増加量は \(\frac{(x + h)^2 - x^2}{h}\) となる。

また
\[
(x + h)^2 = x^2 + 2hx + h^2,
\]
であるから
\[
\frac{(x + h)^2 - x^2}{h} = \frac{2hx + h^2}{h} = 2x + h.
\]

したがって、\( 2x + h \) は引数の単位増加に対する平均増加量であり、この平均は \( x \) から \( x + h \) の区間で取られる。

しかし、\( 2x + h \) は \( h \) の値に依存しており、区间的な大きさである。引数の値 \( x \) における増加率を得るためには、\( h \) を徐々に小さくしていく必要がある。したがって、\( h \) が「無限に減少」する極限において、\( 2x \) は引数の値 \( x \) における \( x^2 \) の増加率であると言える。

ここでは、再び無限小量という概念に直面しているように思われます。"hが無限に減少したときの極限において"といった表現を使う際にはです。ライプニッツは、いかに神秘的であっても実際に存在する無限小量やそれに応じた無限小数があると主張しました。ニュートンの言葉と概念は現代的なものでしたが、明確な説明を提供してレービンツの考えを直接的に表現しているわけではありませんでした。

十九世紀中頃にヴェイエルハ从根本理论上讲，魏尔斯特拉斯和柏林数学学派首次与える了この主题的真实解释。然而，在莱布尼茨与魏尔斯特拉斯之间，围绕これ些神秘的无无限小量，数学发现了它達而哲学家達开始解释它達，已经发展出大量的数学和哲学文献。

例えば、ベルキル希主教など一部の哲人正当地否定了整个概念的有效性，但其理由并非如これ里所述。尽管対この主题的基础进行了所有批评，毫无疑问的是，数学程序基本上是正确的。事实上，该主题本身是正确的，而解释则是错误的。正是これ种虽然解释完全错误但却可能正确的情况，使得外部批评——即限制方法追求的方法——在科学进步中往往顕得极为无益和徒劳。训练有素的观察者的直觉以及他達明らかに正在探索某些事物的好奇心，才是更安全的指导。

事实上，微分 calculus 的成功导致了大量的不良哲学产生，围绕无限小量的概念中心展开。これ种陈词滥调的遗迹至今仍可在许多微积分入门教科书的解释中找到。一个稳妥的原则是：当数学或哲学作者以朦胧的深奥性写作时，他就是在胡言乱语。

ニュートンはこの質問を次のように述べたであろう：
\index{関数の極限|EtSeq}\%
関数$f(x+h)$が$h$がゼロに近づくとき、極限として$2x + h$が$2x$になるという。私たちの仕事は、この声明を解釈し、実際にはリーブニッツの無限小量の存在を暗に前提していないことを示すのである。ニュートンの方法論を読み返すと、$h$がゼロであるとき、$2x + h$が$2x$になるという単純さを求めたくなる誘惑がある。しかし、これは許されない；なぜなら、これにより$x$から$x + h$までの区間が消滅し、平均的な増加を計算するための区間がなくなるからである。問題は、長さ$h$の区間を維持しながら、同時に$h$をゼロとして扱う方法を見つけることである。ニュートンは極限という概念によってこれを達成したのであり、現在はウィーラーントスの説明に進む。

まず最初に注意してほしいのは、$2x + h$について議論するとき、$x$の値が固定され、$h$が変化していることを考慮していた点である。つまり、$x$は「定数」として扱われた変数、またはパラメータであり、\ChapRef{IX}で説明されている通りである。したがって、実際には$2x + h$を引数$h$の関数として考えていたのである。従って、質問を一般化し、引数$h$が値ゼロに近づくとき、関数$f(h)$が極限$l$に収束することとは何を意味するかという問いを立てることができる。しかしまた、特殊な値である「ゼロ」は主題の本質には属さないことが明らかになる。さらに一般化し、引数$h$が値$a$に近づくとき、関数$f(h)$が極限$l$に収束することとは何を意味するかという問いを立てることができる。

現在、魏尔斯特拉斯の説明によれば、変数 $h$ が値 $a$ に「近づく」という概念は、我々が目指すものに対する一種の比喩的な図像を提供するかもしれませんが、根本的には完全に外れている。実際、我々が「$h$ が $a$ に近づく」という基本的概念を持ち続ける限り、無限小の概念から逃れることはできないだろう。なぜなら、これは $h$ が $a$ に無限に近いという意味を含んでいるからである。これが我々が放棄したいものだ。

従って、説明すべき表現を再定義し直すことにする。そして、関数 $f(h)$ の $a$ 处の極限が $l$ とは何を意味しているのかを問う。

関数 $f(h)$ の $a$ 处の極限は、$\ChapRef{XI}$ 章で関数の連続性について議論したときと同じ意味での $a$ の近傍の性質である。ここで「近傍」という言葉は $\ChapRef{XI}$ 章で定義されたように用いられる。関数 $f(h)$ が $a$ 处の値は $f(a)$ だが、極限はその概念から異なり、異なる可能性があり、定義されていない場合でも存在することができる。また、「近似基準」という言葉も $\ChapRef{XI}$ 章で定義された意味で用いる。実際、その章の終わりに与えられた「連続性」の定義では、事実上極限を定義している。極限の定義は以下の通りである。

変数 $x$ の値 $a$ における関数 $f(x)$ が極限 $l$ をもつとは、$a$ の近傍においてその値が任意の近似基準内で $l$ に近づくことをいう。

この定義を既に与えられた連続性の定義と比較すると、以下のようである。

変数 $x$ の値 $a$ における関数 $f(x)$ が連続とは、$a$ の近傍においてその値が任意の近似基準内で $a$ 处の値に近づくことをいう。

即刻明白、関数が点$a$で連続であるためには、(i) その関数が点$a$において極限を持ち、(ii) その極限が関数の値と等しい、という二つの条件を満たす必要がある。したがって、第XI章末尾で与えられた連続性の例は、すべて$f(a)$が点$a$における$f(x)$の極限であることを示すものであり、これは$f(x)$と考慮された$a$の値についての証明に向けたものだった。実際、関数が連続でない点での極限を考える方がより教育的である。例えば、第XI章の図20で与えられたグラフを持つ関数を考えてみよう。この関数$f(x)$は、引数が整数$0, 1, 2, 3,\ldots$以外の場合には値$1$を持ち、これらの整数値では値$0$を持っている。今$x = 3$のときの極限を考える。極限の定義においては関数の値（この場合$a=3$における$f(a)$）が除外されていることに注意せよ。しかし、$f(3)$を除くと、$x$が以下の条件を満たす任意の区間に属するとき、即ち(i) $3$が端点でない範囲内にあり、(ii) $2$や$4$には達していない場合、$f(x)$の値はすべて$1$と等しくなる。したがって、これらの値はどの基準でも$1$に近づく。よって、$3$における引数$x$の$f(x)$の極限は$1$だが、定義により$f(3) = 0$である。

これは、引数$x=3$において関数が値と極限を持つが、その値が極限とは異なる例である。第XI章末尾で考慮された関数$x^2$は、引数$x=2$における値が$4$であり、その極限も$4$であることが証明されていた。したがってここでは、値と極限が等しい関数の例がある。

Finally we come到此为止，私たちが考慮すべき本質的に重要なケースに至ります。つまり、特定の引数値において定義された値がないながら極限を持つ関数についてです。このような関数を求める必要はありません；$\dfrac{2x}{x}$は私たちの目的に適います。
今、どの数学書でも、方程式 $\dfrac{2x}{x} = 2$ を素早く且つ無批判的に記述することができます。しかし、この式には難点があります；$x$が0のとき、$\dfrac{2x}{x} = \dfrac{0}{0}$となり、$\dfrac{0}{0}$は定義されていません。したがって、関数 $\dfrac{2x}{x}$ が $x = 0$ における値は定義されていません。しかし、$x$ の他のすべての値において、関数 $\dfrac{2x}{x}$ の値は2です。つまり、$\dfrac{2x}{x}$ が $x = 0$ における極限は2であり、その値はありません。同様に、$x = a$ における $\dfrac{x^{2}}{x}$ の極限は任意の$a$について$a$であり、したがって $x = 0$ における $\dfrac{x^{2}}{x}$ の極限は0です。しかし、$x = 0$ における $\dfrac{x^{2}}{x}$ の値は $\dfrac{0}{0}$の形となり、これは定義されていません。したがって、関数 $\dfrac{x^{2}}{x}$ は極限を持ちますが、$0$において値はありません。

今や私たちは、この話題を始める問題に戻ります。任意の引数$x$における関数 $x^{2}$ の増加率とはどのような定義をするべきでしょうか。私たちの答えは、その増加率は引数$h$が0における関数 $\dfrac{(x + h)^{2} - x^{2}}{h}$ の極限であるということです。（ここで$x$は「定数」であることに注意してください）この答えがどのように機能するかを見てみましょう。私たちは
\[
\dfrac{(x + h)^{2} - x^{2}}{h}
  = \dfrac{2hx + h^{2}}{h}
  = \dfrac{h(2x + h)}{h}\Add{.}
\]

現在、変数$h$の値が$0$に近づくとき、関数$h(2x + h)/h$の極限を求める際には、$h = 0$での関数の値（存在する場合）は除外されます。しかし、$h \neq 0$のすべての値に対して$h$で除算することができます。したがって、$h = 0$における$h(2x + h)/h$の極限は、$h = 0$における$2x + h$の極限と等しいです。

任意の近似基準$k$を選んでみても、$-\frac{1}{2}k$から$+\frac{1}{2}k$までの区間に属する$h$の値に対して、$2x + h$は$2x$から$\frac{1}{2}k$未満（つまり$k$未満）の差異があります。これはどの近似基準$k$に対しても成り立ちます。したがって、$h = 0$の周辺において、$2x + h$は任意の近似基準内で$2x$に近づきます。従って、$h = 0$における$2x + h$の極限は$2x$であり、したがって$\dfrac{(x + h)^{2} - x^{2}}{h}$の$h = 0$における極限も$2x$です。従って、$x^{2}$の増加率は$x$の値において$2x$と等しいことがわかります。

より抽象的な用語として、「微分係数」または「導関数」という一般的な用語がこれまでの「増加率」に相当します。一般の定義は次の通りです：関数$f(x)$の微分係数は、その引数$h$の値が$0$に近づくときの関数$\dfrac{f(x + h) - f(x)}{h}$の極限（存在する場合）です。

どのようにして、この定義と極限の補助的定義によって、「無限に小さな数」という概念を避けてきたかというと、数学的な先駆者たちがそのような概念に悩まされた理由は、一方で平均増加を計算するために区間$x$から$x + h$を使用し、他方で最終的に$h = 0$にしたい必要があったからです。結果として、彼らは実質的には「存在する零の大きさの区間」の概念に陥るような状況に直面しました。

では、この困難をどのように避けるかというと、任意の近似の基準に対応して、ある種の性質を持つ区間を見つけることができるという概念を使うことで避けることができます。違いは、我々が「変数」の概念の重要性を理解した反面、彼らはそのような理解を持たなかったことです。

従って、数学解析の基本的な概念についての説明の終わりに至ると、我々は第II章で調査を開始したときの考え方に戻されるのです。すなわち、数学において根本的に重要な概念は「あるもの」および「どんなもの」であるということです。

\chapter{第16章：Chapter XVI: Geometry}

幾何学は、数学の残りの部分と同様に抽象的です。その中で、形状の性質や事物の相対的な位置が研究されます。しかし、観察者がどのような存在であるか、または彼がそれらを視覚的に、触覚的に、または聴覚的に認識するかどうかなど、具体的な感覚は考慮に入れません。つまり、すべて特定の感覚を無視します。さらに、パラマーソン議事堂や地球のような特定の事物も無視されます。各命題は、このような幾何学的な性質を持つ任意の事物について述べられています。もちろん、球体や円錐、三角形や正方形などの具体的な例を観察することは想像力を助けるかもしれませんが、これらの命題は本に印刷された実際の図形だけでなく、それらと同じような任意の図形にも適用されます。

したがって、幾何学は代数学と同様に「あるもの」や「いくつかのもの」という概念によって支配されています。また、同じように事物の集合間の相互関係を研究します。例えば、任意の二つの三角形 $ABC$ と $DEF$ を考慮します。

これらの三角形がすべて面で等しいためには、どの部分間にどのような関係性が必要でしょうか？これは全ての初等幾何学で最初に行われる調査の一つです。これは、二つの三角形間の特定の可能性のある関連性の研究です。答えは次の通りです：---
（a）一方の三角形の二辺とその間に含まれる角が他方の三角形の二辺とその間に含まれる角に対応するそれぞれ等しい場合、

または（b）一方の三角形の二つの角とそれらを結ぶ辺が他方の三角形の二つの角とそれらを結ぶ辺に対応するそれぞれ等しい場合、

または（c）一方の三角形の三辺が他方の三角形の三辺に対応するそれぞれ等しい場合。

この答えはさらに調査を提起します。
一つの三角形の三つの角が別の三角形の三つの角とそれぞれ等しいとき、その三角形どうしの関連性は何であるかという問いが浮上します。この更なる調査は、類似性\index{類似性}の全理論へと導きます（\cf 第13章）\PageSep{238}、これは別の種類の関連性です。

また、別の例を考えてみましょう。三角形$ABC$の内部構造について考えます。その辺と角は相互に関係しています---大きい角は大きい辺に面し、等脚三角形の底角は等しいです。三角法へ進むと、この関連性は以下の有名な形でより正確に定式化されます：
\[
\frac{\sin A}{a} = \frac{\sin B}{b} = \frac{\sin C}{c},
\]
$a^{2} = b^{2} + c^{2} - 2bc \cos A$ と二つの類似な公式があります。

さらに、三角形の角どうしのより単純な関連性もあります---つまり、その和は二直角に等しいこと；また三辺どうしの関連性もあります---つまり、任意の二辺の長さの和は第三辺の長さよりも大きいことです。

したがって、幾何学を真に研究する方法は、三角形や平行四边形、円などの興味深い単純な図形を考え、その各部分間の関連性を調査することです。幾何学者の頭には、個々の命題ではなく、互いに関係している各部分を持つ図形が存在します。代数において同様に、三角形を多角形へと一般化し、辺を楕円截線へと一般化します。または逆方向に進むと、三角形は等边、等脚、不等脚のいずれかであり、多角形は辺の数によって分類され、楕円截線は双曲線、楕円、放物線のいずれであるかによって分類されます。

preceding examples illustrate how
基本的な幾何学の概念は代数学のものと全く同じであり、ただし代数学は数値を扱い、幾何学は線分、角度、面積などの幾何学的要素を扱うだけです。この基本的な同一性が、多くの幾何学的真理が代数的形式で表現できる理由の一つとなっています。したがって、$\triangle ABC$ の角 $A$, $B$, $C$ がそれぞれ度数で表されるとき、角度間の関係は方程式
\[
A + B + C = 180°
\]
によって表されます；また、三辺 $a$, $b$, $c$ がそれぞれフィートで表されるとき、辺の関係は $a < b + c$, $b < c + a$, $c < a + b$ によって表されます。また、上記の三角関数公式も同様の事実を示すものです。したがって、変数の概念と変数間の関係性は幾何学においても代数学と同じく重要です。

しかし、幾何学と代数学の並行性はさらに推し進められることがあり、これは長さ、面積、体積、角度が測定可能であるという事実によるものです。したがって、例えば任意の長さの大きさは、その長さが特定の単位を含む回数（整数でない場合もある）によって決定されます；同様に面積や体積もまた同じです。上記の三角関数公式はこの事実の例です。しかし、これは解析幾何学において最高潮に達します。この偉大な分野はしばしば誤って「解析的楕円錐」などと呼ばれ、その一部に焦点が当たるようになります。これは人類学という大きな科学が鼻だけを研究しているように、実質的な内容から外れています。

几何学と代数学における数的手続きは本質的に同一であり、その発展において密接に絡み合っているが、空間の性質と数の性質之间必然的な根本的な区別がある---実際には空間と数のすべての重要な違いである。空間の「空間性」と数の「数性」は本質的に異なるものであり、直接的に理解しなければならない。代数学を幾何学に適用したり、幾何学を代数学に適用しても、この重要な区別を完全に消滅させる手がかりにはならない。

空間と数との間の非常に明確な違いの一つは前者が後者よりもずっと抽象的で基本的ではないように見えることです。天使の数は彼らが存在するからこそ数えることができます。一度その名前がラファエル、ガブリエル、ミカエルであることが分かって、これらの異なる名前が異なる存在を表すことを知れば、三つあることは疑いの余地なく分かります。天使の存在に関する全世界の微妙な差異も、前提が与えられた上でこの事実を変えることはできません。

しかし、彼らと空間との関係についてはまだ何も分かりません。彼らは空間に存在するのでしょうか？おそらく彼らがここにいる、またはあそこにある、どこかにある、あるいはどこでもないという言葉も無意味であるかもしれません。彼らの存在は空間上の位置とは関連がないかもしれません。従って、数はすべてのものに適用される一方で、空間はそうしなければならないとは限らないのです。

事物の局所性に対する感覚は、多くの場合またはすべての我々の感覚に伴うか、あるいはその中に含まれているかのように見受けられる。これは特定の感覚に依存しないもので、多くの感覚と共に存在するからである。しかし、それは我々が感覚によって認識するものの特別な特徴である。自分たちが他の事物に対する位置について直接認識すること自体は、音、色、味、匂いの認識と同様に、一種独特のものである。したがって最初の見方からすると、数学がその範囲において幾何学を含む限り、それは第I章で述べた抽象性の意味での抽象的ではないように思われる。

しかし、これは誤りであり、実際には空間の「空間性」は全然関係ありません。それは数学家の幾何学的直感において個人的なもので、個々に特有な方法で働きかけます。しかし、論理では空間内の事物の特定の性質や空間を形成する事物の性質のみが関与し、これらは第1章で定義された抽象の意味での完全な抽象です。これらの性質には特別な空間的把握や直感、または感覚が含まれていません。これらは数の数学的な性質と同じ基盤にあります。

したがって、幾何学的研究にとって極めて重要な直感も論理的には無関係です：それは明確に述べられた前提には入らず、推論のどの段階にも関わらないのです。これは実用的な重要性を持つ例であり、私たちの思考を刺激するためのものであります。数に対する思考を刺激するために例は必要不可欠です。

「二」や「三」と考えると、一列に並んだ線画や堆積した球体などの物理的集合が浮かび上がります。幾何学の特異性は、私たちの頭の中にある特定の一例の固定性とその圧倒的な重要性です。完全に述べられた命題の抽象論理形式は、「どのような集合の事物がこのような抽象的な性質を持つなら、それらも他の抽象的な性質を持つ」というものです。しかし、眼前に浮かぶのは空間内の点、線、面、体積の集合であり、これは避けられない例で、命題の興味を生み出す唯一の例です。

ただし、その圧倒的な重要性にもかかわらず、それはただの一例に過ぎません。

幾何学は数学的な科学として見れば、より一般的な秩序の科学の一分野である。それは次元による秩序の科学と呼ばれ得る；「次元」の限定語が付けられたのは、それが一般の秩序の一部に過ぎないという制限が直線と平面との規則的な関係を生み出すためである。

空間に対する実用的重要性は、外部物理的世界の科学的概念形成における重要な役割を理解しやすくする。一方で、我々の空間知覚は様々な感覚と結びつき合い、それらをつなげている。通常、私たちは触れる物体が見える場所と同じ場所にあると判断する；そして異常な場合でも、その物体を同じ空間で触れると考え、これが様々な感覚を結びつける実質的な事実である。したがって、空間知覚は一種の共通部分として我々の感覚に存在している。

また、空間に対する抽象的な性質は空間に関心のある多くの場合の大半を占める。言えるほどではないが、空間のどの性質に対しても抽象的な数学的声明がある。最も不利な例ではあるが、曲線が特定の形状を持つ特殊な美しさを持っているとしても、その形状にはこの形状と一致する抽象的な数学的性質があり、それ以外の性質はない。

要するに：（1）幾何学で調査される空間の性質は、事物自体の性質であり、特別な認識方法への参照なしでも研究対象となるものである；（2）空間知覚は我々の感覚と伴うものであり、おそらく全ての感覚に伴うものであるが、少なくとも多くの感覚には伴うものである。しかし、すべての事物が一つの空間や何らかの空間に存在する必要性とは思えない。

\chapter{第17章：Chapter XVII: Quantity}

前章で述べたように、長さは特定の単位長さに基づいて測定されるものであり、面積は特定の単位面積に基づいて測定されるもの、体積は特定の単位体積に基づいて測定されるものです。

長さを含むような一連のものが与えられた場合、どの一つの長さに基づいて測定されるなら、これらは同じ種類の量であると言います。したがって、長さは同じ種類の量であり、面積も同じ種類の量であり、体積も同じ種類の量です。しかし、面積は長さと同じ種類の量ではなく、また体積とは異なる種類の量です。測定という概念についてもう少し考えてみましょう（ここでは長さを例にとって）。

長さは尺で測定されます。尺を移動させることにより長さの等しさを判断します。また、三つの隣接する長さが各々1フィートずつある場合、これらは3フィートの一つの全体的な長さとなります。したがって、長さを測定するには、長さの等しさと長さの加算を決定する必要があります。何かテストが行われた場合、例えば尺を移動させた場合、その長さは等しいと言います；また何かプロセスが行われた場合、すなわち長さが連続し重複しないように確保された場合、その長さは一つの全体的な長さとして加算されたと言います。しかし、任意のテストを等しさのテストとし、任意のプロセスを加算のプロセスとはできません。加算操作の結果と等しさの判断結果は特定の前定められた条件に従う必要があります。例えば、二つの長い長さの加算により得られる長さは、二つの短い長さの加算により得られる長さよりも大きいということが一つの例です。これらの前定められた条件を正確に述べたものは量の公理と呼ばれます。その真偽について唯一問題となるのは、公理が満足された場合に普通の人々が呼ぶような量を得ることができるかどうかです。もし得られなければ、「量の公理」という名前は適切でないというだけのことです。

これらの量の公理は完全に抽象的であり、空間の数学的な性質と同じく、特定の感覚形態を前提としません。量という概念に関連するアイデアは、線や面積、体積のような連続体を明確な部分に分割するための手段です。これらの部分が数えられるので、数を使って連続的な全体の正確な性質を決定することができます。

時間の流れと事件の順序に対する当時の認識は、これらの量の概念の適用の主要な例です。時間の測定（周期性について考慮したときには述べた通り）は、類似のイベントの繰り返しによって行われます——連続する等長のろうそくの1インチずつの燃焼、地球が固定星に対して回転するとき、時計の針が回転するときなどはそのような繰り返しがあります。これらの種類のイベントは長さに対する尺の役割を果たします。各再発時にこれらのいずれかの種類のイベントが正確に等しい時間ではないと仮定する必要はありません。必要なのは、ある種類の2つの例の相対的な持続時間を表現できる規則があることだけです。例えば、もし私たちが望むのであれば、地球の回転速度が減少していると仮定し、各日は前日より若干の秒分の長さであるとします。このような規則はどの日の長さを他のどの日の長さと比較するのに使用できます。しかし、重要なのは、例えば連続する日のような一連の繰り返しが基準となる一連として取り扱われるべきであり、その系列内の各イベントが等しい持続時間ではない場合でも、各日に与えられる持続時間を他のどの日の持続時間で規制する規則が述べられているべきであるということです。

どのような規則が求められるかと言えば、まず第一に、その規則は普通の感覚が等しい期間だと判断する出来事に対してほぼ等しい期間を指定することにつながるべきである。日が激しく異なる長さになり、明らかに類似した操作の速度がその表面上の微細な差異に対する比に全く比例しないようにされるような規則は不適切である。従って最初の要求条件は、一般的な感覚と一致することである。しかし、これだけでは絶対的に規則を決定するためには十分でない。普通の感覚は粗雑な観察者であり、非常に容易に満足するからである。次の要求条件は、規則の微細な調整が自然法則の最も単純な表現を可能にするようにすることである。例えば、天文学者は地球の自転が遅くなることを教えており、そのため各日は不思議なほど小さな秒数の部分で長くなるという。その主張の唯一の理由（より詳しく論じられる周期性に関する議論の中で述べられている）は、それなしではニュートンの運動法則を放棄しなければならないことである。運動法則を単純にするために時間の測定を変更する必要がある。これは完全に理解されている限り正当な手順である。

以上の文章を忠実に日本語へ翻訳しました：

空間の抽象的な数学的性質について述べたことは、時間の数学的性質にも適切な文言の変更を加えて適用できます。時間の流れという感覚は私たちのすべての感覚と認識とを伴いますし、時間に関する私たちが特に興味を持つほとんどのものは、それを抽象的な数学的性質に並行させることができます。

逆に、一日の長さを決定する規則について述べた二つの要件も、ヤード棒の長さを決定する規則にも適用されます。すなわち、ヤード棒は移動しても同じ長さを保つように見えます。したがって、どのような規則でも、微細な変化を除いて恒常的な長さを維持していることを明らかにしなければなりません。また、二つの要件のもう一つは、最も単純な形で自然法則を表すための微細な変化に関する明確な規則を述べることです。例えば、第二の要件に基づいて、ヤード棒は温度の変化に応じて物質によって異なる拡大収縮します。

私たちの感覚が場所と時間との認識とを伴い、線、面積、体積、そして時間はそれぞれ量であるという事実以外では、数論は宇宙の法則を探求する上で非常に下位的な役割しか果たしません。現状では、物理学は主に数、量、空間、時間を基盤とする主要な概念に依存しています。それらに関連した数学的科学はすべての数学を形成するものではありませんが、現在存在する数学物理の基礎となっています。

付録：注記

\Note{A} (\Pageref{60}).---これらの式を読む際には、括弧が数学的な記号で使用され、その中にある演算が最初に行われるという意味であることに注意しなければなりません。つまり $(1 + 3) + 2$ はまず $3$ を $1$ に加え、その後 $2$ を結果に加えることを指示します。また $1 + (3 + 2)$ はまず $2$ を $3$ に加え、その結果を $1$ に加えることを指示します。また、式 \Eq{(5)} の例として以下の数値計算があります：
\[
2 × (3 + 4) = (2 × 3) + (2 × 4).
\]
まず括弧内の演算を行い、
\[
2 × 7 = 6 + 8
\]
となり、これは明らかに成り立ちます。

\Note{B} (\Pageref{136}).---基本的な比 $\dfrac{SP}{PN}$ はこの曲線の離心率と呼ばれています。曲線の形状（尺度や大きさとは異なる）はその離心率の値によって決定されます。したがって、楕円全体や双曲線全体をそれぞれどちらの場合も一つの特定の形状として考えるのは誤りです。異なる離心率を持つ楕円は異なる形状を持ち、その大きさは主軸の長さに依存します。離心率が小さい楕円は非常に近似して円になりますし、離心率が1より僅かに小さい楕円は長い扁平な卵形となります。すべての放物線は同じ離心率を持ち、したがって同じ形状を持ちますが、描く尺度は異なる場合があります。
\PageSep{251}

もしすべての項が正である級数が収束する場合、ある一定の規則に従っていくつかの項を正にして他の項を負にする改変された級数も収束します。このように得られる系列のそれぞれ、元の系列を含むものは「絶対収束」と呼ばれます。しかし、部分的に正で部分的に負である項を持つ級数が収束する場合でも、そのすべての項が正である改変された級数は発散することがあります。例えば、
\[
1 - \frac{1}{2} + \frac{1}{3} - \frac{1}{4} + \text{等}
\]
は収束しますが、私たちは既に証明したように
\[
1 + \frac{1}{2} + \frac{1}{3} + \frac{1}{4} + \text{等}
\]
は発散します。このような収束するが絶対収束でない級数は、絶対収束する級数よりも扱いが難しいです。

付録[数学の研究に関する補足説明]{参考文献}

初学者がこの科学を学ぶ際に直面する困難は、基本教科書に蓄積された大量の技術的な詳細によって重要なものから目を逸らすことによるものです。

初めの研究対象は、前提となる算術知識以外では、基本的な幾何学と基本的な代数学が含まれます。両方の科目の講座は短く、必要な概念だけを与えます；代数はグラフィカルに学習することで、実践的に基礎的な座標幾何学のアイデアも吸収します。次に扱うべき科目は基本的な三角法と直線および円の座標幾何学です。後者の科目は短く、実際には代数に合流しています。その後、学生は放物線などの円錐曲線への準備が整います；これは非常に短い幾何学的円錐曲線の講座と長い解析的な円錐曲線の講座です。しかし、これらのすべての講座においては、必要以上の詳細を頭に押し込めるのではなく、基本的なアイデアの説明のために必要な詳細だけを与えるべきです。

微分積分学が残されており、その後は同様のシステムで積分積分学も攻めます。良い教師は既に代数学と座標幾何学の課程の中で特別なケースを考慮することでこれらの概念を示唆してきました。3次元幾何に関する短い本も読むべきです。

この基本的な数学科目は、一部の専門職キャリアには十分です。また、その科目自体の興味のために学ぶ人にとっては必要な前段階でもあります。彼は今やより広範な課程に着手することができますが、しかし一挙に全てをマスターするのを期待すべきではありません。科学は非常に大きな規模になっており、おそらく生き残っている数学者の中で誰もこれを行うことを主張することはできません。

本稿の前段階の講座を終え、真剣な論文を読むには、次のようなものがあげられます：クリモーナの『純粋幾何学』（オックスフォード・クラレンデンプレスによる英語版）、ホブソンの『三角法論』、クリスタルの『代数学論』（2巻），サロモンの『楕円軌跡論』、ランブの『微分法論』、あるいは微分方程式に関する本。学生はおそらくすべてのこれらの科目に等しい注意を向けず、自分の興味に従って一つまたは複数のものを学ぶでしょう。その結果、より高度な著作を選択し、主題の高い部分に没頭する準備ができることになります。分析に興味があるなら、現在では複素変数関数論の基礎的な論文をマスターすべきです；幾何学に特化したいなら、3次元解析幾何学の標準的な著作に進むべきです。しかし、彼の学習キャリアにおけるこの段階では、特にこのメモからの助言は必要ありません。

私は故意に基本的な著作を挙げませんでした。それらは非常に多く、それぞれ異なる優れた点を持っていますが、どの本も他のすべてのものよりも顕著な優越性があるわけではありませんので、特別な名前で紹介する必要はありませんでした。

アベル 156

 абсцисса 95

絶対収束 251

幾何学の抽象性|連続 242

抽象性（定義） 9, 13

アダムス 220

加法則 212

アヘムス 71

アレクサンドル大王 128, 129

代数学の基本法則 60

アンペール 34

解析的楕円軌跡 240

ペルガモンのオポリオン 131, 134

近似|連続 197

アラビア数字表記|連続 58

アキメデス|連続 37

関数の引数 146

アリストテレス 30, 42, 128

天文学者 137, 173, 174

軸 125

量の公理|連続 246

ベーカン 156

ボール、W. W. R. 58

ベーコンフィールド伯爵 41

ベルキルシー主教 226

バッシャラ 58

カントール、ゲオルク 79

円 120, 130

円|連続 180

円柱 143

マクスウェル卿 34, 35

コロンブス 122

コンパクトな級数 76

複素量 109

放物線|連続 128

定数 69, 117

連続関数@連続関数 150

連続関数@連続関数（定義） 162

絶対収束 251

収束|連続 203

座標幾何学|連続 112

座標 95

コペルニクス 45, 137

コサイン|連続 182

コウ 33

交比 140

ダーウィン 138, 220

導関数 234

 Descartes 95, 113, 116, 122, 218

微分法|連続 217

微分係数 234

直線 135

不連続関数|連続 150

距離 30

収束しない|連続 203
**ページ254**

動力学的説明 13, 14

動力学的説明|連続 47

動力学 30

動力学|連続 43

偏心率 250

電流 33

電気|連続 32

電磁気学|連続 31

楕円 45, 120

楕円|連続 130

ユークリッド 114

指数級数|連続 211

ファラデイ 34

フェルマー 218

流数 219

焦点 120, 135

力 30

代数的形式@代数的形式|連続 66

代数的形式 82, 117

フーリエの定理 191

分数|連続 71

フレンチル 32, 122

関数|連続 144

ガリレオ@ガリレオ|連続 42

ガリレオ 30, 122

ガヴェイ 33

数学における普遍性 82

幾何級数|連続 206

幾何学 36

幾何学|連続 236

ギルバート博士 32

グラフ|連続 148

重力 29, 139

ハーレイ 139

調和解析 192

ハリオット、トマス 66

ヘルツ 35

ヒエロ 38

ヒッパルコス 173

双曲線|連続 131

虚数|連続 87

虚量 109

無比率|連続 72

極小量|連続 226

積分法 222

区間|連続 158

ケプラー 45, 46, 137, 138

ケプラーの法則 138

ラプタ 10

運動の第一法則 43
**ページ255**

近傍|連続 159

ニュートン 10, 16, 30, 34, 37, 38, 43, 46, 139

ニュートン@ニュートン|連続 218

非一様収束|EtSeq 208

通常誤差、曲線 214

オーレステッド@Öersted 34

順序|EtSeq 194

順序、種類@Order, Type of|EtSeq 75

順序、種類 196

順序対|EtSeq 93

縦座標 95

原点 95, 126

パッパス 135, 136

放物線|EtSeq 131

平行四辺形の法則@Parallelogram Law|EtSeq 51

平行四辺形の法則 99, 126

パラメータ 69, 117

束 140

周期 170

周期@Period|EtSeq 189

周期性@Periodicity|EtSeq 164

周期性 188, 216

ピット、ウィリアム 194

ピサロ 122

プラUTERSキ 37

正数と負数|EtSeq 83

射影幾何学 139

ポタメイウス 137, 173

ピタゴラス 18

量|EtSeq 245

関数の増加率|EtSeq 220

比|EtSeq 72

実数|EtSeq 73

長方形 57

変数間の関係|EtSeq 18

共振 170, 171

ローズバリー、ロード 194

地図のスケール 178

シーデル 210

級数|EtSeq 74, 194

シェリリア（引用） 217

相似性@Similarity|EtSeq 177

相似性 237

正弦|EtSeq 182

密度 41

円の正方形化 187

近似の基準|EtSeq 159, 201, 229

手順@Steps|EtSeq 79

手順 96

シュタフィル 85

スコット Baron ジョージ 210

無限大までの和|EtSeq 201

測量|EtSeq 176

スウィフト 10

接線 221, 222

テイラーの定理 156, 157

時間|EtSeq 166, 247

輸送、ベクトルの|EtSeq 54

三角形@Triangle|EtSeq 176

三角形 237

三角測量 177

三角法|EtSeq 173

一様収束|EtSeq 208

未知数 17, 23
\% ** Page 256

関数の値 146

変数、そのもの 18, 24, 49, 82, 234, 239

変数関数 147

ベクトル@Vectors|EtSeq 51

ベクトル 85, 96

頂点 134

ボルタ 33

ウォールケイ 220

ヴィエトリス 156, 226, 228

ゼロ@Zero|EtSeq 63

ゼロ 103
\%[** TN: End of index]
\fi

\% ハズル、ウォルトン \& バイニー、Ltd., ロンドンおよびアイルズバリーで印刷されました。
\%[** TN: OCRによる書誌情報の原始出力が続きます]

無謬の確信をもって、この図書館シリーズは国家の教養に関心を持つ大衆の一部に対する真摋な文献の先駆者であると断言できる。デイリー・メイル。

この図書館シリーズが実質的な成功を収めやすいのは、単に無知を知識で満たそうとするだけでなく、その巻物はしばしば不適切に部分的に満たされた渇望をさらに消化不良の結果を伴うものから救う役割も果たしているためである。供給される食料は従って頻繁に薬的効果を持つとともに栄養価がある；これは確かに時代が求めていることである。マンチェスター・ガーディアン。

各巻は優れた脳の話す力と、専門家が自分の分野について語るときのような容易さと興味深い自由さで三時間の交通を代表するものである。…このシリーズは実質的な社会サービスを提供することを約束している。タイムズ。

私たちは、より高価な

id prol
シリーズに代わって賞与として授与されたら、それほど適切なものはないと考える。もし出版者がこのシリーズの賞与日には生徒たちがより多くの利益を得られるようにするための方法を探しているなら、彼らは次のように言えるべきである：「小さな部屋にも無限の財宝がある」。アイルランド教育雑誌。

この計画は始まりから成功を収めやすいのは、真剣な読者たちが満たすべき欲求に合致していたためだが、より広範で持続的な成功は、その大部分で Taste を創り出し、確実にそれを洗練させたことによるものであると断言できる。デイリー・クリスチョン。

ここには世界の学識が詰まっているし、部屋もあまり狭くないだろう！デイリーテレグラフ。

このシリーズを布製で

256ページ

2/6ポンドで革製で

提供する。

歴史と地理

3. フランス革命

ヒラリー・ベロック博士著（地図付き）「作者の性質に色を付けたものである」とデイリーニュース。

4. 战争与和平的历史

G.H. フェリス著。ジャmes ブライア卿（右栄典）は記している。「私はこの本を非常に興味深く、楽しく読んだ。あなたが多くの事実と見解を小さな巻きに収めることの技巧に感心した。」

8. 北極探検

W.S. ブルース博士（ロイヤル・スコットランド工学会会員）著、「スコティア」探検隊のリーダー。「非常に新鮮な文体で、興味深い記録である。」タイムズ。「魅力的な本である。」ポータスモス・タイムズ。

12. アフリカの開拓

H.H. ジョンストン卿（大英帝国勲章、大英帝国功労勲章、博士号、英国皇家学会会員）著。「この優れた著作により、ホーム・ユニバーシティ図書館は豊かさを増した。」デイリーマイル。

13. 中世ヨーロッパ

H.W.C. デavis先生（博士号）著。「デavis先生は主題について簡潔に記述するための完全な専門家であるという事実がまた一つ証明された。」マンチェスター・ガーディアン。

14. パピアと近代史（1303-1870）

ウィリアム・バリー博士（ドクター号）著。「バリー博士は広範な知識を持ち、芸術家の選択力を持っている。」マンチェスター・ガーディアン。

23. 1885-1911年の歴史

G.P. グー奇先生（博士号）著。「グー奇さんは自分の物語に生命を吹き込み、最近の出来事の骨と肉を我々に与える方法を見つけ出した。」オブザーバー。

25. 中国文明

H.A. ギレス博士（ラテド号）、ケンブリッジ大学の中国語教授著。「ギレス博士は大量の事実の中でも退屈することなく、読者に娯楽を提供する幽霊話や街頭冒険話を常に用意している。」スペクターター。

29. 歴史の Dawn

J.L. ミリース先生（ファースト・スコットランド考古学会会員）、オックスフォードの古代歴代学教授著。「その中で一つのページも退屈なものがない。」マンチェスター・ガーディアン。

33. 英国史

A.F. ポラード教授（博士号）著、年代順テーブル付き。「ポラード教授は一瞬でも早くイギリス歴史の権威的な著作の一冊として位置づけられた。」オブザーバー。

34. 加拿大

アーサー・G・ブラッドリー作。カナダを熟知している人物は誰よりもマスターズ・エイジェント・A.G.ブラッドリーでしょう。
デイリーチロニカル。この本は、カナダについて鮮やかで正確な知識を得たい人に対して即座に魅力を発揮します。カナディアン・ゲゼイル。

37. ペーパーズ 6* インドの問題

ティム・W・ホルネス卿作。インド事務所の収入、統計、商務省大臣。今日の新聞読者にとって必要な本で、驚くべき内容の豊かさがあります。
パールマール・ゲゼイル。

42. ローマ

ヴィンセント・W・フォールスハース作。ローマ人の性格とその世界への貢献についての優れた概要です。
スペクタター。この作家から期待する明瞭さと魅力がありますと言います。マンチェスター・ガゼイル。

48. アメリカ合衆国 Civil War

F.L.パックスン作。ウィスコンシン大学アメリカ歴史教授。地図付き。「興奮を誘う研究」です。
ガーディアン。

51. ブリテンの戦争

ヒラール・ベルロック作。MA。過去の大規模な戦闘がブリテン本土でどのように、そしてどこで行われたかについての記述。道路や物理的条件が島国の戦略を決定し、城塞、要塞都市など。

55. マスターマリンナーズ

J.R.スピアース作。海のロマンス、大航海時代の探検航海、海軍戦闘、水手の英雄性、船の発展から古代まで。

準備中

古代ギリシャ。ギルバート・ムーリー教授作。D.Litt., LL.D., F.B.A
古代エジプト。F. L. グリフィス作。MA。
古代東方。D.G.ホーカース作。MA., F.B.A.
ヨーロッパの短い歴史。ハーバート・フィシャー教授作。MA., F.B.A.
前史的ブリテン。ロバート・ムーン作。MA., M.D., LL.D.
東ローマ帝国。ノーマン・H. ベインズ作。
改革運動。LL.D. リンドスアイ長官作。
ナポレオン。ハーバート・フィシャー教授作。MA., F.B.A.
ロシアの短い歴史。ミルユコフ教授作。
現代トルコ。D.G.ホーカース作。MA.
今日のフランス。アルベルト・トーマス作。
今日のドイツ。チャールズ・タワー作。
海軍と海上力。デイビッド・ハンナ作。
スコットランドの歴史。R.S.レイ特教授作。
南アメリカ。W.R.シェパード教授作。
ロンドン。Sir ラーレンス・ゴム作。F.S.A.

歴史とスペインの文学。

2.シェイクスピア

ジョン・メイスフェルド著。『この本は素晴らしいです。最近数年間、半打以上もの学術的なシェイクスピアに関する書籍が出版されましたがあれほど賢明なものは一つもありません』。マンチェスター・ガーディアン。

27.現代の英文学

G.H.メイアール著。『一貫性のある新しい本です』。オルスター。

35.フランス文学の landmarks

G.L.ストラシー著。『ストラシーは勇敢で成功したことを祝いたいと思います。彼が提供した内容を考えれば、250ページほどの小さな文章の中でより良いフランス文学の説明を想像するのは難しいです』。タイムズ。

39-建築学

W.R.レタビー教授著。（40以上の図版付き）『一般的な建築ガイドブックは通常価値がありませんが、この本は歓迎すべき例外です』。ビルディング・ニュース。「快適に読める内容です」。クリスチャン・ワールド。

43.中世の英文学

W.P.ケル教授著。『ケル先生は長い間、我々の英語学界で最も堅実な学者の一人として証明してきましたし、彼は未訓練された一般読者に英語中世文学の概要を提供するのに適した人物です。彼の知識と趣味は疑いの余地がなく、その文体は効果的で簡潔ですが決して乾燥したものではありません』。アスメニュム。

45.英語の言語

L.ペアスール・スミス著。『大きな英語言語の流れを形成した異なる流れについての完全に魅力的な研究です』。デイリーニュース。

52.アメリカの大作家たち

J.E.エクスキン教授とW.P.トレント教授共著。二名の主要な専門家による人気概説。

準備中

古代の芸術と儀式。ジェーン・ハリソン博士（ドクターリテラトゥア）著。

ギリシャ文学。ギルバート・ムーリー教授（ドクターリテラトゥア）著。

ラテン文学。J.S.フィリモア教授著。

チャUCERとその時代。ジェイムズ・エドワード・ハドウ小姐著。

人文主義。R.A.タイラー夫人著。

以下の原著英文を忠実に日本語へ翻訳しました：

イタリア Renaissance の芸術。ロジャー・エー・フライ、M.A.
絵画の芸術。サー・フレデリック・ウェドモア。
ジョナサン・ジョンソンとその circle。ジョン・ベイルリー、M.A.
維多利亚時代。G.K.チェストerton。
英語作文。プロフェッサー・ウォム・T・ブロストー。
ロシアの偉大な作家たち。C.T.H.ライツ、LL.D.
ドイツの文学と文学。プロフェッサー・J.G.ロバートソン、M.A., Ph.D.
スカンジナビアの歴史と文学。T.C.スノウ、M.A.

科学

7. 現代地理学
マリオン・ニュービギン博士。『地理学：かつては退屈でつまらない科目だったが、マリオン・ニュービギン博士はその乾き切った骨にロマンチックな興味を注入し、科学の妖精書として地理学を見直しています』。デイリー・テレグラフ。

9. 植物の進化
D.H.スコット博士、M.A., F.R.S.（ジョドレル研究所元名誉所長）。『本書が提供する情報は、第一級の知識がもたらすものほど信頼性が高いです。…スコット博士の率直で親しみやすい文体は、しばしば退屈で時には不快なテーマを魅力的かつ容易に説明します』。ガーデンナーズ・クリシナル。

17. 健康と疾病
エディンバラ地方自治局医師部長のW.レスリー・マッケイナック博士。
公衆衛生管理の科学は、マッケイナック博士ほど優れた魅力的な解説者はありませんでした。彼は問題に対する深い理解に加えて、洞察力のある文体としばしば退屈で時には不快なテーマを引き立てる方法を持っています』。エコノミスト。

18. 数学入門
A.N.ホワイトヘッド博士、Sc.D., F.R.S.（図表付き）。『ホワイトヘッド博士は、彼が非常に適格であると認められる任務を著しく成功裏に果たしました。彼は科学の基礎に関する当世の大いなる権威であり、読者に対してその目的を提示するためには必要な広い視野を持っています。彼の解説は明瞭で印象的です』。ウェストミンスター・ゲゼルシャフト。

19. 动物界

教授F.W. ガムبل博士、D.Sc., F.R.S. 著。オリバー・ロッジ卿による序文。多数の図版付き。「動物（および植物）の生活の魅力的で教育的な要約です。…最も魅惑的で示唆的な概観です。」モーニングポスト。

20. 演化論

教授J. アーサー・ Thomsonと教授パトリック・ゲデス共著。「多色でロマンチックな風景画、我々が知る他のどの本よりも、世界の発展の理不尽な視覚を開示します。」ベルファストニュースレター。

22. 犯罪と精神障害

ドクターセバスチャン・A. メシヤール博士、F.R.C.P., F.R.C.S. 著、「精神病学的法医学の最高位に立つ人物から得られる貴重な情報が提供されています。」アズイル NCJVS。

28. 心霊研究

そして彼が言及する思考読み取り、催眠術、心電感応、水晶視覚、スピリチュアルズム、占いなどについては、「熱心に読まれるでしょう。」ダンディー。

31. 天文学

アーサー・R. ヒンクス博士、M.A., カムブリッジ天文台長。「思想は革新的で、内容は eclectic で、対処は批判的です。…より良い小さな本はありません。」スクールワールド。

32. 科学入門

教授J. アーサー・ Thomson、M.A., 阿伯丁大学自然史名誉教授。「トムソン教授の魅力的な文学スタイルは有名で、ここでは科学の方法とその哲学、芸術、宗教、実践的生活との関係について新鮮かつ容易に語っています。」ア伯丁ジャーナル。

36.

ホー・N. ディクソン博士、D.Sc. オックスフォード, M.A., F.R.S.E.、皇家气象学会会长；雷丁大学地理学教授。（図表付き）「著者は、大気の動きとより安定した風の原因を非常に明瞭で快適な方法で提示しました。」マンチェスター・ガーディアン。

41. 人間科学

R.R. マレット博士、M.A., 著。「絶対に完璧な手引きです。魅力的かつ人間味があり、それは…」

44. 生理学の原理

ジェームズ・G・マッケンドリック教授、M.D.「これは、すべての人にとって重要であるが、技術的な説明には適さない題材についての、魅力的かつ広範な扱いです。…この小さな本は単なる知識の宝庫ではなく、その中の毎ページに創造的な想像力の烙印が押されています。」ゴアーズバリー・ヘラルド。

ローランド・R・マレット、M.A.、オックスフォード大学社会人類学講師。「絶対的に完璧なハンドブックで、子供でも理解できるほど明瞭であり、興味深く人間的で、小説を破り尽くすほど優れています。」モーニングリーダー。

46. 物質とエネルギー

F.S.ソディ、M.A., F.R.S.「物理科学の大切な事実についての、魅力的かつ教育的な解説で、近年の知識の進歩は驚くべきものでした。」ブックセラー。

49. 心理学、行動の研究

W.マクドゥガル教授、F.R.S., M.B.「非技術的な手法による難解な科学の扱いの良い例で、推測よりも断定的ではありません。より深い研究への興味を掻き立てるでしょう。」クリスチャン・ワールド。

53. 地球の形成

J.W.グレゴリー教授、F.R.S.（図版38枚付）グラスゴーの地質学教授は、地球の起源、表面と構造の形成と変化、地質史、生命の最初の出現、そしてその影響について記述しています。

57. 人間の身体

A.ケイチ、M.D., LL.D., 博物館保存者、ハンター教授、ロイヤル外科学会。「解剖室での作業が描写され、体の発達、異形と奇形、若年期と老年期の変化、性差は増加しているのか減少しているのか、人種的特徴、身体的な特徴が精神的性格の指標であるか、退化と再生、そして人類の系譜と古代についても扱われています。」（図版付き）

58. 電気

ギスベルト・カップ博士、D.Eng., M.I.E.E., M.I.C.E.、バーミンガム大学電気工学教授。摩擦電気と接触電気、電位、機械的手段による electrification、電流、電流の動力学、交流電流、電力の配分などについて論じています。

準備中

化学。R. メルドラー博士、F.R.S.

鉱物世界。ティム・H・ホランド卿、K.C.I.E., D.Sc.

植物学。J.B. ファーマー博士、F.R.S.

神経系。D.F. ハリス博士、M.D., D.Sc.

性についての研究。J.A. シュモーン教授とパトリック・ゲデス教授。

ヨーロッパの成長。グレンビル・コール教授

哲学と「宗教」

15. イスラム教

D.S. マルゴリウス博士、M.A., D.Litt. 「この豪華な一シリング分の知恵。... 視野が広く、ユーモアがあり、責任感のある教授による論説」デイリー・メイル。

40. 哲学の問題

バートランド・ルイス博士、F.R.S. 「街の人々はすぐにこの恩恵である本を認識するでしょう。... 全体的に明瞭で技術的なものではありません」クリスチャン・ワールド。

47. 華厳教

W.B. セルビー准教授、M.A. 「歴史的部分はその明瞭さと適切なバランスで輝いており、非妥協派の非妥協主義の現在の状況に関する後期章では、セルビー博士が健全で中庸的な見解を説く理想的な解釈者であることが証明されます」クリスチャン・ワールド。

54. 倫理学

G.E. モア教授、M.A. カムブリッジ大学道徳科学講師。利己主義、道徳判断の客観性、善と悪の試金石、自由意志、内在的価値について論じています。

56. 新約聖書の形成

プロフェッサーB.W.バコン、LL.D., D.D.著。現代の批判的研究による新約聖書の起源に関する公式なまとめ、「意識的な啓発がまだ全盛期だった形成期における二大思潮の混ざり合い」を示し、「パウロ派と「使徒的」、ギリシャキリスト教のイエスの福音、そしてユダヤキリスト教のイエスの福音、精神の福音と権威の福音が交差している」という時期。

伝道：その発展と進歩
マーガレットクリアighton著。改革後の現代的伝道の始まりとその成長を追跡し、現在の活動についての記述を含む。

準備中

旧約聖書
プロフェッサージョージムーれ、D.D., LL.D.著

旧約聖書と新約聖書間
R.H.チャールズ、D.D.著

比較宗教
プロフェッサーJ.エスタリンカーペンター、D.Litt.著

思想の自由の歴史
プロフェッサーJ.B.バリー、LL.D.著

哲学の歴史
Clemont Webb、M.A.著

社会科学

国会
議会の歴史、憲法と実践。Sir Courtenay P.イルバート、K.C.B., K.C.S.I., 議院書記官著。「Bagehotの「憲法」以来、議会の歴史と実践に関する最高の本である」という評価。ヤークシャーポスト。

証券取引所
F.W.ヒースト、『エコノミスト』編集者著。「財政に疎い心には啓示となる。…この本はBagehotの「ロムブール通り」ほど明快で力強いものはない」という評価。モーニングリーダー。

アイルランドの民族性
Mrs J.R.グリーン著。「学術的かつ燃えるような内容。もっとも適切な書籍である。…ゲール精神の正当な活力を強調する壮麗な証明」という評価。フリーマンズジャーナル。

社会主義運動

RAMSAY MACDONALD, M.T. "適切に説明のために作られました。" タイムズ新聞。"マクドナルド氏は非常に明晰な解説者です。…この本は、この国における社会主義の動向に関する誤った印象を払拭するのに大いに役立ちます。" ナショナル。

i. コンサーバatism

ヒュー・シルト卿（M.A., M.P.） "一世代に一度しか現れないような偉大な小冊子の一つです。…" モーニングポスト。

16. 貧困の科学

J.A.ホュソン（M.A.） "生きてる経済学者の中でも、J.A.ホブソンは一 blackColorを占めています。…この教科書は全く傑出しています。新鮮で合理的で、啓発的です。" ナショナル。

21. 自由主義

ロンドン大学社会学教授のL.T.ホブハウス（M.A.） "希少な質の本です。…この本には、基本原理から導き出される主張の迅速かつ巧みな要約が大部分を占めています。それらは我々にとって絶賛すべきものです。" ウェストミンスター・ガゼッタ。

24. 産業の進化

リーズ大学政治経済学教授のD.H.マクグレガー（M.A.） "感情を帯びずに読むことができる本は、不安な現在状況に興味を持つすべての人にとって有益です。" ア伯ンデーン・ジャーナル。

26. 農業

F.L.S.会員のW.ソムerville教授 "大学での実験室での研究結果を実践的な農民に容易に提供します。" エテュニアム。

30. 英国法の要素

オックスフォード大学英法律学ヴィナー座長のW.M.ゲルダート（M.A., B.C.L.） "英法律の基本原則を非常に明確に説明した内容を含んでおり、これらの基本原則について最小限の努力で理解したい人には強く推奨します。" スコットランド法時報。

38. 学校

教育研究への入門。

マンチェスター大学教育教授のJ.J.フィンダイル（M.A., Ph.D.） "驚くほど包括的な本です。…それは、簡潔で印象的な表現と共に主題の範囲も広大な驚異的なパフォーマンスです。" モーニングポスト。

-59. 政治経済学の要素

サミュエル・ジェイ・チャプマン、マ):?>
大学政治経済学教授による「最新の経済思想の光の中で需要と供給の動作原理、独占の性質、金銭と国際貿易、賃金、利益、利息、地租との関係、労働組合の影響を簡潔に説明した著作；またアダム・スミス以来の経済学研究の概要を付記している」。

準備中

犯罪者と共同体。ビスコント・スタリーズ、マ):?>
による著作。

法律における常識。ピーター・ヴィノグラドフ、ドクター・オブ・クラウン Laws による著作。

公務員制度。グレアム・ウォラス、マramento による著作。

実践的イデアル主義。モーリス・ヒューレットによる著作。

新聞。ジーニー・ビニー・ディブルーによる著作。

イングランドの村生活。エドワード・N・ベネット、マramento による著作。

共済と利益共有。ジェイ・P・ネウリンによる著作。

社会集落。ジェーン・アダムズとR.A.ウォーズによる著作。

偉大な発明。J.L.マイレス、マramento およびF.S.A.会員による著作。

都市計画。レイモンド・ウンwinによる著作。

イギリスの政治思想：ベネットからジェイ・エス・ミルまで。W.L.ダビッドソン、マramento による著作。

イギリスの政治思想：ハーバート・スピナーズから現在まで。エアーンスト・バッカー、マramento による著作。

ロンドン：ウィムツ\&ノーグテ

すべての書店とブックスタンドでご購入いただけます。

----------------------
ギュンター版ライセンス
以下はプロジェクト・ジーニアス電子書籍「数学への紹介」、アルフレッド・ノース・ホワイトヘッド著の著作権情報です。

\end{document}
